\paragraph{宏观完备而微观不可识别：固定纤维谱的折叠实现模空间与查询下界}
定理 \ref{thm:pom-oracle-tomography-completeness-f-Iu-Iw} 表明：在指数尺度上，容量曲线 \(\Gamma_{\mathrm{orc}}\) 足以反演计数谱 \(f\) 及两类速率函数。
然而 \(\Gamma_{\mathrm{orc}}\)（乃至完整的纤维谱 \(d_m\)）并不决定\emph{折叠实现本身}：同一纤维谱对应一个阶乘级的实现模空间，并诱导出截面选择的指数自由度与识别所需的查询下界。
本段将这些“宏观完备/微观不可识别”的事实写成可直接引用的闭式计数与信息论不等式。

\begin{definition}[固定纤维谱的微正则折叠实现类]\label{def:pom-microcanonical-fold-class}
固定分辨率 \(m\)，并记
\[
\Omega_m=\{0,1\}^m,
\qquad
X_m\ \text{为稳定类型集合},
\qquad
N:=\card{\Omega_m}=2^m,
\qquad
n:=\card{X_m}.
\]
给定纤维多重度函数 \(d:X_m\to\NN_{\ge 1}\) 满足 \(\sum_{x\in X_m}d(x)=N\)。
定义具有指定纤维谱 \(d\) 的折叠实现类为
\[
\mathfrak{Fold}_m(d)
:=\Bigl\{F:\Omega_m\to X_m:\ \card{F^{-1}(x)}=d(x)\ \ \forall x\in X_m\Bigr\}.
\]
并定义归一化纤维谱分布 \(w\in\Delta(X_m)\) 为 \(w(x):=d(x)/N\)。
\end{definition}

\begin{theorem}[同像计数闭式]\label{thm:pom-microcanonical-fold-class-count}
\leanverified{paper\_pom\_microcanonical\_fold\_class\_count}
在定义 \ref{def:pom-microcanonical-fold-class} 的设定下，有严格等式
\[
\boxed{\;
\card{\mathfrak{Fold}_m(d)}=\frac{N!}{\prod_{x\in X_m} d(x)!}\; } .
\]
\end{theorem}
\begin{proof}
构造 \(F\in\mathfrak{Fold}_m(d)\) 等价于将 \(\Omega_m\) 的 \(N\) 个元素划分为 \(n\) 个带标签的块 \(F^{-1}(x)\)（\(x\in X_m\)），且各块大小分别为 \(\{d(x)\}_{x\in X_m}\)。
带标签的多项式分割计数即为多项式系数
\(
\binom{N}{(d(x))_{x\in X_m}}=N!/\prod_x d(x)!
\)。
证毕。
\end{proof}

\leanverified{paper\_pom\_microcanonical\_fold\_entropy}
\begin{theorem}[微正则折叠熵：同谱模空间的指数率]\label{thm:pom-microcanonical-fold-entropy}
令 \(w(x)=d(x)/N\)，并记 Shannon 熵（自然对数）
\[
H(w):=-\sum_{x\in X_m} w(x)\log w(x).
\]
则有精确渐近式
\[
\boxed{\;
\log\card{\mathfrak{Fold}_m(d)}=N\,H(w)+O(n\log N)\;}
\]
并且在 \(n\log N=o(N)\) 的条件下，
\[
\boxed{\;
\frac{1}{N}\log\card{\mathfrak{Fold}_m(d)}=H(w)+o(1)\;}
\qquad(m\to\infty).
\]
上述条件在本文主线 \(n=\card{X_m}\asymp \varphi^m\)、\(N=2^m\) 下成立。
\end{theorem}
\begin{proof}
由定理 \ref{thm:pom-microcanonical-fold-class-count}，
\(
\log\card{\mathfrak{Fold}_m(d)}=\log N!-\sum_x\log d(x)!
\)。
对任意整数 \(k\ge 1\) 取 Stirling 展开
\(
\log k!=k\log k-k+\tfrac12\log(2\pi k)+O(1)
\)，得到
\[
\log\card{\mathfrak{Fold}_m(d)}
=N\log N-\sum_x d(x)\log d(x)+O\!\left(\log N+\sum_x\log d(x)\right).
\]
由 \(\sum_x d(x)=N\) 且 \(1\le d(x)\le N\) 得 \(\sum_x\log d(x)\le n\log N\)，从而误差项为 \(O(n\log N)\)。
主项化简为
\[
N\log N-\sum_x d(x)\log d(x)
=-\sum_x Nw(x)\log w(x)
=N H(w).
\]
于是得到第一式；若 \(n\log N=o(N)\) 则将误差除以 \(N\) 得第二式。证毕。
\end{proof}

\begin{corollary}[本体描述长度下界：仅知纤维谱时指定折叠的最小信息量]\label{cor:pom-microcanonical-fold-description-length-lower-bound}
在仅给定纤维谱 \(d\) 的前提下，任何能唯一指定某个 \(F\in\mathfrak{Fold}_m(d)\) 的编码所需比特数 \(L_m(d)\) 必满足
\[
\boxed{\;
L_m(d)\ \ge\ \log_2\card{\mathfrak{Fold}_m(d)}
=\frac{1}{\log 2}\log\card{\mathfrak{Fold}_m(d)}\; } .
\]
特别地，在定理 \ref{thm:pom-microcanonical-fold-entropy} 的条件下，
\[
L_m(d)\ \ge\ \frac{N}{\log 2}H(w)+o(N).
\]
\end{corollary}
\begin{proof}
若可用的码字总数小于 \(\card{\mathfrak{Fold}_m(d)}\)，则不可能对 \(\mathfrak{Fold}_m(d)\) 中所有对象给出单射编码；由鸽巢原理得到矛盾。证毕。
\end{proof}

\begin{corollary}[单点稀有性与低复杂度类的双指数稀薄性]\label{cor:pom-microcanonical-fold-singleton-and-low-complexity-rarity}
在定义 \ref{def:pom-microcanonical-fold-class} 的设定下，令 \(F\sim\mathrm{Unif}(\mathfrak{Fold}_m(d))\)。
则：
\begin{enumerate}
  \item 对任意固定 \(F_0\in\mathfrak{Fold}_m(d)\)，有
  \[
  \mathbb P(F=F_0)=\frac{1}{\card{\mathfrak{Fold}_m(d)}},
  \qquad
  \log\mathbb P(F=F_0)=-N H(w)+O(n\log N).
  \]
  特别地，只要 \(H(w)\) 在 \(m\to\infty\) 下不趋于 \(0\) 且 \(n\log N=o(N)\)，则该单点概率为 \(\exp(-\Theta(N))=\exp(-\Theta(2^m))\) 的双指数小量。
  \item 对任意函数类 \(\mathcal C_m\subseteq\mathfrak{Fold}_m(d)\)，若满足 \(\log\card{\mathcal C_m}\le L(m)\)，则
  \[
  \mathbb P(F\in\mathcal C_m)\le \exp\!\bigl(L(m)-N H(w)+O(n\log N)\bigr).
  \]
  因而当 \(L(m)=\mathrm{poly}(m)\) 且 \(n\log N=o(N)\) 时，有 \(\mathbb P(F\in\mathcal C_m)\le \exp(-\Theta(2^m))\)。
\end{enumerate}
\end{corollary}
\begin{proof}
由均匀先验，(1) 的第一式为定义，取对数并代入定理 \ref{thm:pom-microcanonical-fold-entropy} 得第二式；其后的双指数结论由 \(N=2^m\) 与 \(n\log N=o(N)\) 直接推出。
\par
(2) 由 \(\mathbb P(F\in\mathcal C_m)=\card{\mathcal C_m}/\card{\mathfrak{Fold}_m(d)}\) 与定理 \ref{thm:pom-microcanonical-fold-entropy} 代入立得。证毕。
\end{proof}

\paragraph{有限状态流式可实现性的指数稀有性：同谱实现类中的宽度下界}
在固定纤维谱的实现类 \(\mathfrak{Fold}_m(d)\) 中（定义 \ref{def:pom-microcanonical-fold-class}），
同像约束并不强迫“可由小状态机流式实现”：
即便只允许读取一次输入位，宽度受限的有限状态程序所能覆盖的函数族亦仅占极小比例。
本段给出纯计数上界，并由此导出微正则先验下的典型宽度下界。

\begin{definition}[只读一次的宽度-\texorpdfstring{$S$}{S} 流式分支程序]\label{def:pom-oblivious-widthS-branching-program}
固定长度 \(m\) 与状态数 \(S\ge 1\)。
称一组转移表与输出标记
\[
\delta_i:[S]\times\{0,1\}\to [S]\quad(1\le i\le m),
\qquad
\lambda:[S]\to X_m
\]
为\emph{只读一次（oblivious）}宽度-\(S\) 流式分支程序（或等价地：读入位序固定的确定性宽度-\(S\) 流式程序）。
取初始状态 \(s_0:=1\)，对输入 \(\omega=(\omega_1,\dots,\omega_m)\in\Omega_m=\{0,1\}^m\) 定义递推
\[
s_i:=\delta_i(s_{i-1},\omega_i)\quad(1\le i\le m),
\]
并令其计算的函数为
\[
F_{\delta,\lambda}(\omega):=\lambda(s_m)\in X_m.
\]
\end{definition}

\begin{theorem}[宽度受限流式程序的可实现函数计数上界]\label{thm:pom-oblivious-widthS-branching-program-count}
记 \(n:=\card{X_m}\)。
则宽度至多 \(S\) 的长度 \(m\) 只读一次流式分支程序所能实现的函数总数满足
\[
\boxed{\;
\mathcal N^{\mathrm{bp}}_{m,S}(n)
\ \le\ n^{S}\,S^{2Sm}\;}
\]
其中 \(\mathcal N^{\mathrm{bp}}_{m,S}(n)\) 表示所有可由某个宽度 \(\le S\) 程序实现的 \(F:\Omega_m\to X_m\) 的基数上界。
\leanverified{paper\_pom\_oblivious\_widthS\_branching\_program\_count}
\end{theorem}
\begin{proof}
对每个 \(i\)，转移表 \(\delta_i\) 在 \([S]\times\{0,1\}\) 的每个输入上取值于 \([S]\)，故其可能数为
\(
S^{2S}
\)；
于是 \((\delta_1,\dots,\delta_m)\) 的可能数至多为 \(S^{2Sm}\)。
输出标记 \(\lambda:[S]\to X_m\) 的可能数为 \(n^S\)。
不同程序产生的函数至多不同，故可实现函数总数受上述程序数目上界控制。证毕。
\end{proof}

\begin{theorem}[微正则实现类内的指数稀有性：小宽度程序几乎从不出现]\label{thm:pom-microcanonical-streaming-program-rarity}
\leanverified{paper\_pom\_microcanonical\_streaming\_program\_rarity}
在定义 \ref{def:pom-microcanonical-fold-class} 的设定下，令 \(F\sim\mathrm{Unif}(\mathfrak{Fold}_m(d))\)。
则对任意 \(S\ge 1\)，有
\[
\mathbb{P}\Bigl(F\ \text{可由某个宽度}\le S\ \text{的长度 }m\ \text{流式分支程序实现}\Bigr)
\ \le\
\exp\!\Bigl(2Sm\log S+S\log n-\log\card{\mathfrak{Fold}_m(d)}\Bigr).
\]
并且结合定理 \ref{thm:pom-microcanonical-fold-entropy}（令 \(w(x)=d(x)/N\)）得到
\[
\boxed{\;
\mathbb{P}(\cdot)
\ \le\
\exp\!\Bigl(2Sm\log S+S\log n-NH(w)+O(n\log N)\Bigr)\;}
\qquad(N=\card{\Omega_m}=2^m).
\]
\end{theorem}
\begin{proof}
令 \(\mathcal{C}_{m,S}\) 为所有可由某个宽度 \(\le S\) 流式分支程序实现的函数集合。
由定理 \ref{thm:pom-oblivious-widthS-branching-program-count} 得
\(
\card{\mathcal{C}_{m,S}}\le n^S S^{2Sm}
\)。
由 \(F\) 在 \(\mathfrak{Fold}_m(d)\) 上均匀，得到
\[
\mathbb{P}(F\in\mathcal{C}_{m,S})
=\frac{\card{\mathcal{C}_{m,S}\cap\mathfrak{Fold}_m(d)}}{\card{\mathfrak{Fold}_m(d)}}
\le \frac{\card{\mathcal{C}_{m,S}}}{\card{\mathfrak{Fold}_m(d)}}.
\]
取对数并指数化即得第一式；第二式由定理 \ref{thm:pom-microcanonical-fold-entropy} 代入得到。证毕。
\end{proof}

\begin{corollary}[典型最小宽度的指数下界]\label{cor:pom-microcanonical-streaming-program-typical-width-lb}
设存在常数 \(c_0>0\) 使 \(H(w)\ge c_0\log n\)。
并假设 \(\log n=O(m)\) 且 \(n\log N=o(N)\)。
则存在常数 \(c_1,c_2>0\)，使得对所有满足
\[
S\ \le\ c_1\,\frac{N}{m\log N},
\]
都有
\[
\boxed{\;
\mathbb{P}\Bigl(\min\{S:\ F\ \text{可由宽度}\le S\ \text{流式分支程序实现}\}\le S\Bigr)
\ \le\ \exp(-c_2N)\;}
\qquad
F\sim\mathrm{Unif}(\mathfrak{Fold}_m(d)).
\]
特别地，在本文主线 \(N=2^m\) 下，典型最小宽度满足
\[
\min\{S:\ F\ \text{可由宽度}\le S\ \text{流式分支程序实现}\}
\ \ge\ \Omega\!\left(\frac{2^m}{m^2}\right)
\]
以概率 \(1-\exp(-\Theta(2^m))\) 成立。
\end{corollary}
\begin{proof}
由定理 \ref{thm:pom-microcanonical-streaming-program-rarity}，主指数为
\(
2Sm\log S+S\log n-NH(w)+O(n\log N)
\)。
在 \(\log n=O(m)\) 且 \(\log S\le \log N=O(m)\) 下，有
\[
2Sm\log S+S\log n\ \le\ C\,S\,m\log N,
\]
其中常数 \(C\) 仅依赖于 \(\log n=O(m)\) 的隐含常数。
另一方面，由 \(H(w)\ge c_0\log n\) 且 \(n\ge 2\)（从而 \(\log n\ge \log 2\)）得
\(
NH(w)\ge c_0N\log n\ge c_0N\log 2
\)。
取 \(c_1\) 足够小，使当 \(S\le c_1 N/(m\log N)\) 时有 \(2Sm\log S+S\log n\le \tfrac{1}{2}c_0N\log 2\)；
再用 \(n\log N=o(N)\) 吸收误差项，即得指数严格小于 \(-c_2N\) 的概率界。
最后一句由 \(\log N=\Theta(m)\) 代入化简。证毕。
\end{proof}

\endinput


\begin{theorem}[截面选择的指数自由度：右逆空间的精确计数]\label{thm:pom-microcanonical-section-count}
\leanverified{paper\_pom\_microcanonical\_section\_count}
对任意 \(F\in\mathfrak{Fold}_m(d)\)，定义其右逆（截面）集合
\[
\mathfrak{S}(F):=\Bigl\{s:X_m\to \Omega_m:\ F\circ s=\mathrm{id}_{X_m}\Bigr\}.
\]
则有闭式
\[
\boxed{\;
\card{\mathfrak{S}(F)}=\prod_{x\in X_m} d(x)\; } ,
\qquad
\boxed{\;
\log\card{\mathfrak{S}(F)}=\sum_{x\in X_m}\log d(x)\; } .
\]
若记 \(\nu:=\mathrm{Unif}(X_m)\)，并定义
\(
\kappa_\nu(d):=\mathbb{E}_{\nu}[\log d(X)]=\frac{1}{n}\sum_x\log d(x)
\)，则
\[
\boxed{\;
\log\card{\mathfrak{S}(F)}=n\,\kappa_\nu(d).\;}
\]
\end{theorem}
\begin{proof}
构造 \(s\in\mathfrak{S}(F)\) 等价于对每个 \(x\in X_m\) 从纤维 \(F^{-1}(x)\) 中选择一个代表元；各纤维选择彼此独立，计数为乘积 \(\prod_x d(x)\)。
取对数并改写为 \(\nu\) 下的均值即得。证毕。
\end{proof}

\begin{theorem}[谱不可导出截面：仅凭纤维谱无法产生普适有效的规范右逆]\label{thm:pom-spectrum-cannot-induce-universal-section}
\leanverified{paper\_pom\_spectrum\_cannot\_induce\_universal\_section}
设存在一“谱\(\to\)截面”规则 \(\mathcal{A}\)，其输入仅为纤维谱 \(d\)，输出为映射 \(s_d:X_m\to\Omega_m\)。
称 \(\mathcal{A}\) 为\emph{普适有效}，若对所有 \(F\in\mathfrak{Fold}_m(d)\) 都有 \(F\circ s_d=\mathrm{id}_{X_m}\)。
若满足 \(n=\card{X_m}\ge 2\) 且存在某个 \(x_0\in X_m\) 使 \(d(x_0)\ge 2\)，则不存在普适有效的 \(\mathcal{A}\)。
\end{theorem}
\begin{proof}
反设存在普适有效的 \(s_d\)。
任取 \(F\in\mathfrak{Fold}_m(d)\) 并取 \(x_0\in X_m\) 使 \(d(x_0)\ge 2\)。
令 \(\omega_0:=s_d(x_0)\in\Omega_m\)，则由普适有效性知 \(F(\omega_0)=x_0\)。
\par
由于 \(n\ge 2\) 且 \(d(x)\ge 1\) 对所有 \(x\) 成立，存在某个 \(x_1\neq x_0\) 以及 \(\omega_1\in\Omega_m\) 使 \(F(\omega_1)=x_1\neq x_0\)。
取置换 \(\sigma\) 交换 \(\omega_0\) 与 \(\omega_1\)，并令 \(F':=F\circ\sigma^{-1}\)。
则 \(\card{(F')^{-1}(x)}=\card{F^{-1}(x)}=d(x)\) 对一切 \(x\) 成立，即 \(F'\in\mathfrak{Fold}_m(d)\)。
但
\[
F'(s_d(x_0))=F'(\omega_0)=F(\sigma^{-1}(\omega_0))=F(\omega_1)\neq x_0,
\]
从而 \(F'\circ s_d\neq \mathrm{id}_{X_m}\)，与普适有效性矛盾。证毕。
\end{proof}

\begin{theorem}[学习/审计的查询下界：识别同谱折叠所需查询次数]\label{thm:pom-microcanonical-fold-query-lower-bound}
\leanverified{paper\_pom\_microcanonical\_fold\_query\_lower\_bound}
在黑箱查询模型中，未知 \(F\in\mathfrak{Fold}_m(d)\) 作为 oracle 给出；算法可自适应选择
\(\omega^{(1)},\dots,\omega^{(k)}\in\Omega_m\)，并观察响应 \(F(\omega^{(i)})\in X_m\)。
若算法在确定性意义下对所有 \(F\in\mathfrak{Fold}_m(d)\) 都能唯一确定 \(F\)，则其查询次数必满足
\[
\boxed{\;
k\ \ge\ \left\lceil \frac{\log\card{\mathfrak{Fold}_m(d)}}{\log\card{X_m}}\right\rceil\; } .
\]
结合定理 \ref{thm:pom-microcanonical-fold-entropy}，在 \(n\log N=o(N)\) 下有
\[
k\ \ge\ \frac{N\,H(w)}{\log n}+o\!\left(\frac{N}{\log n}\right).
\]
在本文主线 \(n\asymp\varphi^m\) 下，\(\log n=\Theta(m)\)，从而 \(k=\Omega(2^m/m)\)。
\end{theorem}
\begin{proof}
任意 \(k\) 次查询产生的响应轨迹至多有 \(n^k=\card{X_m}^k\) 种可能。
若 \(n^k<\card{\mathfrak{Fold}_m(d)}\)，则存在 \(F\neq F'\) 产生同一轨迹（鸽巢原理），算法无法区分，因而不可能精确识别。
故必须有 \(n^k\ge \card{\mathfrak{Fold}_m(d)}\)，取对数即得第一式。
其余结论由定理 \ref{thm:pom-microcanonical-fold-entropy} 代入并用 \(\log n=\Theta(m)\) 得到。证毕。
\end{proof}

\begin{remark}[计数下界与查询刚性]
定理 \ref{thm:pom-microcanonical-fold-query-lower-bound} 仅使用“响应字母表大小为 \(n\)”这一粗信息，
并未利用同谱约束 \(\card{F^{-1}(x)}=d(x)\) 所诱导的无放回结构。
在微正则语义下，精确识别还满足更强的刚性：只要仍存在至少两个未查询点，即可在保持计数约束的前提下交换其标签而不改变全部已观测响应。
\end{remark}

\begin{theorem}[最坏情形查询刚性：非平凡纤维谱下精确识别需 \(N-1\) 次查询]\label{thm:pom-microcanonical-fold-worstcase-query-rigidity}
\leanverified{paper\_pom\_microcanonical\_fold\_worstcase\_query\_rigidity}
设 \(n=\card{X_m}\ge 2\)。
考虑确定性自适应查询算法：每次选择尚未查询的 \(\omega\in\Omega_m\) 并观察 \(F(\omega)\)，在 \(k\) 次查询后输出估计 \(\widehat F\)。
若算法满足对一切 \(F\in\mathfrak{Fold}_m(d)\) 都有 \(\widehat F=F\)，则必有
\[
\boxed{\;k\ge N-1\;}
\qquad(N=\card{\Omega_m}=2^m).
\]
\end{theorem}
\begin{proof}
取任意确定性算法 \(\mathcal{A}\)，反设其在最坏情形下仅查询 \(k\le N-2\) 次。
由 \(n\ge 2\) 且 \(d(x)\ge 1\)（对所有 \(x\)）可取不同标签 \(x_\star\neq y_\star\in X_m\) 使 \(d(x_\star),d(y_\star)\ge 1\)。
\par
以对抗方式生成回答：在前 \(k\) 次查询中，从每个标签 \(x\) 具有 \(d(x)\) 个“可用计数”的多重集中依次取出回答，并预先保留 \(x_\star\) 与 \(y_\star\) 各一单位不被使用。
由于除去这两单位后仍有 \(N-2\) 个计数可用，而 \(k\le N-2\)，该策略可行且不会违反任意标签的计数上限。
由此得到查询点集合 \(Q\subset\Omega_m\)（\(\card{Q}=k\)）及其在 \(Q\) 上的响应轨迹。
\par
令 \(U:=\Omega_m\setminus Q\)，则 \(\card{U}=N-k\ge 2\)。
在 \(U\) 中任选 \(\omega_1\neq\omega_2\)。
在保持 \(Q\) 上响应轨迹不变的前提下，构造两映射 \(F_0,F_1\in\mathfrak{Fold}_m(d)\)：
令
\[
F_0(\omega_1)=x_\star,\quad F_0(\omega_2)=y_\star,\qquad
F_1(\omega_1)=y_\star,\quad F_1(\omega_2)=x_\star,
\]
并将 \(U\setminus\{\omega_1,\omega_2\}\) 处的取值按剩余计数（总和为 \(\card{U}-2\)）任意补全。
该补全可行是因为上述回答策略保证了 \(x_\star,y_\star\) 的各一单位仍未使用，且其余标签的未使用计数总和恰为 \(\card{U}-2\)。
\par
由于 \(\mathcal{A}\) 未查询 \(\omega_1,\omega_2\)，故在 \(Q\) 上 \(F_0\) 与 \(F_1\) 给出完全相同的观测轨迹，从而 \(\mathcal{A}\) 的输出必相同；
但 \(F_0\neq F_1\)，与“对所有 \(F\)”精确识别矛盾。
\end{proof}

\paragraph{微正则先验下的后验模空间：计数闭式、轨迹概率与 Bayes 识别率}
以下把“同谱折叠实现模空间”赋予均匀先验，并将任意有限查询轨迹的后验规模与 Bayes 最优精确识别率写成完全显式的组合闭式。

\begin{theorem}[后验模空间的精确计数与轨迹概率]\label{thm:pom-microcanonical-fold-posterior-count-and-prob}
\leanverified{paper\_pom\_microcanonical\_fold\_posterior\_count\_and\_prob}
令 \(F\sim\mathrm{Unif}(\mathfrak{Fold}_m(d))\)。
对任意互异查询点集 \(Q=\{\omega_1,\dots,\omega_k\}\subset\Omega_m\) 与任意响应序列 \(\mathbf{a}=(a_1,\dots,a_k)\in X_m^k\)，记其计数向量
\[
c_{\mathbf{a}}(x):=\card{\{i\in\{1,\dots,k\}:\ a_i=x\}},\qquad x\in X_m,\qquad \sum_{x}c_{\mathbf{a}}(x)=k.
\]
则满足 \(F(\omega_i)=a_i\ (1\le i\le k)\) 的实现数为
\[
\boxed{\;
\card{\Bigl\{F\in\mathfrak{Fold}_m(d):\ F(\omega_i)=a_i,\ 1\le i\le k\Bigr\}}
=\frac{(N-k)!}{\prod_{x\in X_m}\bigl(d(x)-c_{\mathbf{a}}(x)\bigr)!}\;},
\]
其中若存在 \(x\) 使 \(c_{\mathbf{a}}(x)>d(x)\) 则右端按 \(0\) 解释。
相应地，该轨迹在均匀先验下出现的概率为
\[
\boxed{\;
\mathbb{P}(\mathbf{a})
=\frac{(N-k)!}{N!}\prod_{x\in X_m}\frac{d(x)!}{\bigl(d(x)-c_{\mathbf{a}}(x)\bigr)!}
=\frac{1}{(N)_k}\prod_{x\in X_m}\bigl(d(x)\bigr)_{c_{\mathbf{a}}(x)}\;},
\]
其中 \((u)_r:=u(u-1)\cdots(u-r+1)\) 为下降阶乘，且 \((u)_0:=1\)，并约定 \((N)_k:=N(N-1)\cdots(N-k+1)\)。
因此 \(\mathbb{P}(\mathbf{a})\) 与兼容实现数仅依赖于计数向量 \(c_{\mathbf{a}}\) 与查询次数 \(k\)，与 \(Q\) 的具体位置无关。
\end{theorem}
\begin{proof}
固定 \(Q\) 上的响应后，剩余 \(N-k\) 个微态必须按剩余计数
\(
d(x)-c_{\mathbf{a}}(x)
\)
分配到各标签（\(x\in X_m\)），该分配数为带标签的多项式系数
\(
(N-k)!/\prod_x(d(x)-c_{\mathbf{a}}(x))!
\)；
若某 \(x\) 使 \(c_{\mathbf{a}}(x)>d(x)\)，则不存在可行分配，计数为 \(0\)。
\par
轨迹概率由定理 \ref{thm:pom-microcanonical-fold-class-count} 的
\(
\card{\mathfrak{Fold}_m(d)}=N!/\prod_x d(x)!
\)
作比值得到，并用恒等式
\(
d(x)!/(d(x)-r)!= (d(x))_{r}
\)
与
\(
N!/(N-k)!=(N)_k
\)
化简即得。
最后一条“仅依赖计数”的结论由公式中不出现 \(Q\) 的具体结构直接推出。证毕。
\end{proof}

\paragraph{自适应无增益与充分统计：微正则先验下的无放回抽样与计数后验}
以下在均匀微正则先验 \(F\sim\mathrm{Unif}(\mathfrak{Fold}_m(d))\) 下，
把任意确定性自适应查询过程严格等价为对多重集 \(\{x^{\times d(x)}\}_{x\in X_m}\) 的无放回抽样；
并指出计数向量 \(C_k\) 在后验意义下构成关于 \(F\) 的最小充分统计量。

\begin{definition}[自适应互异查询轨迹与计数统计]\label{def:pom-microcanonical-adaptive-trajectory-counts}
令 \(\mathcal{A}\) 为任意确定性自适应算法，进行 \(k\) 次互异查询：第 \(i\) 步选择
\[
\omega_i\in\Omega_m\setminus\{\omega_1,\dots,\omega_{i-1}\},
\]
并观察响应
\[
T_i:=F(\omega_i)\in X_m.
\]
记轨迹 \(T_{1:k}:=(T_1,\dots,T_k)\in X_m^k\)，并定义计数向量
\[
C_k(x):=\card{\{i\in\{1,\dots,k\}:\ T_i=x\}},\qquad x\in X_m.
\]
\end{definition}

\begin{theorem}[自适应无增益：轨迹分布与逐步无放回抽样律]\label{thm:pom-microcanonical-adaptive-no-gain-without-replacement}
在定义 \ref{def:pom-microcanonical-adaptive-trajectory-counts} 的设定下：
\begin{enumerate}
  \item \textbf{（轨迹分布与策略无关）}
  对任意 \(t=(t_1,\dots,t_k)\in X_m^k\)，其出现概率仅依赖于 \(k\) 与纤维谱 \(d\)，并满足
  \[
  \boxed{\;
  \mathbb{P}_{\mathcal{A}}(T_{1:k}=t)
  =\frac{1}{(N)_k}\prod_{x\in X_m}\bigl(d(x)\bigr)_{c_t(x)}\;},
  \]
  其中 \(c_t(x)=\card{\{i:\ t_i=x\}}\) 为 \(t\) 的计数向量，\((u)_r:=u(u-1)\cdots(u-r+1)\) 为下降阶乘，
  并约定 \((N)_k:=N(N-1)\cdots(N-k+1)\)。
  \item \textbf{（逐步无放回抽样）}
  对每个 \(i\ge 1\) 与任意 \(x\in X_m\)，有条件概率公式
  \[
  \boxed{\;
  \mathbb{P}_{\mathcal{A}}\!\bigl(T_i=x\mid T_{1:i-1}\bigr)
  =\frac{d(x)-C_{i-1}(x)}{N-i+1}\; } .
  \]
  因而 \((T_i)\) 与从多重集 \(\{x^{\times d(x)}\}_{x\in X_m}\) 的无放回抽样过程同分布。
\end{enumerate}
\end{theorem}
\begin{proof}
对任意给定轨迹 \(t\in X_m^k\)，算法 \(\mathcal{A}\) 的查询点序列 \((\omega_i)\) 是 \(t\) 的前缀的确定性函数；
因此事件 \(\{T_{1:k}=t\}\) 等价于某个互异点集 \(Q_t=\{\omega_1(t),\dots,\omega_k(t)\}\subset\Omega_m\) 上满足
\(
F(\omega_i(t))=t_i
\)。
由定理 \ref{thm:pom-microcanonical-fold-posterior-count-and-prob} 的轨迹概率公式，
\(\mathbb{P}(T_{1:k}=t)\) 仅依赖于 \(c_t\) 并给出所示表达；该表达中不含 \(Q_t\) 的几何信息，故亦与 \(\mathcal{A}\) 的选择规则无关。
\par
逐步条件概率由比值计算得到：对任意历史前缀 \(t_{1:i-1}\) 与任意 \(x\in X_m\)，
将定理 \ref{thm:pom-microcanonical-fold-posterior-count-and-prob} 的概率公式分别应用于长度 \(i\) 的轨迹 \((t_{1:i-1},x)\) 与长度 \(i-1\) 的轨迹 \(t_{1:i-1}\)，并取比值得
\[
\mathbb{P}(T_i=x\mid T_{1:i-1}=t_{1:i-1})
=\frac{(d(x))_{c_{t_{1:i-1}}(x)+1}}{(d(x))_{c_{t_{1:i-1}}(x)}}\cdot\frac{(N)_{i-1}}{(N)_{i}}
=\frac{d(x)-C_{i-1}(x)}{N-i+1},
\]
其中使用了恒等式 \((u)_{r+1}=(u-r)(u)_r\)。
结论成立。证毕。
\end{proof}
\leanverified{paper\_pom\_microcanonical\_adaptive\_no\_gain\_without\_replacement}

\begin{theorem}[计数充分统计与后验均匀性]\label{thm:pom-microcanonical-count-sufficient-statistic-posterior-uniform}
在定义 \ref{def:pom-microcanonical-adaptive-trajectory-counts} 的设定下，设算法在 \(k\) 次查询后得到轨迹 \(t\in X_m^k\)，并令 \(Q_k:=\{\omega_1,\dots,\omega_k\}\)、
\(U_k:=\Omega_m\setminus Q_k\)。
\begin{enumerate}
  \item \textbf{（后验仅依赖计数）}
  对任意两条轨迹 \(t,t'\in X_m^k\)，若 \(c_t=c_{t'}\)，则诱导的后验分布同构：
  存在域置换 \(\sigma\in\mathrm{Sym}(\Omega_m)\) 使
  \[
  \mathrm{Law}\bigl(F\mid T_{1:k}=t\bigr)\ \cong\ \mathrm{Law}\bigl(F\circ\sigma^{-1}\mid T_{1:k}=t'\bigr).
  \]
  因而轨迹对后验的全部影响可压缩为计数向量 \(C_k\)，并且 \(C_k\) 在该意义下为关于 \(F\) 的最小充分统计量。
  \item \textbf{（未查询部分的条件后验均匀性）}
  在给定轨迹 \(t\) 时，限制在未查询集 \(U_k\) 上的后验为均匀分布：
  \[
  \mathrm{Law}\bigl(F|_{U_k}\mid T_{1:k}=t\bigr)
  =\mathrm{Unif}\bigl(\mathfrak{F}_{U_k}(d-c_t)\bigr),
  \]
  其中 \(\mathfrak{F}_{U_k}(d-c_t)\) 表示把 \(U_k\) 映射到 \(X_m\) 且纤维大小为 \(d(x)-c_t(x)\) 的全部映射集合。
\end{enumerate}
\end{theorem}
\leanverified{paper\_pom\_microcanonical\_count\_sufficient\_statistic\_posterior\_uniform}
\begin{proof}
对任意给定的查询点集 \(Q_k\) 与其上标签赋值 \(t\)，兼容集合
\[
\Bigl\{F\in\mathfrak{Fold}_m(d):\ F(\omega_i)=t_i,\ 1\le i\le k\Bigr\}
\]
仅要求在 \(U_k\) 上实现剩余计数 \(d-c_t\)；
由于先验在 \(\mathfrak{Fold}_m(d)\) 上均匀，条件化后在该兼容集合上仍均匀。
将上述均匀性限制到 \(U_k\) 即得第 (2) 点。
\par
若两条轨迹 \(t,t'\) 具有相同计数向量，则存在一个置换把 \(Q_k\) 上的标签分配从 \(t\) 搬运到 \(t'\) 而不改变每个标签的使用次数；
由先验在 \(\Omega_m\) 上的置换不变性可将两后验识别为同构，从而得到第 (1) 点。
最小充分性由“后验仅依赖 \(c_t\)”与“剩余计数 \(d-c_t\) 唯一确定后验可行集合”直接推出。证毕。
\end{proof}

\endinput


\paragraph{截面可学习性的覆盖时刻：策略不变性、精确尾公式与 \texorpdfstring{$n\log n$}{n log n} cutoff}
在微正则先验下，任意自适应互异查询产生的类型轨迹与无放回抽样同分布（定理 \ref{thm:pom-microcanonical-adaptive-no-gain-without-replacement}）。
因此，“首次收齐全部稳定类型”的覆盖时刻与策略无关，并可写成包含--排除闭式；
在均衡纤维条件下，该覆盖时刻呈现 \(n\log n\) 标度与 Gumbel cutoff。
这为“获得一条右逆截面”与“精确识别整张映射”之间的复杂度分离提供了可审计的精确模型。

\begin{definition}[覆盖时刻与缺失类型数]\label{def:pom-microcanonical-section-cover-time}
令 \(F\sim\mathrm{Unif}(\mathfrak{Fold}_m(d))\)，并令 \(\mathcal{A}\) 为任意确定性自适应算法进行互异查询（定义 \ref{def:pom-microcanonical-adaptive-trajectory-counts}），得到轨迹 \(T_1,T_2,\dots\in X_m\)。
定义覆盖时刻
\[
\tau_{\mathcal A}:=\min\Bigl\{t\in\{0,1,\dots,N\}:\ \{T_1,\dots,T_t\}=X_m\Bigr\},
\]
并对任意 \(t\in\{0,1,\dots,N\}\) 定义缺失类型数
\[
M(t):=\card{X_m\setminus\{T_1,\dots,T_t\}}.
\]
\end{definition}

\begin{theorem}[覆盖时刻的策略不变性]\label{thm:pom-microcanonical-cover-time-strategy-invariance}
在定义 \ref{def:pom-microcanonical-section-cover-time} 的设定下，对任意两种自适应互异查询策略 \(\mathcal{A},\mathcal{A}'\)，有
\[
\boxed{\;
\tau_{\mathcal A}\ \stackrel{d}{=}\ \tau_{\mathcal A'}\;}
\qquad\text{并且}\qquad
\boxed{\;
M_{\mathcal A}(t)\ \stackrel{d}{=}\ M_{\mathcal A'}(t)\ \ \forall\,t\;}
\]
其中 \(M_{\mathcal A}(t)\) 表示策略 \(\mathcal A\) 下的缺失类型数过程。
\end{theorem}
\begin{proof}
由定理 \ref{thm:pom-microcanonical-adaptive-no-gain-without-replacement}，
\((T_i)_{i\ge 1}\) 的联合分布与策略无关，且等价于从多重集 \(\{x^{\times d(x)}\}_{x\in X_m}\) 的无放回抽样。
覆盖时刻 \(\tau_{\mathcal A}\) 与缺失过程 \(M_{\mathcal A}(t)\) 均是轨迹 \((T_i)\) 的对称泛函，
从而其分布亦与策略无关。证毕。
\end{proof}

\leanverified{paper\_pom\_microcanonical\_cover\_time\_strategy\_invariance}

\begin{theorem}[覆盖时刻的精确尾概率：包含--排除闭式]\label{thm:pom-microcanonical-cover-time-tail-inclusion-exclusion}
沿用定义 \ref{def:pom-microcanonical-fold-class} 的记号，并对任意子集 \(A\subseteq X_m\) 记
\[
d(A):=\sum_{x\in A}d(x).
\]
令 \(\tau\) 为无放回抽样覆盖时刻（与 \(\tau_{\mathcal A}\) 同分布，见定理 \ref{thm:pom-microcanonical-cover-time-strategy-invariance}）。
则对任意 \(t\in\{0,1,\dots,N\}\)，成立严格恒等式
\[
\boxed{\;
\mathbb{P}(\tau>t)
=\sum_{\emptyset\neq A\subseteq X_m}(-1)^{\card{A}+1}\,
\frac{\binom{N-d(A)}{t}}{\binom{N}{t}}\;}
\]
其中当 \(N-d(A)<t\) 时相应项按 \(0\) 解释。
\end{theorem}
\begin{proof}
事件 \(\{\tau>t\}\) 等价于“存在某些类型在前 \(t\) 次抽样中从未出现”。
对固定非空 \(A\subseteq X_m\)，事件“\(A\) 中全部类型在前 \(t\) 步均缺失”当且仅当
从 \(N-d(A)\) 个非 \(A\) 元素中无放回取出 \(t\) 个，故其概率为 \(\binom{N-d(A)}{t}/\binom{N}{t}\)。
对“至少一个类型缺失”应用包含--排除即可得到所示公式。证毕。
\end{proof}
\leanverified{paper\_pom\_microcanonical\_cover\_time\_tail\_inclusion\_exclusion}

\begin{corollary}[覆盖时刻的期望闭式]\label{cor:pom-microcanonical-cover-time-mean}
在定理 \ref{thm:pom-microcanonical-cover-time-tail-inclusion-exclusion} 的设定下，
\[
\boxed{\;
\EE[\tau]
=\sum_{t=0}^{N-1}\mathbb{P}(\tau>t)
=\sum_{t=0}^{N-1}\ \sum_{\emptyset\neq A\subseteq X_m}(-1)^{\card{A}+1}\,
\frac{\binom{N-d(A)}{t}}{\binom{N}{t}}\;}
\]
并且对任意策略 \(\mathcal A\) 有 \(\EE[\tau_{\mathcal A}]=\EE[\tau]\)。
\end{corollary}
\begin{proof}
由非负整数值随机变量的恒等式 \(\EE[\tau]=\sum_{t\ge 0}\mathbb{P}(\tau>t)\) 且 \(\tau\le N\)，得第一式。
第二句由定理 \ref{thm:pom-microcanonical-cover-time-strategy-invariance} 立得。证毕。
\end{proof}
\leanverified{paper\_pom\_microcanonical\_cover\_time\_mean}

\begin{theorem}[均衡纤维下的 \texorpdfstring{$n\log n$}{n log n} 标度]\label{thm:pom-microcanonical-cover-time-nlogn-scale}
\leanverified{paper\_pom\_microcanonical\_cover\_time\_nlogn\_scale}
令 \(n=\card{X_m}\to\infty\)，并假设均衡条件：存在常数 \(0<c_-\le c_+<\infty\) 使对所有 \(x\in X_m\) 有
\[
c_-\frac{N}{n}\ \le\ d(x)\ \le\ c_+\frac{N}{n}.
\]
再假设 \(n\log n=o(N)\)。
则存在常数 \(C_1,C_2>0\)（仅依赖 \(c_\pm\)）使
\[
\boxed{\;
n\log n-C_1n\ \le\ \EE[\tau]\ \le\ n\log n+C_2n\;}
\]
并且有概率型集中
\[
\boxed{\;
\tau=n\log n+O_{\mathbb{P}}(n)\;}
\qquad(n\to\infty).
\]
\end{theorem}
\begin{proof}
对任意单点集合 \(\{x\}\) 由无放回抽样得
\[
\mathbb{P}(x\ \text{在前 }t\ \text{步缺失})
=\frac{\binom{N-d(x)}{t}}{\binom{N}{t}}
=\prod_{j=0}^{t-1}\left(1-\frac{d(x)}{N-j}\right).
\]
当 \(t\le n\log n+O(n)\) 且 \(n\log n=o(N)\) 时，右端与
\(\exp(-t\,d(x)/N)\) 在指数阶上一致。
由均衡性 \(d(x)/N\asymp 1/n\)，对 \(t=n(\log n+s)\) 有
\[
\mathbb{P}(x\ \text{缺失})\asymp \frac{e^{-s}}{n}.
\]
\par
\textbf{上界：}并集界给出 \(\mathbb{P}(\tau>t)\le\sum_{x}\mathbb{P}(x\ \text{缺失})\le C e^{-s}\)，
对 \(t\) 求和（或积分比较）得到 \(\EE[\tau]\le n\log n+O(n)\)。
\par
\textbf{下界：}令 \(M(t)\) 为时刻 \(t\) 的缺失类型数。
在无放回抽样下，缺失指示族满足负相关性，从而 \(\Var(M(t))\le \EE[M(t)]\)。
取 \(t=n(\log n-s)\) 使 \(\EE[M(t)]\ge c e^{s}\)。
由 Paley--Zygmund 不等式得到 \(\mathbb{P}(M(t)>0)\) 有界离零，
从而 \(\mathbb{P}(\tau>t)\) 在该尺度上不趋于 \(0\)；
再对 \(t\) 求和得到 \(\EE[\tau]\ge n\log n-O(n)\)。
概率型集中由对 \(\mathbb{P}(\tau>n(\log n+s))\) 的指数尾界与相同尺度的下界合并推出。证毕。
\end{proof}

\begin{corollary}[Fibonacci 尺度下的指数分离：截面获得 \texorpdfstring{$\ll 2^m$}{<<2^m} 而全识别需 \texorpdfstring{$2^m$}{2^m} 级别]\label{cor:pom-microcanonical-cover-time-fibonacci-separation}
在本文折叠主线 \(N=\card{\Omega_m}=2^m\)、\(n=\card{X_m}=F_{m+2}\asymp\varphi^m\)（见 \S\ref{subsec:pom-fold-compression}）下，
若均衡条件在该尺度族上成立，则由定理 \ref{thm:pom-microcanonical-cover-time-nlogn-scale} 得
\[
\EE[\tau]=n\log n+O(n)=\Theta(m\varphi^m),
\qquad
\frac{\EE[\tau]}{2^m}=\Theta\!\bigl(m(\varphi/2)^m\bigr)\xrightarrow[m\to\infty]{}0.
\]
因此，在微正则先验下“获得一条右逆截面（每类取一微态代表）”的典型查询规模为 \(\Theta(m\varphi^m)\)，
而“精确识别整张映射”在最坏情形至少需要 \(N-1=2^m-1\) 次查询（定理 \ref{thm:pom-microcanonical-fold-worstcase-query-rigidity}），两者发生指数分离。
\leanverified{compression\_growth}
\leanverified{compression\_ratios}
\leanverified{paper\_pom\_microcanonical\_cover\_time\_fibonacci\_separation}
\end{corollary}

\begin{theorem}[无放回均衡收集的 Gumbel 极限：覆盖时刻出现 cutoff]\label{thm:pom-microcanonical-cover-time-gumbel-cutoff}
设 \(N\to\infty\)、\(n\to\infty\)，并给定计数 \(d_1,\dots,d_n\in\NN\) 满足 \(\sum_{i=1}^n d_i=N\)。
假设均衡性
\[
\max_{1\le i\le n}\left|\frac{d_i}{N}-\frac{1}{n}\right|=o\!\left(\frac{1}{n}\right),
\qquad
n(\log n)^2=o(N).
\]
令 \(\tau\) 为无放回抽样下首次出现全部 \(n\) 种类型的时间（定理 \ref{thm:pom-microcanonical-cover-time-strategy-invariance} 表明：微正则先验下任意自适应互异查询的覆盖时刻与该 \(\tau\) 同分布）。
对任意固定 \(c\in\RR\)，令
\[
t_n(c):=\left\lfloor n(\log n+c)\right\rfloor.
\]
则成立极限
\[
\boxed{\;
\lim_{N\to\infty}\mathbb{P}(\tau\le t_n(c))=\exp(-e^{-c})\;},
\qquad
\frac{\tau-n\log n}{n}\xrightarrow{d}G,
\]
其中 \(G\) 为标准 Gumbel 分布（\(\mathbb{P}(G\le c)=\exp(-e^{-c})\)）。
\leanverified{paper\_pom\_microcanonical\_cover\_time\_gumbel\_cutoff}
\end{theorem}
\begin{proof}
对任意非空子集 \(A\subseteq[n]\)，记 \(d(A):=\sum_{i\in A}d_i\)。
由无放回抽样，事件“\(A\) 中全部类型在前 \(t\) 步缺失”的概率为
\[
\frac{\binom{N-d(A)}{t}}{\binom{N}{t}}
=\prod_{j=0}^{t-1}\left(1-\frac{d(A)}{N-j}\right).
\]
在 \(t=t_n(c)=O(n\log n)\) 与 \(n(\log n)^2=o(N)\) 下，有一致展开
\[
\log\prod_{j=0}^{t-1}\left(1-\frac{d(A)}{N-j}\right)
=t\log\!\left(1-\frac{d(A)}{N}\right)+o(1).
\]
由均衡性得 \(d(A)=\frac{\card{A}}{n}N(1+o(1))\)，从而
\[
\frac{\binom{N-d(A)}{t}}{\binom{N}{t}}
=\exp\!\Bigl(-\frac{\card{A}}{n}t+o(1)\Bigr)
=\exp\!\bigl(-\card{A}(\log n+c)+o(1)\bigr).
\]
代入定理 \ref{thm:pom-microcanonical-cover-time-tail-inclusion-exclusion} 的包含--排除表示并按 \(\card{A}=r\) 分组，得到
\[
\mathbb{P}(\tau>t_n(c))
=\sum_{r=1}^n(-1)^{r+1}\binom{n}{r}\exp\!\bigl(-r(\log n+c)+o(1)\bigr)
=1-\Bigl(1-\frac{e^{-c}}{n}+o\!\bigl(\tfrac1n\bigr)\Bigr)^n+o(1),
\]
从而 \(\mathbb{P}(\tau\le t_n(c))\to \exp(-e^{-c})\)。证毕。
\end{proof}

\begin{theorem}[缺失类型数的 Poisson 极限：cutoff 的更强刻画]\label{thm:pom-microcanonical-missing-types-poisson}
\leanverified{paper\_pom\_microcanonical\_missing\_types\_poisson}
在定理 \ref{thm:pom-microcanonical-cover-time-gumbel-cutoff} 的假设下，取 \(t=t_n(c)\) 并记缺失类型数为 \(M(t)\)（定义 \ref{def:pom-microcanonical-section-cover-time}）。
则对任意固定 \(c\in\RR\)，有分布收敛
\[
\boxed{\;
M(t)\xrightarrow{d}\mathrm{Poisson}(e^{-c})\;},
\qquad
\mathbb{P}(\tau\le t)=\mathbb{P}(M(t)=0)\to e^{-e^{-c}}.
\]
\end{theorem}
\begin{proof}
对任意固定整数 \(r\ge 1\)，考虑下降阶乘矩
\[
\EE\bigl[(M(t))_r\bigr]
=\sum_{\substack{A\subseteq[n]\\ \card{A}=r}}
\mathbb{P}\bigl(\text{\(A\) 中全部类型在前 \(t\) 步缺失}\bigr)
=\sum_{\card{A}=r}\frac{\binom{N-d(A)}{t}}{\binom{N}{t}}.
\]
由定理 \ref{thm:pom-microcanonical-cover-time-gumbel-cutoff} 证明中的一致近似，
当 \(\card{A}=r\) 固定时
\(
\binom{N-d(A)}{t}/\binom{N}{t}
=\exp(-rt/n+o(1))
=\exp(-r(\log n+c)+o(1))
\)。
故
\[
\EE\bigl[(M(t))_r\bigr]
=\binom{n}{r}n^{-r}e^{-rc}(1+o(1))\to \frac{e^{-rc}}{r!}.
\]
Poisson 分布由全部下降阶乘矩刻画，从而 \(M(t)\Rightarrow\mathrm{Poisson}(e^{-c})\)。
最后一式由 \(\mathbb{P}(\tau\le t)=\mathbb{P}(M(t)=0)\) 得到。证毕。
\end{proof}

\endinput


\begin{theorem}[剩余 \(t\) 点的 Bayes 最优精确识别率：超几何剩余计数闭式]\label{thm:pom-microcanonical-fold-bayes-success-Nminus-t}
\leanverified{paper\_pom\_microcanonical\_fold\_bayes\_success\_nminus\_t}
令 \(F\sim\mathrm{Unif}(\mathfrak{Fold}_m(d))\)。
考虑任意查询协议，在不重复的前提下查询 \(k=N-t\) 个点并观测全部响应；记未查询集 \(U\) 的大小为 \(\card{U}=t\)。
令剩余计数向量 \(R=(R_x)_{x\in X_m}\) 定义为
\[
R_x:=\card{\{\omega\in U:\ F(\omega)=x\}},\qquad \sum_{x\in X_m}R_x=t.
\]
则在给定全部观测的条件下，\(U\) 上所有满足剩余计数 \(R\) 的标签赋值在后验下等可能，其数目为 \(t!/\prod_x R_x!\)。
因此任何算法在该轨迹下精确输出整张映射 \(F\) 的最优成功概率为 \(\prod_x R_x!/t!\)，从而
\[
\boxed{\;
\mathrm{Succ}^{\mathrm{Bayes}}_{N-t}(d)
=\EE\Bigl[\frac{\prod_{x\in X_m}R_x!}{t!}\Bigr]\;},
\]
其中 \(R\) 服从由总体计数 \(d\) 与样本量 \(t\) 参数化的多元超几何分布。
并且存在等价闭式
\[
\boxed{\;
\mathrm{Succ}^{\mathrm{Bayes}}_{N-t}(d)
=\frac{1}{(N)_t}\sum_{\substack{r_x\in\ZZ_{\ge 0}\\ \sum_x r_x=t}}\ \prod_{x\in X_m}\bigl(d(x)\bigr)_{r_x}\; } .
\]
\end{theorem}
\begin{proof}
在固定 \(Q=\Omega_m\setminus U\) 上的全部响应之后，\(U\) 上唯一尚未被决定的是“把多重集 \(\{x^{\times R_x}\}\) 赋给 \(t\) 个已命名位置”的排列；
可行赋值数为 \(t!/\prod_x R_x!\)。
由于先验在 \(\mathfrak{Fold}_m(d)\) 上均匀，且对 \(U\) 内的任意置换都保持不变，故在给定观测与 \(R\) 的条件下，上述 \(t!/\prod_x R_x!\) 个赋值在后验下等可能；
于是最优（MAP）成功概率为其倒数 \(\prod_x R_x!/t!\)，从而得第一式。
\par
对闭式：对任意 \((r_x)_x\) 满足 \(\sum_x r_x=t\)，由超几何分布
\[
\mathbb{P}(R=r)=\frac{\prod_{x\in X_m}\binom{d(x)}{r_x}}{\binom{N}{t}}
\]
得
\[
\EE\Bigl[\frac{\prod_x R_x!}{t!}\Bigr]
=\sum_{r:\ \sum r_x=t}\frac{\prod_x r_x!}{t!}\cdot \frac{\prod_x\binom{d(x)}{r_x}}{\binom{N}{t}}
=\frac{1}{t!\binom{N}{t}}\sum_{r:\ \sum r_x=t}\ \prod_x\Bigl(r_x!\binom{d(x)}{r_x}\Bigr).
\]
用恒等式 \(r!\binom{d}{r}=(d)_r\) 与 \(t!\binom{N}{t}=(N)_t\) 化简即得。证毕。
\end{proof}

\paragraph{比例查询下的后验模空间指数律：线性熵收缩与典型不可恢复阈值}
在固定纤维谱的微正则语义下，查询仅在“移除若干已命名位置并揭示其标签”的意义上提供信息；
当查询次数与 \(N\) 同阶时，后验可行实现数的对数规模以 \((1-\beta)\) 的比例线性收缩。

\begin{definition}[比例查询的后验模空间规模]\label{def:pom-microcanonical-posterior-moduli-size-beta}
在定义 \ref{def:pom-microcanonical-adaptive-trajectory-counts} 的设定下，给定轨迹 \(t\in X_m^k\) 与其计数 \(c_t\)，定义后验可行实现数（后验模空间）
\[
\mathcal{M}_k(t)
:=\card{\Bigl\{F\in\mathfrak{Fold}_m(d):\ F(\omega_i)=t_i,\ 1\le i\le k\Bigr\}}
=\frac{(N-k)!}{\prod_{x\in X_m}\bigl(d(x)-c_t(x)\bigr)!},
\]
其中若存在 \(x\) 使 \(c_t(x)>d(x)\) 则按 \(0\) 解释。
\end{definition}

\begin{theorem}[后验熵的线性收缩率：比例查询下的确定性指数律]\label{thm:pom-microcanonical-posterior-entropy-linear-law}
\leanverified{paper\_pom\_microcanonical\_posterior\_entropy\_linear\_law}
设 \(n=\card{X_m}\ge 2\)，并令 \(w(x):=d(x)/N\in\Delta(X_m)\)。
取常数 \(\beta\in(0,1)\) 并令 \(k=\lfloor \beta N\rfloor\)。
则在 \(N\to\infty\) 且 \(w\) 固定时，有依概率收敛
\[
\boxed{\;
\frac{1}{N}\log \mathcal{M}_k(T_{1:k})
\ \xrightarrow{\ \mathbb{P}\ }\ (1-\beta)\,H(w)\;},
\qquad
H(w):=-\sum_{x\in X_m}w(x)\log w(x).
\]
\end{theorem}
\begin{proof}
由定理 \ref{thm:pom-microcanonical-adaptive-no-gain-without-replacement}，
\((T_i)\) 等价于从多重集 \(\{x^{\times d(x)}\}\) 的无放回抽样；
因此对每个 \(x\)，计数 \(C_k(x)\) 为超几何随机变量，满足
\(
\EE[C_k(x)]=kw(x)
\)
且 \(\Var(C_k(x))=O(N)\)。
由 Chebyshev 不等式得到
\(
C_k(x)/N\to \beta w(x)
\)
（对每个 \(x\)）并由有限维性推出向量收敛 \(C_k/N\to \beta w\) 依概率成立；
于是剩余计数向量
\(
R:=d-C_k
\)
满足
\(
R/(N-k)\to w
\)
依概率成立。
\par
对 \(\log\mathcal{M}_k=\log(N-k)!-\sum_x\log R_x!\) 作 Stirling 展开：
\[
\log (N-k)!
=(N-k)\log(N-k)-(N-k)+O(\log N),
\]
\[
\sum_x\log R_x!
=\sum_x\bigl(R_x\log R_x-R_x\bigr)+O(n\log N).
\]
合并并除以 \(N\) 得
\[
\frac{1}{N}\log\mathcal{M}_k
=\frac{N-k}{N}\Bigl(\log(N-k)-\sum_x \frac{R_x}{N-k}\log R_x\Bigr)+o(1)
=(1-\beta)\,H\!\left(\frac{R}{N-k}\right)+o(1),
\]
其中 \(o(1)\) 使用了 \(n\) 固定（或更一般地 \(n\log N=o(N)\)）以及 \(\log(N-k)=O(\log N)=o(N)\)。
由 \(R/(N-k)\to w\) 与 \(H(\cdot)\) 的连续性，得到极限 \((1-\beta)H(w)\)。证毕。
\end{proof}

\begin{theorem}[对数后验模空间的高斯涨落：比例查询下的中心极限定理]\label{thm:pom-microcanonical-posterior-moduli-clt}
\leanverified{paper\_pom\_microcanonical\_posterior\_moduli\_clt}
在定理 \ref{thm:pom-microcanonical-posterior-entropy-linear-law} 的设定下，进一步假设字母表大小 \(n=\card{X_m}\) 固定，且 \(w\) 固定并满足 \(w(x)\in(0,1)\)。
令
\[
V(w):=\Var_{X\sim w}\bigl(\log w(X)\bigr)
=\sum_{x\in X_m}w(x)\bigl(\log w(x)\bigr)^2-\Bigl(\sum_{x\in X_m}w(x)\log w(x)\Bigr)^2.
\]
则当 \(N\to\infty\) 时，
\[
\boxed{\;
\frac{\log \mathcal{M}_k(T_{1:k})-(1-\beta)NH(w)}{\sqrt{\beta(1-\beta)N}}
\ \xrightarrow{d}\ \mathcal N\bigl(0,V(w)\bigr)\;}
\]
其中 \(k=\lfloor\beta N\rfloor\) 且 \(H(w)=-\sum_x w(x)\log w(x)\)。
\end{theorem}
\begin{proof}
沿用定理 \ref{thm:pom-microcanonical-posterior-entropy-linear-law} 的记号，令 \(R=d-C_k\) 且 \(\widehat w:=R/(N-k)\)。
当 \(n\) 固定时 Stirling 展开给出一致误差界
\[
\log\mathcal{M}_k=(N-k)H(\widehat w)+O(1).
\]
无放回抽样下的多元超几何中心极限定理给出
\[
\sqrt N\,(\widehat w-w)\ \xrightarrow{d}\ \mathcal N\!\left(0,\frac{\beta}{1-\beta}\bigl(\mathrm{diag}(w)-ww^{\mathsf T}\bigr)\right).
\]
对熵函数 \(H\) 作 delta 方法：在约束 \(\sum_x \widehat w(x)=1\) 下，
\[
H(\widehat w)-H(w)=-\sum_{x\in X_m}\bigl(\widehat w(x)-w(x)\bigr)\log w(x)+o_p(N^{-1/2}),
\]
故
\[
\sqrt N\bigl(H(\widehat w)-H(w)\bigr)\ \xrightarrow{d}\ \mathcal N\!\left(0,\frac{\beta}{1-\beta}V(w)\right).
\]
再乘以 \(N-k=(1-\beta)N\) 并吸收 \(O(1)\) 余项（相对 \(\sqrt N\) 为 \(o(\sqrt N)\)），即得所示结论。证毕。
\end{proof}

\begin{corollary}[典型后验精确恢复成功率的双指数衰减]\label{cor:pom-microcanonical-posterior-map-success-double-exponential}
在定理 \ref{thm:pom-microcanonical-posterior-entropy-linear-law} 的设定下，
由定理 \ref{thm:pom-microcanonical-count-sufficient-statistic-posterior-uniform} 的后验均匀性，
给定轨迹 \(T_{1:k}\) 时精确输出整张映射 \(F\) 的 Bayes 最优条件成功概率恰为
\[
\mathrm{Succ}^{\mathrm{Bayes}}\bigl(F\mid T_{1:k}\bigr)=\frac{1}{\mathcal{M}_k(T_{1:k})}.
\]
因此依概率成立
\[
\boxed{\;
\frac{1}{N}\log \mathrm{Succ}^{\mathrm{Bayes}}\bigl(F\mid T_{1:k}\bigr)
\ \xrightarrow{\ \mathbb{P}\ }\ -(1-\beta)\,H(w)\; } .
\]
特别地，在本文主线 \(N=2^m\) 下，上式给出典型条件成功率在 \(m\) 上的双指数衰减
\(
\mathrm{Succ}^{\mathrm{Bayes}}(F\mid T_{1:k})
=\exp(-(1-\beta)H(w)\,2^m+o(2^m))
\)。
\end{corollary}

\begin{remark}[整体 Bayes 成功率与典型指数率的区分]\label{rem:pom-microcanonical-global-vs-typical-bayes-success}
由定义，整体 Bayes 成功率为 \(\EE[\mathcal{M}_k(T_{1:k})^{-1}]\)。
由于后验在每条轨迹的兼容集合上均匀，
对任意固定 \(k\) 也有严格恒等式
\[
\EE\!\left[\frac{1}{\mathcal{M}_k(T_{1:k})}\right]
=\frac{\card{\{t\in X_m^k:\ c_t(x)\le d(x)\ \forall x\}}}{\card{\mathfrak{Fold}_m(d)}}.
\]
该恒等式表明整体成功率仅由“可实现轨迹的总数/微正则实现类大小”的比值决定；
而推论 \ref{cor:pom-microcanonical-posterior-map-success-double-exponential} 刻画的是典型轨迹下的条件指数率。
\end{remark}

\endinput


\paragraph{信息线性守恒与涨落：比例查询下的互信息、中心极限定理与大偏差}
在微正则先验 \(F\sim\mathrm{Unif}(\mathfrak{Fold}_m(d))\) 下，
任意确定性自适应互异查询产生的类型轨迹等价于从多重集 \(\{x^{\times d(x)}\}_{x\in X_m}\) 的无放回抽样（定理 \ref{thm:pom-microcanonical-adaptive-no-gain-without-replacement}）。
因此，“信息获得”可在不涉及几何细节的前提下由纯组合量刻画：后验模空间 \(\mathcal{M}_k\) 的对数收缩即为信息密度；
其热力学极限在指数尺度上线性，并在 \(\sqrt N\) 尺度呈现高斯涨落；同时可由多元超几何的 KL--对称大偏差率函数给出全大偏差描述。

\begin{definition}[信息密度与归一化信息率]\label{def:pom-microcanonical-information-density}
在定义 \ref{def:pom-microcanonical-adaptive-trajectory-counts} 与定义 \ref{def:pom-microcanonical-posterior-moduli-size-beta} 的设定下，
令
\[
\mathsf{I}_k:=\log\card{\mathfrak{Fold}_m(d)}-\log \mathcal{M}_k(T_{1:k})
\]
为\emph{信息密度}（亦即观测轨迹的 self-information，见定理 \ref{thm:pom-microcanonical-information-identity}）。
并定义归一化量
\[
\iota_N:=\frac{1}{N}\mathsf{I}_k,
\qquad
\ell_N:=\frac{1}{N}\log \mathcal{M}_k(T_{1:k}).
\]
\end{definition}

\begin{theorem}[信息密度的精确模空间表示]\label{thm:pom-microcanonical-information-identity}
\leanverified{paper\_pom\_microcanonical\_information\_identity}
在上述微正则先验与确定性自适应互异查询的设定下，对任意 \(k\le N\) 有严格恒等式
\[
\boxed{\;
\mathsf{I}_k
=-\log\mathbb{P}(T_{1:k})
=\log\card{\mathfrak{Fold}_m(d)}-\log \mathcal{M}_k(T_{1:k})\;}
\]
并且互信息满足
\[
\boxed{\;
I(F;T_{1:k})=H(T_{1:k})
=\EE[\mathsf{I}_k]
=\log\card{\mathfrak{Fold}_m(d)}-\EE\bigl[\log \mathcal{M}_k(T_{1:k})\bigr]\;}
\]
其中 \(H(T_{1:k})\) 为轨迹分布的 Shannon 熵（自然对数）。
\end{theorem}
\begin{proof}
由于查询策略确定且每次查询点由历史响应唯一确定，轨迹 \(T_{1:k}\) 是 \(F\) 的确定性函数，故 \(I(F;T_{1:k})=H(T_{1:k})\)。
\par
另一方面，由定理 \ref{thm:pom-microcanonical-fold-posterior-count-and-prob} 的轨迹概率公式可知：
对任意可实现轨迹 \(t\in X_m^k\)，其出现概率恰为
\[
\mathbb{P}(T_{1:k}=t)=\frac{\card{\{F'\in\mathfrak{Fold}_m(d):F'(\omega_i)=t_i,\ 1\le i\le k\}}}{\card{\mathfrak{Fold}_m(d)}}
=\frac{\mathcal{M}_k(t)}{\card{\mathfrak{Fold}_m(d)}}.
\]
取 \(t=T_{1:k}\) 即得
\(
-\log\mathbb{P}(T_{1:k})
=\log\card{\mathfrak{Fold}_m(d)}-\log\mathcal{M}_k(T_{1:k})
\)，从而得到第一式。
对两边取期望并注意 \(\EE[-\log\mathbb{P}(T_{1:k})]=H(T_{1:k})\)，得到第二组恒等式。证毕。
\end{proof}

\begin{theorem}[信息线性守恒律：比例查询下的指数率]\label{thm:pom-microcanonical-information-linear-law}
取常数 \(\beta\in[0,1]\) 并令 \(k=\lfloor\beta N\rfloor\)。
则在 \(N\to\infty\) 且 \(w\) 固定时，有依概率收敛
\[
\boxed{\;
\frac{1}{N}\mathsf{I}_k\ \xrightarrow{\ \mathbb{P}\ }\ \beta\,H(w)\;},
\qquad
\boxed{\;
\frac{1}{k}\mathsf{I}_k\ \xrightarrow{\ \mathbb{P}\ }\ H(w)\;}
\quad(\beta\in(0,1]),
\]
其中 \(w(x)=d(x)/N\) 且 \(H(w)=-\sum_x w(x)\log w(x)\)。
此外，互信息满足同一指数率
\[
\frac{1}{N}I(F;T_{1:k})=\frac{1}{N}\EE[\mathsf{I}_k]\ \to\ \beta\,H(w).
\]
\leanverified{paper\_pom\_microcanonical\_information\_linear\_law}
\end{theorem}
\begin{proof}
由定理 \ref{thm:pom-microcanonical-information-identity}，
\(
\mathsf{I}_k=\log\card{\mathfrak{Fold}_m(d)}-\log\mathcal{M}_k(T_{1:k})
\)。
由定理 \ref{thm:pom-microcanonical-fold-entropy}，
\[
\frac{1}{N}\log\card{\mathfrak{Fold}_m(d)}=H(w)+o(1),
\]
而由定理 \ref{thm:pom-microcanonical-posterior-entropy-linear-law} 有
\[
\frac{1}{N}\log\mathcal{M}_k(T_{1:k})\ \xrightarrow{\ \mathbb{P}\ }\ (1-\beta)\,H(w).
\]
两式相减即得 \(\frac{1}{N}\mathsf{I}_k\to \beta H(w)\) 依概率收敛。
再除以 \(k/N\to\beta\) 得第二式。
最后一式由 \(\EE[\mathsf{I}_k]=I(F;T_{1:k})\) 与 \(\mathsf{I}_k/N\) 的一致有界性推出：由于 \(\log\mathcal{M}_k(T_{1:k})\ge 0\)，有
\[
0\le \frac{1}{N}\mathsf{I}_k\le \frac{1}{N}\log\card{\mathfrak{Fold}_m(d)}=H(w)+o(1),
\]
从而 \(\mathsf{I}_k/N\) 在 \(L^1\) 上一致有界；结合依概率收敛可得
\(
\EE[\mathsf{I}_k]/N\to \beta H(w)
\)。
证毕。
\end{proof}

\begin{corollary}[互信息涨落的中心极限定理]\label{cor:pom-microcanonical-information-clt}
\leanverified{paper\_pom\_microcanonical\_information\_clt}
在定理 \ref{thm:pom-microcanonical-posterior-moduli-clt} 的设定下（特别地 \(n\) 固定、\(w\) 固定且 \(\beta\in(0,1)\)），
记
\[
V(w):=\Var_{X\sim w}\bigl(\log w(X)\bigr).
\]
则当 \(N\to\infty\) 时，
\[
\boxed{\;
\frac{\mathsf{I}_k-\beta N H(w)}{\sqrt{\beta(1-\beta)N}}
\ \xrightarrow{d}\ \mathcal N\bigl(0,V(w)\bigr)\;}
\qquad\left(k=\lfloor\beta N\rfloor\right).
\]
\end{corollary}
\begin{proof}
由定义 \ref{def:pom-microcanonical-information-density} 与定理 \ref{thm:pom-microcanonical-fold-entropy}，
\[
\mathsf{I}_k
=\log\card{\mathfrak{Fold}_m(d)}-\log\mathcal{M}_k(T_{1:k})
=NH(w)-\log\mathcal{M}_k(T_{1:k})+O(n\log N).
\]
将定理 \ref{thm:pom-microcanonical-posterior-moduli-clt} 的中心极限定理代入并注意在 \(n\) 固定时 \(O(n\log N)=o(\sqrt N)\)，即得结论。证毕。
\end{proof}

\begin{theorem}[多元超几何的大偏差：样本组成与剩余组成的 KL--对称分解]\label{thm:pom-microcanonical-hypergeometric-ldp-kl-sym}
\leanverified{paper\_pom\_microcanonical\_hypergeometric\_ldp\_kl\_sym}
假设 \(n\) 固定、\(w\) 固定且 \(\beta\in(0,1)\)。
令 \(k=\lfloor\beta N\rfloor\)，并定义样本组成与剩余组成
\[
Q_N(x):=\frac{C_k(x)}{k}\in\Delta(X_m),
\qquad
R_N(x):=\frac{d(x)-C_k(x)}{N-k}\in\Delta(X_m).
\]
则 \(Q_N\) 满足大偏差原理，其良速率函数可显式写为
\[
J_\beta(q)=
\begin{cases}
\beta\,D(q\|w)+(1-\beta)\,D(r(q)\|w), & r(q):=\dfrac{w-\beta q}{1-\beta}\in\Delta(X_m),\\[6pt]
+\infty, & \text{否则},
\end{cases}
\]
其中 \(D(\cdot\|\cdot)\) 为 KL 散度（自然对数）。
并且具有对称性
\[
\boxed{\;
J_\beta(q)=J_{1-\beta}\bigl(r(q)\bigr)\;}
\]
对应于质量守恒 \(w=\beta q+(1-\beta)r\) 的交换不变性。
\end{theorem}
\begin{proof}
由超几何分布的概率质量函数（定理 \ref{thm:pom-microcanonical-adaptive-no-gain-without-replacement} 的计数形式），
对整数向量 \(c\) 满足 \(\sum_x c(x)=k\) 有
\[
\mathbb{P}(C_k=c)=\frac{\prod_{x\in X_m}\binom{d(x)}{c(x)}}{\binom{N}{k}}.
\]
取 \(c(x)=kq(x)\) 并对阶乘用 Stirling 公式展开，得到
\[
-\frac1N\log \mathbb{P}(Q_N\approx q)
=\beta\,D(q\|w)+(1-\beta)\,D(r(q)\|w)+o(1),
\]
其中 \(r(q)=(w-\beta q)/(1-\beta)\) 由守恒关系强制出现。
该展开在 \(q\) 落在单纯形内点且 \(r(q)\in\Delta(X_m)\) 时成立并给出良速率函数；若 \(r(q)\notin\Delta\) 则概率为 \(0\)，速率为 \(+\infty\)。
对称性由把 \(\beta q\) 与 \((1-\beta)r\) 交换而保持 \(w\) 不变直接推出。证毕。
\end{proof}

\begin{theorem}[后验模空间与信息密度的全大偏差：收缩原理的显式变分式]\label{thm:pom-microcanonical-moduli-information-ldp}
\leanverified{paper\_pom\_microcanonical\_moduli\_information\_ldp}
在定理 \ref{thm:pom-microcanonical-hypergeometric-ldp-kl-sym} 的假设下，\(\ell_N\) 与 \(\iota_N\) 均满足大偏差原理。
其速率函数可写为
\[
K_\ell(s)=\inf\Bigl\{J_\beta(q):\ s=(1-\beta)\,H(r(q))\Bigr\},
\]
\[
K_\iota(s)=\inf\Bigl\{J_\beta(q):\ s=H(w)-(1-\beta)\,H(r(q))\Bigr\},
\]
其中 \(r(q)=(w-\beta q)/(1-\beta)\)，且 \(H(\cdot)\) 为 Shannon 熵（自然对数）。
最小点对应 \(q=w\)（从而 \(r(w)=w\)），并给出典型值
\[
\ell_N\ \to\ (1-\beta)H(w),
\qquad
\iota_N\ \to\ \beta H(w),
\]
与定理 \ref{thm:pom-microcanonical-information-linear-law} 一致。
\end{theorem}
\begin{proof}
在事件 \(\{Q_N\approx q\}\) 上，剩余组成满足 \(R_N\approx r(q)\)。
由 Stirling 展开（见定理 \ref{thm:pom-microcanonical-posterior-entropy-linear-law} 的同型推导）得到
\[
\frac{1}{N}\log\mathcal{M}_k(T_{1:k})
=(1-\beta)H(R_N)+o(1)
=(1-\beta)H(r(q))+o(1).
\]
因此在有效域内，映射 \(q\mapsto (1-\beta)H(r(q))\) 连续。
对 \(Q_N\) 的大偏差原理应用收缩原理得到 \(K_\ell\)。
由定理 \ref{thm:pom-microcanonical-fold-entropy}，有 \(\frac{1}{N}\log\card{\mathfrak{Fold}_m(d)}=H(w)+o(1)\)。
因此
\[
\iota_N=\frac{1}{N}\log\card{\mathfrak{Fold}_m(d)}-\ell_N
=H(w)-\ell_N+o(1),
\]
而确定性 \(o(1)\) 平移不改变 \(N\) 尺度的大偏差率函数；
从而由 \(K_\ell\) 立得 \(K_\iota\) 的变分式（亦可视为 \(K_\iota(s)=K_\ell(H(w)-s)\) 的等价写法）。
最小点由 \(J_\beta\) 的严格凸性（在内点上）与 \(J_\beta(w)=0\) 给出。证毕。
\end{proof}

\begin{theorem}[熵--信息的 Doob 分解：条件增量与非渐近浓缩]\label{thm:pom-microcanonical-information-doob}
\leanverified{paper\_pom\_microcanonical\_information\_doob}
对任意 \(i\in\{0,1,\dots,k\}\) 定义剩余组成
\[
p_i(x):=\frac{d(x)-C_i(x)}{N-i}\in\Delta(X_m),
\qquad C_0\equiv 0,
\]
并令滤过 \(\mathcal F_i:=\sigma(T_1,\dots,T_i)\)。
则成立：
\begin{enumerate}
  \item \textbf{（条件增量）}对每个 \(1\le i\le k\)，有
  \[
  \mathbb{P}(T_i=x\mid \mathcal F_{i-1})=p_{i-1}(x),
  \qquad
  \mathsf{I}_k=\sum_{i=1}^k\bigl(-\log p_{i-1}(T_i)\bigr)\ \ \text{a.s.}
  \]
  且
  \[
  \EE\!\left[-\log\mathbb{P}(T_i\mid\mathcal F_{i-1})\ \big|\ \mathcal F_{i-1}\right]=H(p_{i-1}).
  \]
  \item \textbf{（鞅结构）}\((p_i)_{i\ge 0}\) 为 \(\Delta(X_m)\)-值鞅：\(\EE[p_i\mid\mathcal F_{i-1}]=p_{i-1}\)。
  \item \textbf{（Azuma 型浓缩）}令
  \[
  M_k:=\mathsf{I}_k-\sum_{i=1}^k H(p_{i-1}),
  \]
  则 \((M_i)\) 为鞅且其差分满足 \(|M_i-M_{i-1}|\le \log N\)。
  因而对任意 \(t>0\) 有非渐近界
  \[
  \boxed{\;
  \mathbb{P}\Bigl(\bigl|\mathsf{I}_k-\sum_{i=1}^k H(p_{i-1})\bigr|\ge t\Bigr)
  \le 2\exp\!\Bigl(-\frac{2t^2}{k(\log N)^2}\Bigr)\; } .
  \]
\end{enumerate}
\leanverified{paper\_pom\_microcanonical\_query\_plus\_bits\_phase}
\end{theorem}
\begin{proof}
第 (1) 点由定理 \ref{thm:pom-microcanonical-adaptive-no-gain-without-replacement} 的逐步条件分布公式直接得到。
由链式法则，
\(
\mathbb{P}(T_{1:k})=\prod_{i=1}^k\mathbb{P}(T_i\mid\mathcal F_{i-1})
\)，取负对数即得 \(\mathsf{I}_k=\sum_i -\log p_{i-1}(T_i)\)。
条件期望恒等式由
\(
\EE[-\log p_{i-1}(T_i)\mid\mathcal F_{i-1}]
=\sum_x p_{i-1}(x)(-\log p_{i-1}(x))
=H(p_{i-1})
\)
给出。
\par
第 (2) 点：对任意 \(x\)，由 \(C_i(x)=C_{i-1}(x)+\mathbf{1}_{\{T_i=x\}}\) 与
\(\EE[\mathbf{1}_{\{T_i=x\}}\mid\mathcal F_{i-1}]=p_{i-1}(x)\) 得
\[
\EE[p_i(x)\mid\mathcal F_{i-1}]
=\EE\!\left[\frac{d(x)-C_{i-1}(x)-\mathbf{1}_{\{T_i=x\}}}{N-i}\ \Big|\ \mathcal F_{i-1}\right]
=\frac{(N-i+1)p_{i-1}(x)-p_{i-1}(x)}{N-i}
=p_{i-1}(x).
\]
\par
第 (3) 点：由第 (1) 点，\(M_i-M_{i-1}=-\log p_{i-1}(T_i)-H(p_{i-1})\) 条件均值为 \(0\)，故 \(M_i\) 为鞅。
并且由于 \(p_{i-1}(T_i)\ge 1/(N-i+1)\)（观测到的类型在剩余多重集中至少出现一次），
得 \(-\log p_{i-1}(T_i)\le \log(N-i+1)\le \log N\)，而 \(0\le H(p_{i-1})\le \log n\le \log N\)；
故 \(|M_i-M_{i-1}|\le \log N\)。
对 \((M_i)\) 应用 Azuma--Hoeffding 即得所示界。证毕。
\end{proof}

\begin{theorem}[查询 + 外置信息比特：最优恢复率与高斯临界窗]\label{thm:pom-microcanonical-query-plus-bits-phase}
在给定 \(k\) 次查询轨迹 \(T_{1:k}\) 后，允许额外获得 \(B\) 比特的外置信息：
编码器可取任意映射 \(\mathrm{enc}:\mathfrak{Fold}_m(d)\to\{0,1\}^B\)，解码器以 \((T_{1:k},\mathrm{enc}(F))\) 为输入输出 \(\widehat F\)。
\begin{enumerate}
  \item \textbf{（最优成功率闭式）}在均匀先验下，最优精确恢复成功率满足严格恒等式
  \[
  \boxed{\;
  \mathrm{Succ}^{\mathrm{opt}}_{k,B}
  =\EE\!\left[\min\!\left(1,\frac{2^B}{\mathcal{M}_k(T_{1:k})}\right)\right]\;}
  \]
  其中 \(\mathcal{M}_k(T_{1:k})\) 为后验兼容集合大小（定义 \ref{def:pom-microcanonical-posterior-moduli-size-beta}）。
  \item \textbf{（阈值与高斯窗）}
  假设 \(n\) 固定、\(w\) 固定且 \(\beta\in(0,1)\)，令 \(k=\lfloor\beta N\rfloor\) 并取
  \[
  B_N=\frac{(1-\beta)NH(w)}{\log 2}
  +\frac{t}{\log 2}\sqrt{\beta(1-\beta)N V(w)},
  \qquad t\in\RR.
  \]
  则当 \(N\to\infty\) 时，
  \[
  \boxed{\;
  \mathrm{Succ}^{\mathrm{opt}}_{k,B_N}\ \longrightarrow\ \Phi(t)\;}
  \]
  其中 \(\Phi\) 为标准正态分布函数。
  特别地，若 \(B<(1-\beta)NH(w)/\log 2-\omega(\sqrt N)\)，则 \(\mathrm{Succ}^{\mathrm{opt}}_{k,B}\to 0\)；
  若 \(B>(1-\beta)NH(w)/\log 2+\omega(\sqrt N)\)，则 \(\mathrm{Succ}^{\mathrm{opt}}_{k,B}\to 1\)。
\end{enumerate}
\end{theorem}
\begin{proof}
第 (1) 点：在给定轨迹 \(T_{1:k}\) 时，后验在兼容集合上均匀（定理 \ref{thm:pom-microcanonical-count-sufficient-statistic-posterior-uniform}），其大小为 \(\mathcal{M}_k(T_{1:k})\)。
任意 \(B\) 比特编码至多区分 \(2^B\) 个候选映射；条件于轨迹后，最优做法是在后验块内选择至多 \(2^B\) 个映射赋予互异码字，并将其余映射合并到任意码字上。
故条件成功率的最优值为 \(\min(1,2^B/\mathcal{M}_k)\)；对轨迹取期望即得结论。
\par
第 (2) 点：由定理 \ref{thm:pom-microcanonical-posterior-moduli-clt}，
\[
\frac{\log \mathcal{M}_k-(1-\beta)NH(w)}{\sqrt{\beta(1-\beta)N}}
\ \xrightarrow{d}\ \mathcal N\bigl(0,V(w)\bigr).
\]
令 \(B_N\log 2=(1-\beta)NH(w)+t\sqrt{\beta(1-\beta)NV(w)}\)。
则
\[
\min\!\left(1,\frac{2^{B_N}}{\mathcal{M}_k}\right)
=\min\!\left(1,\exp\!\bigl(B_N\log 2-\log\mathcal{M}_k\bigr)\right)
\xrightarrow{d}\ \mathbf{1}_{\{Z\le t\}},
\]
其中 \(Z\sim\mathcal N(0,1)\)。
由于左端始终落在 \([0,1]\)，由支配收敛（或 Portmanteau + 有界性）取期望得到
\(
\mathrm{Succ}^{\mathrm{opt}}_{k,B_N}\to \mathbb{P}(Z\le t)=\Phi(t)
\)。
两侧强相变由 \(t\to\pm\infty\) 得到。证毕。
\end{proof}

% (User input window)

\paragraph{定型类内 Hamming 球的指数体积与部分反演阈值：对角约束最大熵耦合}
微正则后验在“仅固定类型计数/边缘组成”的条件下呈现典型的定型类结构：在给定边缘分布 \(w\) 时，
未查询标签向量在长度 \(t\) 的定型类上均匀。
因此，允许 \(\delta\)-失真时的最优“部分反演”成功率可还原为一个纯组合覆盖问题：
定型类由多少个定型内 Hamming 球即可在指数尺度上覆盖。
本段给出该覆盖率的精确指数阈值，并将其表示为一个对角约束下的最大熵耦合问题；
进一步证明极大耦合的唯一性、指数族结构、标量化求解、以及互信息/对偶导数刻画。

\begin{definition}[定型类与定型内 Hamming 球]\label{def:pom-typeclass-and-hamming-ball}
令 \(X\) 为有限字母表且 \(|X|=n\ge 2\)，取全支撑分布 \(w\in\Delta(X)\)。
令 \(t\to\infty\) 且满足 \(t\,w(x)\in\NN\)（对所有 \(x\in X\)）。
定义定型类
\[
\mathcal T_t(w):=\Bigl\{x^t\in X^t:\ \frac{1}{t}\bigl|\{i:\ x_i=a\}\bigr|=w(a)\ \ \forall a\in X\Bigr\}.
\]
对任意参考序列 \(x^t\in\mathcal T_t(w)\) 与任意失真参数 \(\delta\in[0,1]\)，定义定型内球
\[
\mathcal B_t(w,\delta;x^t):=\Bigl\{y^t\in\mathcal T_t(w):\ t^{-1}d_{\rm H}(x^t,y^t)\le \delta\Bigr\}.
\]
\end{definition}

\begin{definition}[指数体积与覆盖阈值]\label{def:pom-typeclass-volume-and-threshold}
若极限存在且与参考序列 \(x^t\in\mathcal T_t(w)\) 无关，则定义指数体积
\[
\mathcal V_w(\delta):=\lim_{t\to\infty}\frac{1}{t}\log|\mathcal B_t(w,\delta;x^t)|.
\]
并定义相应的覆盖阈值
\[
R_w(\delta):=H(w)-\mathcal V_w(\delta)\ \ge\ 0.
\]
\end{definition}

\begin{theorem}[指数体积的存在性与变分表示]\label{thm:pom-typeclass-hamming-ball-volume-variational}
在定义 \ref{def:pom-typeclass-and-hamming-ball} 的设定下，极限 \(\mathcal V_w(\delta)\) 存在，并满足严格变分式
\[
\boxed{\;
\mathcal V_w(\delta)
=
\sup\Bigl\{H_P(Y|X):\ P\in\Delta(X\times X),\ P_X=P_Y=w,\ \mathbb P_P(X=Y)\ge 1-\delta\Bigr\}\; }.
\]
并且该指数体积与参考序列 \(x^t\in\mathcal T_t(w)\) 的选择无关。
\end{theorem}
\begin{proof}
对任意 \(x^t\in\mathcal T_t(w)\) 与任意 \(y^t\in\mathcal T_t(w)\)，记其联合经验分布（contingency table）为 \(P_{x^t y^t}\in\Delta(X\times X)\)。
定型约束 \(y^t\in\mathcal T_t(w)\) 等价于 \(P_Y=w\)，而 \(x^t\in\mathcal T_t(w)\) 等价于 \(P_X=w\)。
Hamming 约束 \(t^{-1}d_{\rm H}(x^t,y^t)\le \delta\) 等价于 \(P(X=Y)\ge 1-\delta\)。
\par
对固定 \(x^t\)，当 \(y^t\) 的联合型为某个可行 \(P\) 时，
符合该联合型的 \(y^t\) 个数的指数率为 \(H_P(Y|X)\)（标准定型计数：\(|\{y^t:\ P_{x^t y^t}=P\}|=\exp(tH_P(Y|X)+o(t))\)）。
因此 \(|\mathcal B_t(w,\delta;x^t)|\) 在指数尺度上由可行集合上的最大 \(H_P(Y|X)\) 决定，从而得到所示上确界。
\par
由于可行集合在有限维单纯形内紧且 \(H_P(Y|X)\) 连续，最大值存在。
上界（\(\limsup\)）与下界（\(\liminf\)）由对有限个联合型的并集上界与选择近最优联合型的下界配对给出，从而极限存在。
最后，定型类在坐标置换下传递作用，\(|\mathcal B_t(w,\delta;x^t)|\) 仅依赖于 \(x^t\) 的类型而非具体序列，故与参考序列无关。证毕。
\end{proof}

\begin{theorem}[唯一极大耦合与指数族结构]\label{thm:pom-typeclass-hamming-ball-unique-maximizer-exponential-family}
对每个 \(\delta\in(0,1)\)，定理 \ref{thm:pom-typeclass-hamming-ball-volume-variational} 的上确界由唯一耦合 \(P^\star_\delta\) 达到。
存在唯一标量 \(\lambda(\delta)\in\RR\) 与唯一正向量 \(u(\delta)\in\RR_{>0}^{X}\) 使
\[
\boxed{\;
P^\star_\delta(x,y)
=
\frac{u_xu_y\,\exp\!\bigl(\lambda(\delta)\,\mathbf 1_{\{x=y\}}\bigr)}
{\sum_{a,b\in X}u_au_b\,\exp\!\bigl(\lambda(\delta)\,\mathbf 1_{\{a=b\}}\bigr)}\qquad(\forall x,y\in X)\; }.
\]
并且其边缘满足 \(P^\star_{\delta,X}=P^\star_{\delta,Y}=w\)。
记 \(p_0:=\sum_{x\in X}w(x)^2\)。
若 \(1-\delta>p_0\)，则对角质量约束饱和并满足 \(\mathbb P_{P^\star_\delta}(X=Y)=1-\delta\)，且 \(\lambda(\delta)>0\)。
若 \(1-\delta\le p_0\)，则最优耦合为独立耦合 \(P^\star_\delta=w\otimes w\)，并且 \(\lambda(\delta)=0\)、\(\mathcal V_w(\delta)=H(w)\)、\(R_w(\delta)=0\)。
\end{theorem}
\begin{proof}
可行集合由线性约束与不等式约束给出，为紧凸集。
函数 \(H_P(Y|X)=H(P)-H(w)\) 在 \(P\) 上严格凹（因 \(H(P)\) 严格凹且 \(H(w)\) 为常数），
故极大元唯一。
\par
对偶化（KKT）给出指数族形式：
\(
P(x,y)\propto \exp(\alpha_x+\beta_y+\lambda\mathbf 1_{\{x=y\}})
\)。
由于约束要求 \(P_X=P_Y=w\) 且问题在交换 \(X,Y\) 下对称，唯一性强制 \(\alpha=\beta\)（相差常数可并入归一化因子）。
令 \(u_x=e^{\alpha_x}>0\) 得到所示结构。
\par
令 \(P^{\mathrm{ind}}:=w\otimes w\)，则 \(H_{P^{\mathrm{ind}}}(Y|X)=H(w)\) 且 \(\mathbb P_{P^{\mathrm{ind}}}(X=Y)=p_0\)。
由于对任意耦合 \(P\) 恒有 \(H_P(Y|X)\le H(Y)=H(w)\)，当 \(1-\delta\le p_0\) 时 \(P^{\mathrm{ind}}\) 可行并达到全局上界；
因而由唯一性知 \(P^\star_\delta=P^{\mathrm{ind}}\)，并且指数族参数满足 \(\lambda(\delta)=0\)。
\par
当 \(1-\delta>p_0\) 时，独立耦合不可行。
若最优点处约束为严格不等式，则由互补松弛必有 \(\lambda(\delta)=0\)，从而仍为独立耦合，矛盾。
故约束必饱和，且 \(\lambda(\delta)>0\) 并满足 \(\mathbb P_{P^\star_\delta}(X=Y)=1-\delta\)。证毕。
\end{proof}

\begin{theorem}[部分反演的指数阈值：覆盖率由 \texorpdfstring{$R_w(\delta)$}{Rw(delta)} 决定]\label{thm:pom-microcanonical-partial-inversion-threshold}
设 \(Y^t\) 在 \(\mathcal T_t(w)\) 上均匀分布。
允许额外 \(B\) 比特外置信息：编码器取任意映射 \(\mathrm{enc}:\mathcal T_t(w)\to\{1,\dots,2^B\}\)，
解码器以 \(\mathrm{enc}(Y^t)\) 为输入输出 \(\widehat Y^t\in\mathcal T_t(w)\)。
对固定 \(\delta\in(0,1)\) 定义失真事件
\[
\mathcal E_\delta:=\{t^{-1}d_{\rm H}(Y^t,\widehat Y^t)\le \delta\},
\]
并令最优成功率
\[
\mathrm{Succ}^{\mathrm{opt}}_t(B,\delta):=\sup_{\mathrm{enc},\mathrm{dec}}\ \mathbb P(\mathcal E_\delta).
\]
令 \(\mathcal V_w(\delta)\) 与 \(R_w(\delta)\) 如定义 \ref{def:pom-typeclass-volume-and-threshold}。
当 \(t\to\infty\) 且 \(B=\lfloor rt/\log 2\rfloor\)（码率 \(r\ge 0\) 固定）时，成立相变：
\[
\boxed{\;
\lim_{t\to\infty}\mathrm{Succ}^{\mathrm{opt}}_t(B,\delta)=
\begin{cases}
0, & r<R_w(\delta),\\
1, & r>R_w(\delta).
\end{cases}\;}
\]
并且当 \(r<R_w(\delta)\) 时有指数上界
\[
\boxed{\;
\limsup_{t\to\infty}\frac{1}{t}\log \mathrm{Succ}^{\mathrm{opt}}_t(B,\delta)\ \le\ r-R_w(\delta)\;<0\; }.
\]
\end{theorem}
\begin{proof}
由定型计数，\(|\mathcal T_t(w)|=\exp(tH(w)+o(t))\)。
由定理 \ref{thm:pom-typeclass-hamming-ball-volume-variational}，对任意参考点 \(x^t\in\mathcal T_t(w)\) 有
\(|\mathcal B_t(w,\delta;x^t)|=\exp(t\mathcal V_w(\delta)+o(t))\)，且误差项与参考点无关。
\par
\textbf{上界：}
任意编码至多产生 \(M=2^B=\exp(rt+o(t))\) 个解码输出点，
其 \(\delta\)-成功集合被包含于至多 \(M\) 个定型内球的并集，覆盖数不超过
\(
M\exp(t\mathcal V_w(\delta)+o(t))
\)。
除以 \(|\mathcal T_t(w)|\) 得
\[
\mathrm{Succ}^{\mathrm{opt}}_t(B,\delta)
\le
\exp\!\bigl(t(r-\!R_w(\delta))+o(t)\bigr),
\]
推出当 \(r<R_w(\delta)\) 时成功率指数衰减，并得到所示 \(\limsup\) 上界。
\par
\textbf{下界：}
取一个最大化的球覆盖（贪婪选取球心即可），其所需球数的指数率等于
\(|\mathcal T_t(w)|/|\mathcal B_t(w,\delta;\cdot)|=\exp(tR_w(\delta)+o(t))\)。
当 \(r>R_w(\delta)\) 时，允许的码字数 \(M=\exp(rt+o(t))\) 在指数尺度上严格大于所需覆盖数，
故可构造使覆盖比例趋于 \(1\) 的码本，从而 \(\mathrm{Succ}^{\mathrm{opt}}_t(B,\delta)\to 1\)。
证毕。
\end{proof}

\paragraph{最优耦合的标量化：对角乘子的单参求解}
沿用定理 \ref{thm:pom-typeclass-hamming-ball-unique-maximizer-exponential-family} 的指数族参数。
写 \(\kappa:=e^{\lambda}-1\in(-1,\infty)\)，并记 \(A:=\sum_{x\in X}u_x\)。

\begin{theorem}[边缘约束的二次根与自洽方程]\label{thm:pom-typeclass-diagonal-coupling-scalarization}
对任意给定 \(\kappa>-1\) 与任意 \(A>0\)，边缘约束 \(P_X=P_Y=w\) 等价于对每个 \(x\in X\) 满足二次方程
\[
\boxed{\;
\kappa u_x^2+A u_x-w(x)=0\; } .
\]
从而（取正根）有显式表达
\[
u_x(A,\kappa)=
\begin{cases}
\dfrac{-A+\sqrt{A^2+4\kappa w(x)}}{2\kappa}, & \kappa\ne0,\\[10pt]
\dfrac{w(x)}{A}, & \kappa=0.
\end{cases}
\]
并且对每个固定 \(\kappa>-1\)，存在唯一 \(A=A(\kappa)>0\) 使得自洽条件 \(A=\sum_x u_x(A,\kappa)\) 成立；
该唯一解给出唯一向量 \(u(\kappa)\)。
\leanverified{quadConstraint\_seed\_1\_1\_2}
\leanverified{quadConstraint\_seed\_2\_3\_5}
\leanverified{quadConstraint\_seed\_1\_5\_6}
\leanverified{quadConstraint\_seed\_1\_0\_4}
\leanverified{quadConstraint\_seed\_3\_1\_4}
\leanverified{discriminant\_seed\_1\_1\_2}
\leanverified{discriminant\_seed\_2\_3\_5}
\leanverified{discriminant\_seed\_1\_0\_4}
\leanverified{discriminant\_pos\_of\_pos}
\leanverified{quadConstraint\_linear}
\leanverified{linear\_root\_seed}
\leanverified{selfConsistency\_n2}
\leanverified{selfConsistency\_n3}
\leanverified{selfConsistency\_n5}
\leanverified{paper\_pom\_typeclass\_diagonal\_coupling\_seeds}
\leanverified{paper\_pom\_typeclass\_diagonal\_coupling\_package}
\end{theorem}
\begin{proof}
由指数族形式可写（其中 \(A=\sum_x u_x\)）
\[
P(x,y)=\frac{u_xu_y\bigl(1+\kappa\mathbf 1_{\{x=y\}}\bigr)}{Z},
\qquad
Z:=\sum_{a,b\in X}u_au_b\bigl(1+\kappa\mathbf 1_{\{a=b\}}\bigr)=A^2+\kappa\sum_{a\in X}u_a^2.
\]
对固定 \(x\) 求和得边缘
\[
w(x)=\sum_y P(x,y)=\frac{u_x(A+\kappa u_x)}{Z}.
\]
注意 \(u\mapsto cu\ (c>0)\) 不改变 \(P\)（分子与分母同乘 \(c^2\)），故可选取等价标度使得 \(Z=1\)。
在该规范化下，边缘约束等价于 \(w(x)=u_x(A+\kappa u_x)\)，即所示二次方程。
正根显式解立得。
\par
对固定 \(\kappa>-1\)，函数 \(A\mapsto \sum_x u_x(A,\kappa)\) 连续严格递减，并在 \(A\downarrow 0\) 时趋于 \(+\infty\)，在 \(A\to\infty\) 时趋于 \(0\)，
故方程 \(A=\sum_x u_x(A,\kappa)\) 存在唯一解。
唯一性给出 \(u(\kappa)\) 的唯一性。
最后，由 \(w(x)=u_x(A+\kappa u_x)\) 对 \(x\) 求和得 \(A^2+\kappa\sum_x u_x^2=\sum_x w(x)=1\)，从而 \(Z=1\) 与上述规范化一致。证毕。
\end{proof}

\begin{theorem}[对角质量的严格单调性与唯一匹配]\label{thm:pom-typeclass-diagonal-mass-monotone}
令 \(P^\star_{\delta(\kappa)}\) 为由参数 \(\kappa\in(-1,\infty)\) 给出的最优耦合，并记其对角质量
\[
\Delta(\kappa):=\sum_{x\in X}P^\star_{\delta(\kappa)}(x,x)=\mathbb P(X=Y).
\]
则 \(\Delta(\kappa)\) 在 \(\kappa\in(-1,\infty)\) 上严格递增，且满足端点值
\[
\Delta(0)=\sum_x w(x)^2,\qquad
\lim_{\kappa\downarrow-1}\Delta(\kappa)=0,\qquad
\lim_{\kappa\to\infty}\Delta(\kappa)=1.
\]
因此当 \(1-\delta>\sum_x w(x)^2\) 时，存在唯一 \(\kappa(\delta)>0\) 使 \(\Delta(\kappa(\delta))=1-\delta\)；
而当 \(1-\delta\le \sum_x w(x)^2\) 时，最优耦合为独立耦合并对应 \(\kappa(\delta)=0\)。
\leanverified{paper\_pom\_typeclass\_diagonal\_mass\_monotone}
\end{theorem}
\begin{proof}
由指数族形式与边缘约束，可将 \(\Delta(\kappa)\) 表示为 \(u(\kappa)\) 的有理函数；
当 \(\kappa\) 增大时，对角权重 \((1+\kappa)u_x^2\) 相对非对角权重 \(u_xu_y\ (x\neq y)\) 严格上升，
而边缘被固定为 \(w\)，故对角质量必须严格增大。
端点值由 \(\kappa\to 0\) 的连续性与 \(\kappa\downarrow-1,\ \kappa\to\infty\) 的极限（对角项分别被消去/强制）直接得到。
严格单调与端点满射保证：对任意目标对角质量 \(p\in(0,1)\) 存在唯一 \(\kappa\in(-1,\infty)\) 使 \(\Delta(\kappa)=p\)。
结合定理 \ref{thm:pom-typeclass-hamming-ball-unique-maximizer-exponential-family} 中“活跃约束 \(\Leftrightarrow 1-\delta>\sum_x w(x)^2\)”的讨论，即得所示匹配关系。证毕。
\end{proof}

\begin{theorem}[保持--重采样结构：最优耦合对应一参 Gibbs 通道]\label{thm:pom-typeclass-keep-resample-channel}
令 \(P^\star_\delta\) 为最优耦合，并用参数 \((u,\kappa)\) 表示。
记
\[
A=\sum_x u_x,\qquad \pi(x):=\frac{u_x}{A},
\qquad
p_x:=\mathbb P_{P^\star_\delta}(Y=x\mid X=x)=\frac{(1+\kappa)u_x}{A+\kappa u_x}\in(0,1).
\]
则条件分布 \(Y\mid X=x\) 具有分解
\[
\boxed{\;
\mathbb P(Y=x\mid X=x)=p_x,\qquad
\mathbb P(Y=y\mid X=x,\ y\neq x)=(1-p_x)\frac{\pi(y)}{1-\pi(x)}\; }.
\]
\leanverified{paper\_pom\_typeclass\_keep\_resample\_channel}
\end{theorem}
\begin{proof}
由指数族形式，固定 \(x\) 时有
\(
P^\star_\delta(x,x)\propto (1+\kappa)u_x^2
\)
而对 \(y\neq x\) 有
\(
P^\star_\delta(x,y)\propto u_xu_y
\)。
对 \(y\) 归一化并写成 \(\pi\) 的形式即可得到分解。证毕。
\end{proof}

\begin{theorem}[覆盖阈值的互信息刻画：对角约束下的最小互信息]\label{thm:pom-typeclass-rate-distortion-mutual-information}
令
\[
\mathcal C_\delta(w):=\Bigl\{P\in\Delta(X\times X):\ P_X=P_Y=w,\ \mathbb P_P(X=Y)\ge 1-\delta\Bigr\}.
\]
则有严格等式
\[
\boxed{\;
R_w(\delta)=\inf_{P\in\mathcal C_\delta(w)} I_P(X;Y)\; } ,
\]
并且该下确界由唯一耦合 \(P^\star_\delta\) 取到。
\leanverified{paper\_pom\_typeclass\_rate\_distortion\_mutual\_information}
\end{theorem}
\begin{proof}
对任意 \(P\in\mathcal C_\delta(w)\)，有恒等式
\(
I_P(X;Y)=H(w)-H_P(Y|X)
\)。
因此在 \(\mathcal C_\delta(w)\) 上最小化互信息等价于最大化条件熵。
由定理 \ref{thm:pom-typeclass-hamming-ball-volume-variational}，
最大条件熵的值为 \(\mathcal V_w(\delta)\)，且由定理 \ref{thm:pom-typeclass-hamming-ball-unique-maximizer-exponential-family} 的唯一极大耦合达到。
故
\(
R_w(\delta)=H(w)-\mathcal V_w(\delta)=I_{P^\star_\delta}(X;Y)
\)。
证毕。
\end{proof}

\begin{theorem}[对偶导数恒等式：\texorpdfstring{$R_w'(\delta)$}{Rw'(delta)} 等于最优对角乘子的相反数]\label{thm:pom-typeclass-dual-derivative}
记 \(p_0=\sum_{x\in X}w(x)^2\)。
当 \(1-\delta>p_0\)（对角质量约束活跃）时，函数 \(\delta\mapsto R_w(\delta)\) 在该区间上可微，且满足
\[
\boxed{\;
-\frac{d}{d\delta}R_w(\delta)=\lambda(\delta)\; } ,
\]
其中 \(\lambda(\delta)\) 为定理 \ref{thm:pom-typeclass-hamming-ball-unique-maximizer-exponential-family} 中最优耦合的 KKT 乘子（对角质量约束的影子价格）。
当 \(1-\delta\le p_0\) 时，最优耦合为独立耦合，\(R_w(\delta)=0\) 且 \(\lambda(\delta)=0\)。
\end{theorem}
\begin{proof}
将对角质量约束写为
\(
g_\delta(P):=(1-\delta)-\mathbb P_P(X=Y)\le 0
\)。
对最大化问题
\(
\sup\{H_P(Y|X):\ P_X=P_Y=w,\ g_\delta(P)\le 0\}
\)，
取拉格朗日函数
\[
\mathscr L(P,\lambda)
:=H_P(Y|X)-\lambda\,g_\delta(P)
=H_P(Y|X)+\lambda\bigl(\mathbb P_P(X=Y)-1+\delta\bigr),
\qquad \lambda\ge 0.
\]
由定理 \ref{thm:pom-typeclass-hamming-ball-unique-maximizer-exponential-family}：
当 \(1-\delta>p_0\) 时约束活跃并饱和（\(\lambda(\delta)>0\)）；
当 \(1-\delta\le p_0\) 时最优解为独立耦合且 \(\lambda(\delta)=0\)。
两种情形下由严格凹性得到强对偶与乘子唯一性。
因此最优值可写为
\[
\mathcal V_w(\delta)=\max_{P_X=P_Y=w}\ \mathscr L(P,\lambda(\delta)),
\]
且 \(\lambda(\delta)\) 与指数族形式中的对角系数一致。
由包络定理，
\[
\frac{d}{d\delta}\mathcal V_w(\delta)=\frac{\partial}{\partial\delta}\mathscr L(P^\star_\delta,\lambda(\delta))=\lambda(\delta).
\]
而 \(R_w(\delta)=H(w)-\mathcal V_w(\delta)\)，故 \(R_w'(\delta)=-\lambda(\delta)\)，即所示恒等式。证毕。
\end{proof}

\endinput


% (User input window)

\paragraph{两温度失败指数的变分结构：幂律耦合、主序链与临界二次律}
本段把定理 \ref{thm:pom-microcanonical-hypergeometric-ldp-kl-sym} 的 KL--对称速率函数
从 \((q\mapsto J_\beta(q))\) 的单变量表述提升为一个两温度变分问题，并在熵阈值的超临界侧给出唯一极小化器的结构。
其结果可作为“剩余组成熵偏大”事件的指数级失败证书，并与定理 \ref{thm:pom-microcanonical-query-plus-bits-phase} 的临界窗叙述兼容。

\begin{definition}[两温度失败指数：熵阈值的最小 KL--对称代价]\label{def:pom-microcanonical-two-temperature-failure-exponent}
设 \(X\) 为有限集（\(|X|=n\ge 2\)），取全支撑分布 \(w\in\Delta(X)\)，并固定 \(\beta\in(0,1)\)。
对任意 \(q,r\in\Delta(X)\) 满足质量守恒
\[
w=\beta q+(1-\beta)r,
\]
定义两温度散度泛函
\[
J_\beta(q,r):=\beta D(q\|w)+(1-\beta)D(r\|w).
\]
对熵阈值 \(h>H(w)\)，定义超临界失败指数
\[
E_{\mathrm{fail}}(\beta,h)
:=
\inf\Bigl\{J_\beta(q,r):\ q,r\in\Delta(X),\ w=\beta q+(1-\beta)r,\ H(r)\ge h\Bigr\}.
\]
若可行域为空，则约定 \(E_{\mathrm{fail}}(\beta,h)=+\infty\)。
\end{definition}

\begin{theorem}[两温度 KKT 结构：唯一极小化器的幂律耦合]\label{thm:pom-microcanonical-two-temperature-kkt-power-law}
\leanverified{paper\_pom\_microcanonical\_two\_temperature\_kkt\_power\_law}
沿用定义 \ref{def:pom-microcanonical-two-temperature-failure-exponent}，并假设给定 \(h\) 使可行域非空且 \(h>\!H(w)\)。
则极小化问题存在唯一极小点 \((q^\star,r^\star)\)，并且熵约束必饱和：
\[
H(r^\star)=h.
\]
此外，存在唯一乘子 \(\eta^\star>0\)（因而存在唯一指数 \(\gamma>1\)）使得对所有 \(x\in X\),
\[
q^\star(x)=\frac{(r^\star(x))^{\gamma}}{\sum_{y\in X}(r^\star(y))^{\gamma}},
\qquad
\gamma=1+\frac{\eta^\star}{1-\beta},
\]
等价地
\(
\log q^\star=\gamma\log r^\star+\mathrm{const}
\)
（常数在 \(x\) 上不依赖）。
\end{theorem}
\begin{proof}
可行域由线性等式与凹函数 \(H\) 的上水平集给出，故为紧凸集；而 \(J_\beta\) 作为两项 KL 散度的正权和，在 \(\Delta(X)^2\) 的内点上严格凸，从而极小点唯一。
\par
由于 \(w\) 全支撑且 \(\beta\in(0,1)\)，若 \((q,r)\) 可行，则 \(q,r\) 自动全支撑；因此可在内点上写 KKT 系统。
若熵约束不饱和，则对应乘子 \(\eta^\star=0\)，从而极小点必须同时极小化仅含质量守恒的松弛问题；而该松弛问题的唯一极小点为 \(q=r=w\)，与 \(h>H(w)\) 矛盾。
故 \(\eta^\star>0\) 且 \(H(r^\star)=h\)。
\par
对约束写拉格朗日函数并对每个坐标取驻点条件，消去质量守恒的乘子后可得
\(
\log r^\star=\alpha\log q^\star+\mathrm{const}
\),
其中
\(
\alpha=\frac{1-\beta}{1-\beta+\eta^\star}\in(0,1)
\)。
从而 \(q^\star\propto (r^\star)^{1/\alpha}\)，并令 \(\gamma:=1/\alpha=1+\eta^\star/(1-\beta)\) 得到所示幂律耦合。
唯一性迫使 \(\eta^\star\) 与 \(\gamma\) 唯一。证毕。
\end{proof}

\begin{proposition}[幂律轨道的存在唯一性与熵的严格单调性]\label{prop:pom-microcanonical-two-temperature-escort-orbit-entropy-monotone}
沿用定义 \ref{def:pom-microcanonical-two-temperature-failure-exponent}。
对每个 \(\gamma>1\)，定义幂 escort 变换
\[
\mathrm{Escort}_\gamma(r)(x):=\frac{r(x)^\gamma}{\sum_{y\in X}r(y)^\gamma}\qquad(r\in\Delta(X)).
\]
若存在 \((q,r)\in\Delta(X)^2\) 满足 \(w=\beta q+(1-\beta)r\) 且 \(q=\mathrm{Escort}_\gamma(r)\)，则该解唯一；记之为 \((q_\gamma,r_\gamma)\)。
并且在其定义域上，映射
\[
\gamma\longmapsto H(r_\gamma)
\]
连续且严格递增，并满足
\[
\lim_{\gamma\downarrow 1}r_\gamma=w.
\]
令
\[
H_{\max}(\beta,w):=\sup\Bigl\{H(r):\ r\in\Delta(X),\ \frac{w-(1-\beta)r}{\beta}\in\Delta(X)\Bigr\}\in(H(w),\log n],
\]
则 \(\lim_{\gamma\to\infty}H(r_\gamma)=H_{\max}(\beta,w)\)。
特别地，若均匀分布 \(U\) 可行（等价于 \(w(x)\ge (1-\beta)/n\) 对所有 \(x\) 成立），则 \(H_{\max}(\beta,w)=\log n\) 且 \(r_\gamma\to U\)。
因此对任意 \(h\in(H(w),H_{\max}(\beta,w))\) 存在唯一 \(\gamma(h)>1\) 使 \(H(r_{\gamma(h)})=h\)，并且该 \((q_{\gamma(h)},r_{\gamma(h)})\) 与定理 \ref{thm:pom-microcanonical-two-temperature-kkt-power-law} 的极小化器一致。
\end{proposition}
\begin{proof}
对固定 \(\eta>0\) 考虑惩罚型问题
\[
\mathcal{V}(\eta):=\inf\Bigl\{J_\beta(q,r)-\eta H(r):\ q,r\in\Delta(X),\ w=\beta q+(1-\beta)r\Bigr\}.
\]
由于 \(-H\) 为严格凸函数且 \(J_\beta\) 严格凸，故该问题存在唯一极小点 \((q_\eta,r_\eta)\)；
其 KKT 系统给出幂律关系 \(q_\eta=\mathrm{Escort}_\gamma(r_\eta)\) 且 \(\gamma=1+\eta/(1-\beta)>1\)。
因此对每个 \(\gamma>1\) 至多对应一个可行解，记为 \((q_\gamma,r_\gamma)\)。
\par
函数 \(\eta\mapsto \mathcal{V}(\eta)\) 为 \(\eta\) 的下确界包络，因而为凹函数；并且在唯一极小点下可微且满足
\(
\mathcal{V}'(\eta)=-H(r_\eta)
\)。
凹性推出 \(-\mathcal{V}'(\eta)=H(r_\eta)\) 随 \(\eta\)（从而随 \(\gamma\)）严格递增，连续性由隐函数/稳定性与唯一极小性给出。
\par
当 \(\eta\downarrow 0\) 时，惩罚项消失且松弛问题的唯一极小点为 \(q=r=w\)，故 \(r_\gamma\to w\)（\(\gamma\downarrow 1\)）。
当 \(\eta\to\infty\) 时，\(-\eta H(r)\) 主导并迫使 \(r_\eta\) 收敛到“在可行多面体内最大化 \(H\)”的唯一解，从而 \(H(r_\gamma)\uparrow H_{\max}(\beta,w)\)；
若 \(U\) 可行，则最大熵点即为 \(U\) 且极限为 \(\log n\)。
最后，对 \(h\in(H(w),H_{\max})\) 由严格单调性得到唯一 \(\gamma(h)\)，并与定理 \ref{thm:pom-microcanonical-two-temperature-kkt-power-law} 的唯一极小化器匹配。证毕。
\end{proof}

\begin{corollary}[共同排序与主序链：\texorpdfstring{$r^\star\prec w\prec q^\star$}{r* ≺ w ≺ q*}]\label{cor:pom-microcanonical-two-temperature-majorization-chain}
在定理 \ref{thm:pom-microcanonical-two-temperature-kkt-power-law} 的设定下，存在同一排列使得 \(r^\star,w,q^\star\) 的坐标按非增序排列后满足逐段和链
\[
\sum_{i=1}^k r^\star_{(i)}\ \le\ \sum_{i=1}^k w_{(i)}\ \le\ \sum_{i=1}^k q^\star_{(i)}\qquad(k=1,\dots,n-1),
\]
且至少对某个 \(k\) 严格不等；从而 \(r^\star\prec w\prec q^\star\)。
\end{corollary}
\begin{proof}
由幂律耦合 \(q^\star=\mathrm{Escort}_\gamma(r^\star)\)（\(\gamma>1\)），两者具有相同坐标排序，且幂变换在主序意义上将分布严格冷化：\(r^\star\prec q^\star\)（Hardy--Littlewood--Pólya）。
又因 \(w=\beta q^\star+(1-\beta)r^\star\) 为同序向量的凸组合，逐段和对凸组合线性，故得到链式不等式并推出 \(r^\star\prec w\prec q^\star\)。
严格性由 \(h>H(w)\) 排除 \(r^\star=w\)（从而亦排除 \(q^\star=w\)）得到。证毕。
\end{proof}

\begin{proposition}[Jensen--Shannon 结构恒等式：\texorpdfstring{$J_\beta$}{Jbeta} 等于混合熵亏损]\label{prop:pom-microcanonical-two-temperature-js-identity}
\leanverified{paper\_pom\_microcanonical\_two\_temperature\_js\_identity}
在约束 \(w=\beta q+(1-\beta)r\) 下，对任意可行 \((q,r)\) 有恒等式
\[
J_\beta(q,r)=H(w)-\beta H(q)-(1-\beta)H(r).
\]
因此当可行域非空且 \(h>H(w)\) 时，
\[
E_{\mathrm{fail}}(\beta,h)
=H(w)-(1-\beta)h-\beta\sup\Bigl\{H(q):\exists r,\ w=\beta q+(1-\beta)r,\ H(r)=h\Bigr\}.
\]
\end{proposition}
\begin{proof}
由 \(D(p\|w)=-H(p)-\sum_x p(x)\log w(x)\)，并用质量守恒
\(
\beta q+(1-\beta)r=w
\)
消去 \(\sum_x(\beta q+(1-\beta)r)\log w\)，即得第一式。
第二式由定理 \ref{thm:pom-microcanonical-two-temperature-kkt-power-law} 的 \(H(r^\star)=h\) 代入并对 \(q\) 作等价极大化得到。证毕。
\end{proof}

\begin{theorem}[强对偶、唯一乘子与敏感性公式]\label{thm:pom-microcanonical-two-temperature-strong-duality-sensitivity}
\leanverified{paper\_pom\_microcanonical\_two\_temperature\_strong\_duality\_sensitivity}
沿用定义 \ref{def:pom-microcanonical-two-temperature-failure-exponent}，假设给定 \(h\in(H(w),H_{\max}(\beta,w))\) 使可行域非空。
定义对偶值函数（\(\eta\ge 0\)）
\[
\mathcal{V}(\eta):=\inf\Bigl\{J_\beta(q,r)-\eta H(r):\ q,r\in\Delta(X),\ w=\beta q+(1-\beta)r\Bigr\}.
\]
则成立强对偶
\[
E_{\mathrm{fail}}(\beta,h)=\sup_{\eta\ge 0}\bigl(\mathcal{V}(\eta)+\eta h\bigr),
\]
且上确界在唯一 \(\eta^\star>0\) 处达到。
并且函数 \(h\mapsto E_{\mathrm{fail}}(\beta,h)\) 在 \((H(w),H_{\max}(\beta,w))\) 上可微，且
\[
\frac{\partial}{\partial h}E_{\mathrm{fail}}(\beta,h)=\eta^\star,
\qquad
\gamma(h)=1+\frac{\eta^\star}{1-\beta},
\]
其中 \(\gamma(h)\) 为定理 \ref{thm:pom-microcanonical-two-temperature-kkt-power-law} 的唯一幂指数。
\end{theorem}
\begin{proof}
约束 \(H(r)\ge h\) 等价于凸不等式 \(h-H(r)\le 0\)，并且在 \(h< H_{\max}(\beta,w)\) 时满足 Slater 条件（存在严格可行点）。
因此强对偶成立，且由于原问题极小点唯一，KKT 乘子亦唯一且满足 \(\eta^\star>0\)（同定理 \ref{thm:pom-microcanonical-two-temperature-kkt-power-law}）。
可微性与导数公式由凸对偶的一般敏感性定理（包络定理）推出；
最后一式由幂律关系 \(\gamma=1+\eta/(1-\beta)\) 代入得到。证毕。
\end{proof}

\begin{theorem}[临界熵窗的二次律：\texorpdfstring{$h\downarrow H(w)$}{h↓H(w)} 时的精确起跳系数]\label{thm:pom-microcanonical-two-temperature-critical-quadratic-law}
设 \(w\) 全支撑且非均匀，并记
\[
V(w):=\Var_{X\sim w}\bigl(\log w(X)\bigr)>0
\]
（同推论 \ref{cor:pom-microcanonical-information-clt}）。
令 \(\delta:=h-H(w)\downarrow 0^+\) 且取 \(h\in(H(w),H_{\max}(\beta,w))\)。
则有渐近展开
\[
E_{\mathrm{fail}}(\beta,h)
=\frac{1-\beta}{2\beta}\cdot \frac{\delta^2}{V(w)}+o(\delta^2).
\]
若改用比特密度
\(
b:=\frac{(1-\beta)h}{\log 2}
\)
与临界点
\(
b_{\mathrm{crit}}:=\frac{(1-\beta)H(w)}{\log 2}
\),
则等价地
\[
E_{\mathrm{fail}}(\beta,b)
=\frac{(b-b_{\mathrm{crit}})^2(\log 2)^2}{2\beta(1-\beta)V(w)}+o\!\left((b-b_{\mathrm{crit}})^2\right).
\]
\end{theorem}
\begin{proof}
当 \(\delta\downarrow 0\) 时，最优点 \((q^\star,r^\star)\) 逼近 \((w,w)\)。
写 \(r=w+u\)（\(\sum_x u(x)=0\)），则由质量守恒 \(q=w-\frac{1-\beta}{\beta}u\)。
对 KL 在 \(w\) 处作二阶展开可得
\[
J_\beta(q,r)
=\frac{1-\beta}{2\beta}\sum_{x\in X}\frac{u(x)^2}{w(x)}+o(\|u\|_2^2).
\]
另一方面，熵的一阶变分（利用 \(\sum_x u(x)=0\)）给出
\[
H(w+u)-H(w)=-\sum_{x\in X}u(x)\log w(x)+o(\|u\|_2),
\]
故主阶约束为 \(\sum_x u(x)\log w(x)=-\delta\)。
因此主阶问题化为：在二次型 \(\sum u^2/w\) 下最小化线性约束 \(\sum u\log w=-\delta\)。
拉格朗日乘子法给出唯一主方向
\[
u^\star(x)=-\frac{\delta}{V(w)}\,w(x)\Bigl(\log w(x)-\EE_w[\log w]\Bigr),
\]
并满足
\(
\sum_x\frac{(u^\star(x))^2}{w(x)}=\delta^2/V(w)
\)。
代回即得
\(
E_{\mathrm{fail}}(\beta,h)=\frac{1-\beta}{2\beta}\cdot \frac{\delta^2}{V(w)}+o(\delta^2)
\)。
比特密度形式由 \(\delta=\frac{\log 2}{1-\beta}(b-b_{\mathrm{crit}})\) 代换得到。证毕。
\end{proof}


% (User input window)
% (User input window)

\paragraph{查询—失真任务：剩余类型上的均匀定型源、强逆平面与零比特碰撞阈值}
在微正则先验 \(F\sim\mathrm{Unif}(\mathfrak{Fold}_m(d))\) 下，\(k\) 次互异查询只揭示计数信息；
未查询部分在条件后验中退化为一个“剩余类型”约束下的均匀定型源（定理 \ref{thm:pom-microcanonical-count-sufficient-statistic-posterior-uniform}）。
本段将该结构与允许相对 Hamming 失真的恢复任务对齐，得到一个在指数尺度上呈严格平面的强逆相图，
并给出超临界失败指数的收缩变分式与临界窗的高斯转折。

\begin{definition}[未查询标签向量与相对 Hamming 失真]\label{def:pom-microcanonical-query-distortion-setup}
在定义 \ref{def:pom-microcanonical-adaptive-trajectory-counts} 的设定下，
令 \(Q_k=\{\omega_1,\dots,\omega_k\}\subset\Omega_m\)、\(U_k:=\Omega_m\setminus Q_k\)，并记
\[
t:=\card{U_k}=N-k.
\]
固定 \(U_k\) 的任意排序 \(U_k=\{\upsilon_1,\dots,\upsilon_t\}\)，并定义未查询标签向量
\[
Y^t:=(F(\upsilon_1),\dots,F(\upsilon_t))\in X_m^t.
\]
对任意两向量 \(y^t,\widehat y^t\in X_m^t\) 定义 Hamming 距离与相对失真
\[
d_{\mathrm H}(y^t,\widehat y^t):=\card{\{i\in\{1,\dots,t\}:\ y_i\neq \widehat y_i\}},
\qquad
\frac{1}{t}d_{\mathrm H}(y^t,\widehat y^t)\in[0,1].
\]
\end{definition}

\begin{definition}[查询 + 外置信息下的失真成功率]\label{def:pom-microcanonical-query-distortion-success}
给定 \(k\) 次查询轨迹 \(T_{1:k}\) 后，允许额外获得 \(B\) 比特外置信息：
编码器为任意映射 \(\mathrm{enc}:\mathfrak{Fold}_m(d)\to\{0,1\}^B\)，解码器为任意映射
\[
\mathrm{dec}:X_m^k\times\{0,1\}^B\to X_m^t,
\qquad
\widehat Y^t:=\mathrm{dec}(T_{1:k},\mathrm{enc}(F)).
\]
对 \(\delta\in[0,1]\) 定义最优（Bayes）失真成功率
\[
\mathrm{Succ}^{\mathrm{opt}}_{k,B}(\delta)
:=\sup_{\mathrm{enc},\mathrm{dec}}\ \mathbb{P}\!\left(\frac{1}{t}d_{\mathrm H}(Y^t,\widehat Y^t)\le \delta\right).
\]
\end{definition}

\begin{definition}[剩余计数、剩余类型与定型类]\label{def:pom-microcanonical-residual-type-class}
对轨迹 \(T_{1:k}\) 的计数向量 \(C_k\)（定义 \ref{def:pom-microcanonical-adaptive-trajectory-counts}），定义剩余计数向量
\[
R_x:=d(x)-C_k(x),\qquad x\in X_m,
\qquad
\sum_{x\in X_m}R_x=t,
\]
以及剩余类型 \(r\in\Delta(X_m)\)：
\[
r(x):=\frac{R_x}{t}.
\]
定义相应的定型类（类型类）
\[
\mathcal{T}_t(r):=\left\{y^t\in X_m^t:\ \card{\{i:\ y_i=x\}}=R_x\ \ \forall x\in X_m\right\}.
\]
\end{definition}

\begin{proposition}[策略不变的后验结构：未查询部分为均匀定型源]\label{prop:pom-microcanonical-query-distortion-strategy-invariant-posterior}
沿用定义 \ref{def:pom-microcanonical-residual-type-class}。
对任意确定性自适应互异查询策略与任意轨迹 \(T_{1:k}\)，有严格等式
\[
\mathrm{Law}\bigl(Y^t\mid T_{1:k}\bigr)=\mathrm{Unif}\bigl(\mathcal{T}_t(r)\bigr).
\]
特别地，条件后验仅依赖于剩余计数向量 \(R\)（等价于剩余类型 \(r\)），与查询点的自适应选择无关。
\end{proposition}
\begin{proof}
由定理 \ref{thm:pom-microcanonical-count-sufficient-statistic-posterior-uniform}，
\(
\mathrm{Law}(F|_{U_k}\mid T_{1:k})=\mathrm{Unif}(\mathfrak{F}_{U_k}(d-C_k))
\)；
固定 \(U_k\) 的排序后，映射 \(F|_{U_k}\) 与向量 \(Y^t\) 一一对应，且约束 \(\card{(F|_{U_k})^{-1}(x)}=R_x\) 等价于 \(Y^t\in\mathcal{T}_t(r)\)。
故条件分布为 \(\mathrm{Unif}(\mathcal{T}_t(r))\)。
查询策略不变性由上述后验仅依赖 \(C_k\)（从而仅依赖 \(R=d-C_k\)）直接推出。证毕。
\end{proof}

\begin{proposition}[查询—失真任务的充分统计]\label{prop:pom-microcanonical-query-distortion-sufficient-statistic}
在定义 \ref{def:pom-microcanonical-query-distortion-success} 的设定下，
对任意查询策略与任意轨迹 \(T_{1:k}\)，条件最优成功率
\[
\sup_{\mathrm{enc},\mathrm{dec}}\ \mathbb{P}\!\left(\frac{1}{t}d_{\mathrm H}(Y^t,\widehat Y^t)\le \delta\ \Big|\ T_{1:k}\right)
\]
仅为 \((t,r,B,\delta)\) 的函数。
因此在 \(N\to\infty\) 的指数尺度上，\(\mathrm{Succ}^{\mathrm{opt}}_{k,B}(\delta)\) 的相图完全由 \((\beta,w,\delta)\) 与比特密度 \(b=B/N\) 决定。
\end{proposition}
\begin{proof}
由命题 \ref{prop:pom-microcanonical-query-distortion-strategy-invariant-posterior}，给定轨迹后 \(Y^t\) 为 \(\mathrm{Unif}(\mathcal{T}_t(r))\)，且该条件分布不再含 \(\Omega_m\) 的几何信息。
于是任何编码/解码对在条件下的成功率仅由该均匀定型源、码长 \(B\) 与失真阈值 \(\delta\) 决定，即为 \((t,r,B,\delta)\) 的函数。
最后一句由 \(k=\lfloor\beta N\rfloor\)（从而 \(t/N\to 1-\beta\)）以及 \(r\to w\) 的无放回抽样大数律（定理 \ref{thm:pom-microcanonical-information-linear-law} 的证明中使用的同型收敛）推出。证毕。
\end{proof}

\begin{definition}[定型—对角约束速率函数]\label{def:pom-typed-diagonal-rate-function}
设 \(X\) 为有限集、\(p\in\Delta(X)\)、\(\delta\in[0,1]\)。
定义速率函数（自然对数单位）
\[
R_p(\delta)
:=\inf\Bigl\{I_P(X;Y):\ P_X=P_Y=p,\ \mathbb{P}(X=Y)\ge 1-\delta\Bigr\},
\]
其中下确界取遍 \(X\times X\) 上的所有联合分布 \(P\)，并以 \(I_P(X;Y)\) 表示其互信息（自然对数）。
\end{definition}

\begin{theorem}[查询比例—比特密度—失真的强逆平面]\label{thm:pom-microcanonical-query-distortion-strong-converse-plane}
\leanverified{paper\_pom\_microcanonical\_query\_distortion\_strong\_converse\_plane}
沿用定义 \ref{def:pom-microcanonical-query-distortion-success}--\ref{def:pom-typed-diagonal-rate-function}。
固定 \(\delta\in[0,1]\)，取 \(k=\lfloor\beta N\rfloor\) 与 \(B=\lfloor bN\rfloor\)（\(\beta\in[0,1)\)、\(b\ge 0\)），并假设 \(n\log N=o(N)\) 且 \(w\) 固定。
则极限存在并满足
\[
\boxed{\;
\lim_{N\to\infty}\frac{1}{N}\log \mathrm{Succ}^{\mathrm{opt}}_{k,B}(\delta)
=\min\Bigl\{0,\ b\log 2-(1-\beta)\,R_w(\delta)\Bigr\}\; }.
\]
因此临界平面可写为
\[
\boxed{\;
b_{\mathrm{crit}}(\beta,\delta)=\frac{(1-\beta)\,R_w(\delta)}{\log 2}\; } .
\]
当 \(\delta=0\) 时有 \(R_w(0)=H(w)\)，从而 \(b_{\mathrm{crit}}(\beta,0)=(1-\beta)H(w)/\log 2\)，与定理 \ref{thm:pom-microcanonical-query-plus-bits-phase} 的精确恢复阈值一致。
\end{theorem}
\begin{proof}
由命题 \ref{prop:pom-microcanonical-query-distortion-strategy-invariant-posterior}，条件于轨迹后源为 \(\mathrm{Unif}(\mathcal{T}_t(r))\)。
\par
\textbf{上界（覆盖体积）。}
定型类规模满足 \(|\mathcal{T}_t(r)|=\exp(tH(r)+o(t))\)（Stirling 展开），其中 \(H(\cdot)\) 为 Shannon 熵（自然对数）。
对任意重构中心 \(\widehat y^t\)，
\(\mathcal{T}_t(r)\cap\{y^t:\ d_{\mathrm H}(y^t,\widehat y^t)\le \delta t\}\)
可按联合类型分解为有限个条件类型类，其指数体积由 \(\sup H(X|Y)\) 控制；
在边缘约束 \(P_X=P_Y=r\) 与对角约束 \(\mathbb{P}(X=Y)\ge 1-\delta\) 下，
\[
\sup H(X|Y)=H(r)-R_r(\delta)
\]
（由 \(I(X;Y)=H(X)-H(X|Y)\) 的等价改写）。
因此存在函数 \(o(t)\)（对固定 \(r,\delta\)）使对所有 \(\widehat y^t\) 有
\[
\card{\mathcal{T}_t(r)\cap\{y^t:\ d_{\mathrm H}(y^t,\widehat y^t)\le \delta t\}}
\le \exp\bigl(t(H(r)-R_r(\delta))+o(t)\bigr).
\]
任意 \(B\) 比特码至多包含 \(2^B\) 个重构中心，故条件成功率至多为覆盖比例：
\[
\mathrm{Succ}^{\mathrm{opt}}_{k,B}(\delta\mid T_{1:k})
\le
\min\!\left(1,\frac{2^B\exp(t(H(r)-R_r(\delta))+o(t))}{\exp(tH(r)+o(t))}\right)
=\min\!\left(1,\exp\bigl(B\log 2-tR_r(\delta)+o(t)\bigr)\right).
\]
对轨迹取期望，并用无放回抽样下 \(r\to w\) 的大数律与 \(R_p(\delta)\) 的连续性（在单纯形内点上）得到指数上界
\(\limsup \frac1N\log \mathrm{Succ}^{\mathrm{opt}}_{k,B}(\delta)\le \min\{0,b\log 2-(1-\beta)R_w(\delta)\}\)。
\par
\textbf{下界（定型覆盖构造）。}
对任意 \(\varepsilon>0\)，在给定 \(r\) 时可选取一组大小 \(\exp(t(R_r(\delta)+\varepsilon))\) 的重构中心使其 \(\delta\)-球覆盖 \(\mathcal{T}_t(r)\) 的 \(1-o(1)\) 比例（定型覆盖的标准构造；以联合类型达到 \(\sup H(X|Y)\) 的分量为主）。
当 \(B\log 2\ge t(R_r(\delta)+\varepsilon)\) 时可使条件成功率趋于 \(1\)；当 \(B\log 2\le t(R_r(\delta)-\varepsilon)\) 时可使条件成功率达到
\(\exp(B\log 2-tR_r(\delta)-o(t))\) 的指数率。
对随机 \(r\) 取期望，并用 \(r\to w\) 得到匹配的指数下界。
令 \(\varepsilon\downarrow 0\) 即得结论。证毕。
\end{proof}

\begin{theorem}[超临界失败指数：由剩余类型偏离触发的收缩变分式]\label{thm:pom-microcanonical-query-distortion-failure-exponent}
在定理 \ref{thm:pom-microcanonical-query-distortion-strong-converse-plane} 的设定下，
进一步假设处于超临界侧
\[
b\log 2>(1-\beta)\,R_w(\delta).
\]
令 \(\mathrm{Fail}^{\mathrm{opt}}_{k,B}(\delta):=1-\mathrm{Succ}^{\mathrm{opt}}_{k,B}(\delta)\)。
则极限
\[
E_{\mathrm{fail}}(\beta,b,\delta)
:=\lim_{N\to\infty}-\frac{1}{N}\log \mathrm{Fail}^{\mathrm{opt}}_{k,B}(\delta)
\]
存在并满足严格等式
\[
\boxed{\;
E_{\mathrm{fail}}(\beta,b,\delta)
=\inf\Bigl\{J_{1-\beta}(r):\ (1-\beta)\,R_r(\delta)\ge b\log 2\Bigr\}\; } ,
\]
其中 \(J_{1-\beta}\) 为剩余类型 \(R_N=r_N\) 的大偏差速率函数（定理 \ref{thm:pom-microcanonical-hypergeometric-ldp-kl-sym} 的对称形式）。
当 \(\delta=0\) 时 \(R_r(0)=H(r)\)，上述公式退化为“剩余熵阈值”的两温度失败指数（定义 \ref{def:pom-microcanonical-two-temperature-failure-exponent}，取 \(h=b\log 2/(1-\beta)\)）。
\end{theorem}
\begin{proof}
由定理 \ref{thm:pom-microcanonical-query-distortion-strong-converse-plane} 的条件形式，
在给定剩余类型 \(r\) 时，若 \(b\log 2>(1-\beta)R_r(\delta)\) 则条件失败概率可压至 \(e^{-\Theta(N)}\)；
若 \(b\log 2<(1-\beta)R_r(\delta)\)，则条件成功率在指数尺度上衰减，从而条件失败趋于 \(1\)。
因此失败事件在指数尺度上等价于
\[
\{(1-\beta)R_{r_N}(\delta)\ge b\log 2\}.
\]
对随机 \(r_N\) 应用定理 \ref{thm:pom-microcanonical-hypergeometric-ldp-kl-sym} 的 LDP（以 \(R_N\) 为收缩像）并用收缩原理得到所示变分式。证毕。
\end{proof}

\begin{theorem}[临界窗的高斯律：\texorpdfstring{$\sqrt N$}{sqrtN} 比特尺度的正态转折]\label{thm:pom-microcanonical-query-distortion-critical-window-clt}
在定理 \ref{thm:pom-microcanonical-query-distortion-strong-converse-plane} 的设定下，
进一步假设映射 \(p\mapsto R_p(\delta)\) 在 \(p=w\)（单纯形内点）处可 Fr\'echet 可微；
取其梯度代表 \(g_\delta\in\RR^{X_m}\) 并中心化
\[
\widetilde g_\delta:=g_\delta-\langle g_\delta,w\rangle\,\mathbf 1.
\]
令
\[
B_N=\left\lfloor \frac{(1-\beta)R_w(\delta)}{\log 2}\,N+\frac{s}{\log 2}\sqrt N\right\rfloor,\qquad s\in\RR.
\]
则当 \(N\to\infty\) 时有极限
\[
\boxed{\;
\mathrm{Succ}^{\mathrm{opt}}_{k,B_N}(\delta)\ \longrightarrow\ \Phi\!\left(\frac{s}{\sigma_{\beta,\delta}}\right)\; } ,
\]
其中 \(\Phi\) 为标准正态分布函数，且
\[
\boxed{\;
\sigma_{\beta,\delta}^2
=\beta(1-\beta)\ \widetilde g_\delta^{\mathsf T}\bigl(\mathrm{diag}(w)-ww^{\mathsf T}\bigr)\widetilde g_\delta\; } .
\]
当 \(\delta=0\) 时 \(R_p(0)=H(p)\)，可取 \(\widetilde g_0(x)=-\log w(x)+\EE_{X\sim w}[\log w(X)]\)，从而 \(\sigma_{\beta,0}^2=\beta(1-\beta)V(w)\)，与定理 \ref{thm:pom-microcanonical-query-plus-bits-phase} 的高斯临界窗一致。
\end{theorem}
\begin{proof}
由定理 \ref{thm:pom-microcanonical-hypergeometric-ldp-kl-sym} 同型的无放回抽样 CLT（定理 \ref{thm:pom-microcanonical-posterior-moduli-clt} 的证明中使用），有
\[
\sqrt N\,(r_N-w)\ \xrightarrow{d}\ \mathcal N\!\left(0,\frac{\beta}{1-\beta}\bigl(\mathrm{diag}(w)-ww^{\mathsf T}\bigr)\right),
\qquad r_N:=\frac{d-C_k}{N-k}.
\]
对 \(R_{r_N}(\delta)\) 施加 delta 方法得
\[
(1-\beta)N\bigl(R_{r_N}(\delta)-R_w(\delta)\bigr)
\ \xrightarrow{d}\ \mathcal N\bigl(0,\sigma_{\beta,\delta}^2\bigr).
\]
在临界标度下，
\[
B_N\log 2-(1-\beta)N R_{r_N}(\delta)=s\sqrt N-(1-\beta)N\bigl(R_{r_N}(\delta)-R_w(\delta)\bigr)+o_p(\sqrt N),
\]
其符号以极限正态变量决定。
由强逆平面（定理 \ref{thm:pom-microcanonical-query-distortion-strong-converse-plane}）的条件形式可将条件成功率在该尺度上等价为阶跃函数，
取期望即得所示高斯 CDF 极限。证毕。
\end{proof}

\begin{proposition}[零速率阈值：由二阶碰撞概率精确刻画]\label{prop:pom-diagonal-rate-zero-threshold-collision}
设 \(X\) 为有限集、\(p\in\Delta(X)\)，并记二阶碰撞概率
\[
p_2(p):=\sum_{x\in X}p(x)^2.
\]
则对任意 \(\delta\in[0,1]\) 有严格充要条件
\[
\boxed{\;
R_p(\delta)=0\quad\Longleftrightarrow\quad \delta\ge 1-p_2(p)\; } .
\]
\end{proposition}
\begin{proof}
若 \(\delta\ge 1-p_2(p)\)，取独立耦合 \(P(x,y)=p(x)p(y)\)，则 \(\mathbb{P}(X=Y)=p_2(p)\ge 1-\delta\) 且 \(I(X;Y)=0\)，从而 \(R_p(\delta)=0\)。
\par
反之，若 \(R_p(\delta)=0\)，则存在可行耦合使 \(I(X;Y)=0\)，而 \(I=0\) 当且仅当 \(X,Y\) 独立，故必有 \(P(x,y)=p(x)p(y)\)；
其对角质量为 \(p_2(p)\)，因此可行性要求 \(p_2(p)\ge 1-\delta\)，即 \(\delta\ge 1-p_2(p)\)。
证毕。
\end{proof}

\begin{theorem}[速率函数的显式参数化与对偶导数]\label{thm:pom-diagonal-rate-explicit-param-and-derivative}
设 \(X\) 为有限集，取全支撑分布 \(w\in\Delta(X)\)，并固定 \(\delta\in(0,1)\) 使 \(R_w(\delta)>0\)（等价于 \(\delta<1-p_2(w)\)，命题 \ref{prop:pom-diagonal-rate-zero-threshold-collision}）。
则实现 \(R_w(\delta)\) 的最优耦合 \(P^\star_\delta\) 唯一存在，并具有如下结构：
存在唯一参数 \(\kappa>0\) 与唯一向量 \(u\in\RR_{>0}^X\) 使
\[
\boxed{\;
P^\star_\delta(x,y)=\frac{u_xu_y\bigl(1+\kappa\,\mathbf 1_{\{x=y\}}\bigr)}{Z}\;},
\qquad
Z=\sum_{a,b\in X}u_au_b\bigl(1+\kappa\,\mathbf 1_{\{a=b\}}\bigr).
\]
并且边缘约束 \(P_X=P_Y=w\) 等价于对每个 \(x\in X\) 的二次自洽关系
\[
\boxed{\;
w(x)=\frac{u_x(A+\kappa u_x)}{Z}\;},
\qquad
A:=\sum_{y\in X}u_y.
\]
在上述参数化下，速率具有闭式表示
\[
\boxed{\;
R_w(\delta)
=2\sum_{x\in X}w(x)\log\frac{u_x}{w(x)}
-\log Z+(1-\delta)\log(1+\kappa)\; } .
\]
令对偶变量 \(\lambda:=\log(1+\kappa)\ge 0\)，则 \(\delta\mapsto R_w(\delta)\) 在 \((0,1-p_2(w))\) 上可微且满足敏感性公式
\[
\boxed{\;
\frac{\mathrm d}{\mathrm d\delta}R_w(\delta)=-\lambda(\delta)=-\log(1+\kappa(\delta))<0\; } .
\]
并且当 \(\delta\uparrow 1-p_2(w)\) 时有 \(\kappa(\delta)\downarrow 0\) 与 \(R_w(\delta)\downarrow 0\)。
\end{theorem}
\begin{proof}
对角约束问题为有限维凸规划：以 \(P\) 为变量、目标 \(I_P(X;Y)\) 凸，约束线性且在 \(\delta<1-p_2(w)\) 时满足 Slater 条件（存在严格可行点）。
故存在唯一最优解 \(P^\star_\delta\)，并且对角约束必饱和：\(\mathbb{P}_{P^\star_\delta}(X=Y)=1-\delta\)。
\par
写拉格朗日函数并对 \(P(x,y)\) 的驻点条件求解，可得指数族形式
\(
P^\star_\delta(x,y)\propto a_x b_y\exp(\lambda\mathbf 1_{\{x=y\}})
\)；
由对称性与边缘同为 \(w\) 可取 \(a=b=u\)，并将 \(\exp(\lambda)-1\) 记为 \(\kappa\)，得到所示的 \((u,\kappa)\) 参数化与归一化常数 \(Z\)。
边缘条件即为对每个 \(x\) 的方程组 \(w(x)=\sum_y P^\star_\delta(x,y)\)，化简即得二次自洽关系。
\par
闭式表达由
\(
\log P^\star_\delta(x,y)=\log u_x+\log u_y+\log(1+\kappa\mathbf 1_{\{x=y\}})-\log Z
\)
展开，并对 \(w\) 的边缘约束分别在 \(x\) 与 \(y\) 侧取期望得到。
\par
对偶导数由包络定理给出：约束可写为 \((1-\delta)-\mathbb{P}(X=Y)\le 0\)，其最优乘子即为 \(\lambda\)，从而 \(\frac{\mathrm d}{\mathrm d\delta}R_w(\delta)=-\lambda(\delta)\)。
当 \(\delta\uparrow 1-p_2(w)\) 时，独立耦合可行且给出 \(R=0\)，唯一性迫使 \(\lambda\downarrow 0\)（即 \(\kappa\downarrow 0\)）并导致 \(R_w(\delta)\downarrow 0\)。证毕。
\end{proof}

% (User input window)

\begin{theorem}[单标量坍缩：最优耦合由“重叠幅度”完全参数化]\label{thm:pom-diagonal-rate-scalar-collapse}
设 \(X\) 为有限集，取全支撑分布 \(w\in\Delta(X)\)，并记
\[
p_2(w):=\sum_{x\in X}w(x)^2,
\qquad
\delta_0:=1-p_2(w).
\]
当 \(\delta\ge \delta_0\) 时有 \(R_w(\delta)=0\)（命题 \ref{prop:pom-diagonal-rate-zero-threshold-collision}）。
下设 \(\delta\in(0,\delta_0)\)，并沿用定理 \ref{thm:pom-diagonal-rate-explicit-param-and-derivative} 的最优耦合指数族表示。
\par
对任意 \(\kappa>0\) 与任意 \(A>0\)，定义
\[
u_{\kappa,A}(x):=\frac{\sqrt{A^2+4\kappa\,w(x)}-A}{2\kappa}\qquad(x\in X),
\qquad
\Phi_\kappa(A):=\sum_{x\in X}u_{\kappa,A}(x).
\]
则成立：
\begin{enumerate}
  \item \textbf{（唯一不动点与规范化缩放）}
  对每个 \(\kappa>0\)，方程 \(A=\Phi_\kappa(A)\) 在 \(A>0\) 上存在唯一解，记为 \(A(\kappa)\in(0,1)\)。
  令 \(u_\kappa(x):=u_{\kappa,A(\kappa)}(x)\)，并定义耦合
  \[
  P_\kappa(x,y):=
  \begin{cases}
  (1+\kappa)\,u_\kappa(x)^2, & x=y,\\
  u_\kappa(x)u_\kappa(y), & x\ne y.
  \end{cases}
  \]
  则 \(P_\kappa\) 为概率分布且满足双边缘约束 \((P_\kappa)_X=(P_\kappa)_Y=w\)。

  \item \textbf{（\(\kappa\)-失真双射）}
  令
  \[
  \delta(\kappa):=1-\sum_{x\in X}P_\kappa(x,x)\in(0,\delta_0).
  \]
  则映射 \(\kappa\mapsto \delta(\kappa)\) 在 \((0,\infty)\) 上连续且严格递减，并把 \((0,\infty)\) 双射到 \((0,\delta_0)\)；
  其端点极限为
  \[
  \kappa\downarrow 0\ \Longrightarrow\ \delta(\kappa)\uparrow\delta_0,
  \qquad
  \kappa\uparrow\infty\ \Longrightarrow\ \delta(\kappa)\downarrow 0.
  \]
\end{enumerate}
\end{theorem}
\begin{proof}[证明要点]
对固定 \(\kappa>0\)，函数 \(A\mapsto u_{\kappa,A}(x)\) 连续且严格递减，因而 \(\Phi_\kappa\) 连续且严格递减。
同时 \(A\mapsto A-\Phi_\kappa(A)\) 连续严格递增，并满足：
当 \(A\downarrow 0\) 时 \(\Phi_\kappa(A)\to \kappa^{-1/2}\sum_x\sqrt{w(x)}>0\) 从而 \(A-\Phi_\kappa(A)<0\)；
当 \(A\uparrow\infty\) 时 \(u_{\kappa,A}(x)=w(x)/A+O(A^{-3})\) 从而 \(\Phi_\kappa(A)=A^{-1}+o(A^{-1})\) 且 \(A-\Phi_\kappa(A)>0\)。
故存在唯一不动点 \(A(\kappa)\)。
\par
由二次方程 \(\kappa u^2+Au-w=0\) 的正根公式可得
\(
w(x)=A(\kappa)\,u_\kappa(x)+\kappa u_\kappa(x)^2
\)，并由不动点条件 \(\sum_x u_\kappa(x)=A(\kappa)\) 得
\[
\sum_{y\in X}P_\kappa(x,y)=u_\kappa(x)\bigl(A(\kappa)+\kappa u_\kappa(x)\bigr)=w(x),
\]
从而双边缘为 \(w\)。
再由 \(\sum_x w(x)=1\) 推出
\(
1=A(\kappa)^2+\kappa\sum_x u_\kappa(x)^2
=\sum_{x,y}P_\kappa(x,y)
\)，故 \(P_\kappa\) 为概率分布。
\par
\(\delta(\kappa)=1-(1+\kappa)\sum_x u_\kappa(x)^2\)。
当 \(\kappa\downarrow 0\) 时，上述缩放退化到独立耦合并给出 \(\delta\uparrow \delta_0\)；
当 \(\kappa\uparrow\infty\) 时对角增强趋于同一耦合并给出 \(\delta\downarrow 0\)。
严格单调性可由严格凸性与对偶斜率表示推出：由命题 \ref{prop:pom-diagonal-rate-analytic-strict-convexity-curvature}，
\(\delta\mapsto\kappa(\delta)\) 在 \((0,\delta_0)\) 上严格递减，故其逆映射 \(\kappa\mapsto\delta(\kappa)\) 严格递减。证毕。
\end{proof}
\leanverified{paper_pom_diagonal_rate_scalar_collapse}

\begin{corollary}[一维闭式速率与代数回收：\texorpdfstring{$\kappa=\frac{1-A^2}{A^2-\delta}$}{kappa=(1-A^2)/(A^2-delta)}]\label{cor:pom-diagonal-rate-scalar-collapse-rate}
设 \(\delta\in(0,\delta_0)\)，并令 \(\kappa(\delta)\) 为定理 \ref{thm:pom-diagonal-rate-scalar-collapse} 的双射反函数（满足 \(\delta(\kappa(\delta))=\delta\)）。
记 \(A:=A(\kappa(\delta))\) 与 \(u:=u_{\kappa(\delta)}\)。
则 \(R_w(\delta)\) 的唯一极小耦合为 \(P_{\kappa(\delta)}\)，并且速率具有显式一维表示
\[
\boxed{\;
R_w(\delta)=2\sum_{x\in X}w(x)\log\frac{u(x)}{w(x)}+(1-\delta)\log\bigl(1+\kappa(\delta)\bigr)\; }.
\]
此外 \(A\in(\sqrt\delta,1)\)，并满足显式代数关系
\[
\boxed{\;
\kappa(\delta)=\frac{1-A^2}{A^2-\delta}\; } .
\]
\end{corollary}
\begin{proof}[证明要点]
由定理 \ref{thm:pom-diagonal-rate-scalar-collapse}，对给定 \(\delta\in(0,\delta_0)\) 存在唯一 \(\kappa(\delta)\) 使 \(P_{\kappa(\delta)}\) 可行且满足 \(\mathbb P(X=Y)=1-\delta\)；
由严格凸性（定理 \ref{thm:pom-diagonal-rate-explicit-param-and-derivative}）知其即为唯一最优耦合。
\par
由 \(\sum_{x,y}P_{\kappa}(x,y)=1\)（证明中已得）可取规范化常数 \(Z\equiv 1\)，
于是定理 \ref{thm:pom-diagonal-rate-explicit-param-and-derivative} 的闭式
\(
R_w(\delta)=2\sum_x w(x)\log\frac{u_x}{w(x)}-\log Z+(1-\delta)\log(1+\kappa)
\)
化为所示一维表示。
\par
最后由恒等式
\[
1-\delta=\sum_x P_\kappa(x,x)=(1+\kappa)\sum_x u_\kappa(x)^2,
\qquad
1=\sum_x w(x)=A^2+\kappa\sum_x u_\kappa(x)^2,
\]
消去 \(\sum_x u_\kappa(x)^2\) 得
\(
\delta=A^2-\frac{1-A^2}{\kappa}
\)，即 \(\kappa=(1-A^2)/(A^2-\delta)\)；并且 \(\kappa>0\) 迫使 \(A^2>\delta\)，又由 \(1=A^2+\kappa\sum u^2\) 得 \(A<1\)。证毕。
\end{proof}

\paragraph{阈值邻域的局部律与有限层代数闭合}

\begin{theorem}[碰撞阈值处的二次律与显式信息曲率]\label{thm:pom-diagonal-rate-collision-threshold-quadratic-law}
设 \(X\) 为有限集，取全支撑分布 \(w\in\Delta(X)\)，并记碰撞阈值
\[
\delta_0:=1-\sum_{x\in X}w(x)^2.
\]
定义
\[
\sigma_w^2
:=\sum_{x\in X}w(x)^2-2\sum_{x\in X}w(x)^3+\Bigl(\sum_{x\in X}w(x)^2\Bigr)^2.
\]
则 \(\sigma_w^2>0\)，且当 \(\delta\uparrow\delta_0\) 时有二次展开
\[
\boxed{\;
R_w(\delta)=\frac{(\delta_0-\delta)^2}{2\sigma_w^2}+o\!\bigl((\delta_0-\delta)^2\bigr)\; }.
\]
更进一步，令 \(\varepsilon:=\delta_0-\delta\downarrow 0\)，则存在一族可行耦合 \(P_\varepsilon\) 满足
\[
I_{P_\varepsilon}(X;Y)=\frac{\varepsilon^2}{2\sigma_w^2}+O(\varepsilon^3),
\qquad
\PP_{P_\varepsilon}(X=Y)=\sum_{x\in X}w(x)^2+\varepsilon,
\]
并且任何实现 \(R_w(\delta)\) 的最优耦合亦具有同一主项系数。
\leanverified{paper\_pom\_diagonal\_rate\_collision\_threshold\_quadratic\_law}
\end{theorem}
\begin{proof}[证明要点]
令基准独立耦合 \(P_0:=w\otimes w\)。由于 \(I_P(X;Y)=D_{\mathrm{KL}}(P\|w\otimes w)=D_{\mathrm{KL}}(P\|P_0)\)，故
\[
R_w(\delta)=\min\Bigl\{D_{\mathrm{KL}}(P\|P_0):\ P_X=P_Y=w,\ \EE_P[\mathbf 1_{\{X=Y\}}]\ge 1-\delta\Bigr\}.
\]
注意 \(\EE_{P_0}[\mathbf 1_{\{X=Y\}}]=\sum_x w(x)^2=1-\delta_0\)，从而当 \(\delta\ge \delta_0\) 时最优值为 \(0\)；以下仅考虑 \(\delta<\delta_0\) 的邻域。
\par
在 \(X\times X\) 上定义中心化统计量
\[
R(x,y):=\mathbf 1_{\{x=y\}}-w(x)-w(y)+\sum_{z\in X}w(z)^2.
\]
直接计算得到 \(\EE_{P_0}[R\mid X]=\EE_{P_0}[R\mid Y]=0\)，并且
\[
\EE_{P_0}[R]=0,
\qquad
\EE_{P_0}\!\bigl[R(X,Y)^2\bigr]=\sigma_w^2.
\]
由于 \(w\) 全支撑且 \(|X|\ge 2\)，\(R\) 非零非常值，故 \(\sigma_w^2=\Var_{P_0}(R(X,Y))>0\)。
\par
取 \(\varepsilon>0\) 充分小，并构造扰动耦合
\[
P_\varepsilon(x,y):=P_0(x,y)\Bigl(1+\frac{\varepsilon}{\sigma_w^2}R(x,y)\Bigr).
\]
由 \(\EE_{P_0}[R\mid X]=\EE_{P_0}[R\mid Y]=0\) 得两边缘仍为 \(w\)；又由
\[
\EE_{P_0}\!\bigl[\mathbf 1_{\{X=Y\}}R(X,Y)\bigr]=\sigma_w^2
\]
得到 \(\PP_{P_\varepsilon}(X=Y)=\PP_{P_0}(X=Y)+\varepsilon\)，从而 \(\PP_{P_\varepsilon}(X=Y)=1-\delta\)。当 \(\varepsilon\) 足够小时 \(P_\varepsilon\ge 0\)，故其可行。
\par
令 \(Z:=\frac{\varepsilon}{\sigma_w^2}R(X,Y)\)，则 \(\EE_{P_0}[Z]=0\) 且 \(\EE_{P_0}[Z^2]=\varepsilon^2/\sigma_w^2\)。对相对熵作泰勒展开：
\[
D_{\mathrm{KL}}(P_\varepsilon\|P_0)
=\EE_{P_0}\bigl[(1+Z)\log(1+Z)\bigr]
=\frac12\EE_{P_0}[Z^2]+O\!\bigl(\EE_{P_0}[|Z|^3]\bigr)
=\frac{\varepsilon^2}{2\sigma_w^2}+O(\varepsilon^3),
\]
从而给出上界。
\par
反向下界由 KL 在 \(P_0\) 邻域的二次型主导给出：在边缘约束与线性约束 \(\sum_x(P(x,x)-P_0(x,x))=\varepsilon\) 下，
\(
D_{\mathrm{KL}}(P\|P_0)
\)
的二阶变分是加权 \(\ell^2(P_0^{-1})\) 范数，其最小方向与 \(P_0R\) 共线，极小二次值恰为 \(\varepsilon^2/(2\sigma_w^2)\)；
结合上界即得二次律与主项系数的唯一性。证毕。
\end{proof}

\begin{theorem}[有限层分布的代数闭合：最优耦合由有限次开方确定]\label{thm:pom-diagonal-rate-finite-layer-algebraic}
设 \(w\) 仅取 \(r\) 个不同概率值 \(p_1,\dots,p_r\in(0,1)\)，其中 \(p_\ell\) 的重数为 \(m_\ell\)（\(\sum_{\ell=1}^r m_\ell=|X|\)）。
取任意 \(\delta\in(0,\delta_0)\)（\(\delta_0:=1-\sum_x w(x)^2\)）。
则达到 \(R_w(\delta)\) 的唯一最优耦合 \(P^\star_\delta\) 可写成如下闭式：存在唯一 \(t>1\) 与唯一向量 \(v\in(0,\infty)^X\) 使得
\[
\boxed{\;
P^\star_\delta(x,y)=
\begin{cases}
v(x)v(y), & x\ne y,\\[2mm]
t\,v(x)^2, & x=y,
\end{cases}
\qquad
\sum_{y\in X}P^\star_\delta(x,y)=w(x)\ \ (\forall x\in X).\;}
\]
并且 \(v(x)\) 仅依赖于 \(w(x)\)，因而 \(v\) 亦仅取 \(r\) 个不同数值 \(x_1,\dots,x_r\)（分别对应各概率层）。
记
\[
V:=\sum_{x\in X}v(x)=\sum_{\ell=1}^r m_\ell x_\ell,
\]
则 \((t,V)\) 必满足系统
\[
\boxed{\;
V^2=\frac{1+\delta(t-1)}{t},\qquad
\sum_{\ell=1}^r m_\ell\sqrt{V^2+4(t-1)p_\ell}=(|X|+2t-2)V.\;}
\]
因此 \(t\) 为一元代数方程的唯一实根：存在系数属于 \(\QQ(\delta,p_1,\dots,p_r)\) 的多项式 \(P_{w,\delta}\) 使
\[
\deg P_{w,\delta}\le 2^r,\qquad P_{w,\delta}(t)=0,\qquad t\in(1,\infty)\ \text{唯一}.
\]
特别地，在 \(\delta\in(0,\delta_0)\) 上，\(\delta\mapsto R_w(\delta)\) 作为复解析延拓意义下的一支代数函数，其代数度不超过 \(2^r\)。
\end{theorem}
\begin{proof}[证明要点]
由定理 \ref{thm:pom-diagonal-rate-explicit-param-and-derivative}，在 \(\delta\in(0,\delta_0)\) 上最优耦合唯一且具有对角增强的指数族形状；
在 \(Z\equiv 1\) 的规范下可写为 \(P^\star_\delta(x,y)=v(x)v(y)\)（\(x\ne y\)）与 \(P^\star_\delta(x,x)=t\,v(x)^2\)（\(t>1\)），并满足行和约束。
行和方程等价于
\[
w(x)=v(x)\bigl(V+(t-1)v(x)\bigr),
\]
故 \(v(x)\) 是关于 \(w(x)\) 的二次方程正根；因此同一概率层上的 \(v(x)\) 取同一常数 \(x_\ell\)。
由正根公式得到
\[
x_\ell=\frac{-V+\sqrt{V^2+4(t-1)p_\ell}}{2(t-1)}.
\]
将其代入 \(V=\sum_\ell m_\ell x_\ell\) 即得含 \(r\) 个平方根的方程。
\par
另一方面，由 \(\sum_x w(x)=1\) 与对角和约束 \(1-\delta=\sum_x P^\star_\delta(x,x)\) 消去 \(\sum_x v(x)^2\) 可得
\(
V^2=\frac{1+\delta(t-1)}{t}
\)。
将 \(V^2\) 的有理表达代入根号方程，并反复平方消元（至多 \(r\) 次）即可得到关于 \(t\) 的多项式方程，次数不超过 \(2^r\)；
唯一性来自原问题严格凸导致的最优解唯一性。证毕。
\end{proof}
\leanverified{paper\_pom\_diagonal\_rate\_finite\_layer\_algebraic}

\begin{corollary}[两层分布的四次闭合解与显式最优耦合]\label{cor:pom-diagonal-rate-two-layer-quartic}
在定理 \ref{thm:pom-diagonal-rate-finite-layer-algebraic} 的设定下取 \(r=2\)，记两层概率为 \(p_1\ne p_2\)，重数分别为 \(m_1,m_2\)（\(m:=|X|=m_1+m_2\)）。
令 \(s:=t-1>0\) 并记
\[
V^2=\frac{1+s\delta}{1+s},
\qquad
A:=V^2+4sp_1,\qquad B:=V^2+4sp_2.
\]
则最优参数 \(s\) 等价于满足四次方程
\[
\boxed{\;
\Bigl((m+2s)^2V^2-m_1^2A-m_2^2B\Bigr)^2=4m_1^2m_2^2AB.\;}
\]
该方程在 \(s>0\) 上存在唯一解 \(s^\star\)。
由此唯一确定
\[
x_\ell=\frac{-V+\sqrt{V^2+4s^\star p_\ell}}{2s^\star}\quad(\ell=1,2),
\]
并得到唯一最优耦合：若 \(x\) 属于第 \(\ell\) 层，则 \(v(x)=x_\ell\)，且
\[
P^\star_\delta(x,y)=
\begin{cases}
v(x)v(y), & x\ne y,\\
(1+s^\star)v(x)^2, & x=y.
\end{cases}
\]
\end{corollary}
\begin{proof}
由定理 \ref{thm:pom-diagonal-rate-finite-layer-algebraic} 的根号方程在 \(r=2\) 情形化为
\(
m_1\sqrt{A}+m_2\sqrt{B}=(m+2s)V
\)。
两侧平方并再平方以消去 \(\sqrt{AB}\) 即得四次方程；唯一性由最优解唯一性推出。证毕。
\end{proof}

\begin{proposition}[有限层情形的代数相图：奇点集合有限且具 Puiseux 展开]\label{prop:pom-diagonal-rate-finite-layer-algebraic-phase-diagram}
在定理 \ref{thm:pom-diagonal-rate-finite-layer-algebraic} 的 \(r\)-层设定下，存在有限集合 \(\Sigma\subset(0,\delta_0)\) 使得：
\begin{enumerate}
  \item \(\delta\mapsto R_w(\delta)\) 在每个连通分支 \((0,\delta_0)\setminus\Sigma\) 上为实解析函数；
  \item 对每个 \(\delta_*\in\Sigma\)，\(R_w\) 在 \(\delta_*\) 的邻域内存在 Puiseux 型展开
  \[
  R_w(\delta)=\sum_{k\ge k_0} c_k(\delta-\delta_*)^{k/q},
  \]
  其中 \(q\in\NN\) 且 \(q\le 2^r\)。
\end{enumerate}
\end{proposition}
\begin{proof}
由定理 \ref{thm:pom-diagonal-rate-finite-layer-algebraic}，\(R_w(\delta)\) 为代数函数，满足某一不可约多项式关系 \(F(\delta,R)=0\)。
代数函数的分歧仅可能发生在同时满足 \(\partial_R F=\partial_\delta F=0\) 的投影处，等价地由判别式 \(\Disc_R F(\delta)=0\) 的有限根集合给出；
在其补集上由隐函数定理得实解析性。
分歧点的 Puiseux 展开是代数函数的经典局部结构定理之直接推论，且指数分母可由代数度界定（不超过 \(2^r\)）。证毕。
\end{proof}

\begin{theorem}[矩阵缩放对偶：\texorpdfstring{$R_w(\delta)$}{R} 的单参数极大与最优耦合回收]\label{thm:pom-diagonal-rate-scaling-duality}
设 \(X\) 为有限集，取全支撑分布 \(w\in\Delta(X)\)。
对 \(\kappa\ge 0\) 定义对角增强核
\[
K_\kappa(x,y):=1+\kappa\,\mathbf 1_{\{x=y\}},
\]
并对 \(u\in\RR_{>0}^{X}\) 定义
\[
Z_\kappa(u):=\sum_{x,y\in X}u(x)K_\kappa(x,y)u(y)
=(\sum_x u(x))^2+\kappa\sum_x u(x)^2.
\]
定义缩放能量泛函
\[
\mathcal G_\kappa(w):=\inf_{u\succ 0}\Bigl(\log Z_\kappa(u)-2\sum_{x\in X}w(x)\log u(x)\Bigr),
\]
其极小点在整体伸缩（\(u\mapsto cu\)）下唯一。
则对任意 \(\delta\in[0,1]\) 有对偶极大表示
\[
\boxed{\;
R_w(\delta)=\sup_{\kappa\ge 0}\Bigl(2H(w)-\mathcal G_\kappa(w)+(1-\delta)\log(1+\kappa)\Bigr)\;},
\qquad
H(w):=-\sum_x w(x)\log w(x).
\]
并且当 \(\delta\in(0,\delta_0)\)（\(\delta_0:=1-p_2(w)\)）时，上确界由唯一 \(\kappa^\star(\delta)>0\) 达到；
若 \(u^\star\) 为实现 \(\mathcal G_{\kappa^\star}(w)\) 的极小点之一，则
\[
\boxed{\;
P^\star_\delta(x,y):=\frac{u^\star(x)K_{\kappa^\star}(x,y)u^\star(y)}{Z_{\kappa^\star}(u^\star)}\;}
\]
即为 \(R_w(\delta)\) 的唯一极小耦合。
\end{theorem}
\begin{proof}[证明要点]
对固定 \(\kappa\ge 0\)，函数
\(
u\mapsto \log Z_\kappa(u)-2\langle w,\log u\rangle
\)
在变量 \(\log u\) 上凸，并在模去常数方向上严格凸，从而极小点在整体伸缩下唯一。
其一阶条件等价于存在常数 \(Z>0\) 使
\[
w(x)=\frac{u(x)\bigl(A+\kappa u(x)\bigr)}{Z},
\qquad A:=\sum_y u(y),
\]
并因此诱导一个满足双边缘为 \(w\) 的耦合
\(
P(x,y)\propto u(x)K_\kappa(x,y)u(y)
\)；
该结构与定理 \ref{thm:pom-diagonal-rate-explicit-param-and-derivative} 的指数族同型。
\par
对该耦合 \(P\) 展开相对熵并用
\(
\mathcal G_\kappa(w)=\log Z_\kappa(u)-2\sum_x w(x)\log u(x)
\)
消去 \(\log Z_\kappa\) 与 \(\sum w\log u\)，得到
\[
D(P\|w\otimes w)=2H(w)-\mathcal G_\kappa(w)+\mathbb P_P(X=Y)\log(1+\kappa).
\]
将约束 \(\mathbb P_P(X=Y)\ge 1-\delta\) 以乘子 \(\log(1+\kappa)\) 外提并取上确界，即得所示对偶形式；强对偶、唯一性与 \(\kappa^\star>0\) 由原问题严格凸与 Slater 条件给出。
把 \(\kappa^\star\) 与任一对应极小点 \(u^\star\) 回代即可回收唯一最优耦合。证毕。
\end{proof}

\paragraph{速率曲线的指纹化与分布可识别性}
定理 \ref{thm:pom-diagonal-rate-scalar-collapse} 与定理 \ref{thm:pom-diagonal-rate-scaling-duality} 表明：在 \(\delta\in(0,\delta_0)\) 上，最优耦合由单一对偶标量 \(\kappa(\delta)\) 及其规范化缩放向量 \(u_\delta\) 完全刻画。
以下给出一个更强的结论：若把 \(\delta\mapsto R_w(\delta)\) 视为函数级观测，则它本身已构成 \(w\) 的完备指纹（除去标签置换的不变性）。

\begin{theorem}[单曲线可识别性：\texorpdfstring{$\delta\mapsto R_w(\delta)$}{R(delta)} 唯一决定 \texorpdfstring{$w$}{w}（至置换）]\label{thm:pom-diagonal-rate-rate-curve-identifiability}
令 \(X\) 为有限集，\(|X|=n\ge 2\)，并取全支撑分布 \(w\in\Delta(X)\)。
记 \(\delta_0(w):=1-\sum_{x\in X}w(x)^2\)。
若另一全支撑分布 \(v\in\Delta(X)\) 满足：存在 \((0,\delta_1)\subset\bigl(0,\min\{\delta_0(w),\delta_0(v)\}\bigr)\) 使
\[
R_w(\delta)=R_v(\delta)\qquad(\forall \delta\in(0,\delta_1)),
\]
则两分布的权重多重集一致：
\[
\{w(x):x\in X\}=\{v(x):x\in X\},
\]
从而 \(w=v\)（至标签置换）。
\end{theorem}
\begin{proof}[证明要点]
由命题 \ref{prop:pom-diagonal-rate-analytic-strict-convexity-curvature}，\(\delta\mapsto R_w(\delta)\) 与 \(\delta\mapsto R_v(\delta)\) 在各自的 \((0,\delta_0)\) 上实解析。
因此在含聚点区间 \((0,\delta_1)\) 上的恒等式迫使二者在公共定义域上处处相等。
\par
由定理 \ref{thm:pom-diagonal-rate-explicit-param-and-derivative} 的对偶敏感性，
\[
R_w'(\delta)=-\log\bigl(1+\kappa_w(\delta)\bigr),
\qquad
R_v'(\delta)=-\log\bigl(1+\kappa_v(\delta)\bigr),
\]
从而由 \(R_w\equiv R_v\) 得 \(\kappa_w(\delta)\equiv \kappa_v(\delta)\)。
再由定理 \ref{thm:pom-diagonal-rate-scalar-collapse}，映射 \(\kappa\mapsto\delta(\kappa)\) 严格递减并在 \((0,\infty)\) 与 \((0,\delta_0)\) 间成双射，
故其反函数 \(\delta\mapsto \kappa(\delta)\) 唯一；于是两者诱导的 \(\delta(\kappa)\) 在 \(\kappa>0\) 上一致。
\par
令 \(A(\kappa)\) 为定理 \ref{thm:pom-diagonal-rate-scalar-collapse} 的唯一不动点解，并取规范化缩放 \(Z\equiv 1\)（同定理）。
由推论 \ref{cor:pom-diagonal-rate-scalar-collapse-rate} 的代数关系可得
\[
A(\kappa)^2=\frac{1+\kappa\,\delta(\kappa)}{1+\kappa}\qquad(\kappa>0),
\]
从而由已知的 \(\delta(\kappa)\) 唯一确定函数 \(\kappa\mapsto A(\kappa)\)。
\par
在 \(\kappa\downarrow 0\) 时有 \(A(\kappa)\uparrow 1\) 且 \(A\) 可解析延拓到 \(\kappa=0\) 的邻域。
将定理 \ref{thm:pom-diagonal-rate-scalar-collapse} 的不动点方程
\(
A=\Phi_\kappa(A)=\sum_x u_{\kappa,A}(x)
\)
视为关于 \(\kappa\) 的解析恒等式，并对
\(
u_{\kappa,A}(x)=(\sqrt{A^2+4\kappa w(x)}-A)/(2\kappa)
\)
在 \(\kappa=0\) 处作 Taylor 展开，可得到
\[
\Phi_\kappa(A)=\frac{1}{A}-\kappa\,\frac{p_2(w)}{A^3}+O(\kappa^2),
\qquad p_2(w):=\sum_x w(x)^2,
\]
从而由系数比较给出 \(p_2(w)=-2A'(0)\)。
更一般地，\(\Phi_\kappa(A)\) 的 \(\kappa^{q-1}\) 系数线性包含幂和 \(p_q(w)=\sum_x w(x)^q\) 且仅依赖于 \(A\) 的低阶导数与 \(\{p_j(w)\}_{2\le j\le q-1}\)；
因此由归纳可从 \(\kappa\mapsto A(\kappa)\) 在 \(\kappa=0\) 的 Taylor jet 逐阶恢复 \(\{p_q(w)\}_{q=2}^{n}\)。
\par
最后，由 Newton 恒等式，已知幂和 \(\{p_q(w)\}_{q=1}^{n}\)（其中 \(p_1\equiv 1\)）等价于已知基本对称多项式 \(\{e_k(w)\}_{k=1}^{n}\)，从而确定一元多项式
\(
t^n-e_1 t^{n-1}+\cdots+(-1)^n e_n
\)
的全部根；该根多重集恰为 \(\{w(x):x\in X\}\)。
因此 \(R_w\) 唯一决定 \(w\) 的权重多重集（至置换）。证毕。
\end{proof}

\begin{corollary}[两标量证书的完备性：\texorpdfstring{$R_w(\delta)$}{R(delta)} 与 \texorpdfstring{$\lambda_2(K^\star_\delta)$}{lambda2}]\label{cor:pom-diagonal-rate-two-curve-identifiability}
令 \(X\) 为有限集，\(w\in\Delta(X)\) 全支撑。
对 \(\delta\in(0,\delta_0(w))\) 令 \(K^\star_\delta\) 为定理 \ref{thm:pom-maxent-markov-unique-optimal-kernel} 的唯一最优核，并令 \(\lambda_2(K^\star_\delta)\) 为其第二特征值（按 \(\lambda_1=1\ge \lambda_2\ge\cdots\ge \lambda_n\) 排序）。
若已知两条曲线
\(
\delta\mapsto R_w(\delta)
\)
与
\(
\delta\mapsto \lambda_2(K^\star_\delta)
\)
在某个含聚点区间上的全部值，则可唯一确定 \(\{w(x)\}\)（至置换）。
\end{corollary}
\begin{proof}
由定理 \ref{thm:pom-diagonal-rate-rate-curve-identifiability}，\(\delta\mapsto R_w(\delta)\) 已足以唯一确定 \(\{w(x)\}\)；
因此在此基础上再给出 \(\lambda_2(K^\star_\delta)\) 不会引入额外歧义。证毕。
\end{proof}

\begin{proposition}[从速率曲线回收权重多重集：幂和--Newton 重构]\label{prop:pom-diagonal-rate-reconstruct-w-from-rate-curve}
在定理 \ref{thm:pom-diagonal-rate-rate-curve-identifiability} 的设定下，给定 \(n=|X|\) 与函数 \(\delta\mapsto R_w(\delta)\) 在任意含聚点区间上的全部值，可按如下程序回收权重多重集 \(\{w(x)\}\)：
\begin{enumerate}
  \item 由实解析性将 \(R_w\) 延拓到 \((0,\delta_0)\) 的全域，并由对偶敏感性计算
  \[
  \kappa(\delta):=\exp\bigl(-R_w'(\delta)\bigr)-1.
  \]
  \item 反解 \(\delta\mapsto\kappa(\delta)\) 得到 \(\delta(\kappa)\)（\(\kappa\downarrow 0\) 对应 \(\delta\uparrow\delta_0\)），并定义
  \[
  A(\kappa):=\sqrt{\frac{1+\kappa\,\delta(\kappa)}{1+\kappa}}.
  \]
  \item 由不动点方程 \(A(\kappa)=\Phi_\kappa(A(\kappa))\) 在 \(\kappa=0\) 的 Taylor 展开逐阶恢复幂和 \(\{p_q(w)\}_{q=2}^{n}\)（例如 \(p_2(w)=-2A'(0)\)）。
  \item 由 Newton 恒等式由 \(\{p_q(w)\}_{q=1}^{n}\) 计算基本对称多项式 \(\{e_k(w)\}_{k=1}^{n}\)，并以
  \[
  t^n-e_1 t^{n-1}+\cdots+(-1)^n e_n
  \]
  的根多重集作为 \(\{w(x)\}\) 的输出。
\end{enumerate}
该过程仅依赖于 \(\delta\mapsto R_w(\delta)\) 作为函数输入，并在理论上实现对 \(w\) 的排序谱回收。
\end{proposition}


\begin{corollary}[零比特走廊：碰撞概率决定无外置信息区]\label{cor:pom-query-distortion-zero-bit-corridor}
在定理 \ref{thm:pom-microcanonical-query-distortion-strong-converse-plane} 的设定下，
若 \(\delta\ge 1-p_2(w)\)，则 \(R_w(\delta)=0\) 且 \(b_{\mathrm{crit}}(\beta,\delta)=0\) 对所有 \(\beta\in[0,1)\) 成立；
若 \(\delta<1-p_2(w)\)，则 \(R_w(\delta)>0\) 且对任意 \(\beta<1\) 有 \(b_{\mathrm{crit}}(\beta,\delta)>0\)。
\end{corollary}
\begin{proof}
由命题 \ref{prop:pom-diagonal-rate-zero-threshold-collision}，\(\delta\ge 1-p_2(w)\) 当且仅当 \(R_w(\delta)=0\)。
代入定理 \ref{thm:pom-microcanonical-query-distortion-strong-converse-plane} 的临界平面公式即得。
证毕。
\end{proof}

\begin{corollary}[\texorpdfstring{$\ell^2$}{l2} 直方图不变量：零比特失真阈值的可审计性]\label{cor:pom-query-distortion-zero-bit-threshold-auditable}
在 Fold 口径下，推前分布为 \(w_m(x)=d_m(x)/2^m\)。
定义二阶碰撞概率（\(\ell^2\) 直方图）
\[
\mathrm{Col}_m:=\sum_{x\in X_m}w_m(x)^2=\frac{S_2(m)}{2^{2m}}.
\]
则零比特可达失真阈值
\[
\delta_{0,m}:=\inf\{\delta:\ R_{w_m}(\delta)=0\}
\]
满足严格等式
\[
\boxed{\;
\delta_{0,m}=1-\mathrm{Col}_m
=1-\frac{S_2(m)}{2^{2m}}\; } .
\]
若进一步存在极限 \(\mathrm{Col}_m\to \mathrm{Col}_\infty\)，则 \(\delta_{0,m}\to 1-\mathrm{Col}_\infty\)。
\end{corollary}
\begin{proof}
由命题 \ref{prop:pom-diagonal-rate-zero-threshold-collision}，
\(
\delta_{0,m}=1-p_2(w_m)
\)。
而 \(p_2(w_m)=\sum_x w_m(x)^2=\mathrm{Col}_m=2^{-2m}S_2(m)\)，代入即得闭式。
极限结论由连续性立即推出。证毕。
\end{proof}

\begin{theorem}[双边缘对角约束速率的 Fr\'echet 导数：由缩放势给出显式梯度]\label{thm:pom-diagonal-rate-frechet-derivative-scaling-potential}
设 \(X\) 为有限字母表（\(|X|\ge 2\)），取全支撑分布 \(w\in\Delta(X)\)，并令 \(\delta\in(0,1-p_2(w))\)。
令 \(P^\star_\delta\) 为定理 \ref{thm:pom-diagonal-rate-explicit-param-and-derivative} 的唯一最优耦合，并取其指数族表示
\[
P^\star_\delta(x,y)=\frac{u_\delta(x)u_\delta(y)\,\exp\!\bigl(\lambda_\delta\,\mathbf 1_{\{x=y\}}\bigr)}{Z_\delta},
\qquad
u_\delta\in\RR_{>0}^{X},\ \lambda_\delta>0,
\]
其中 \(u_\delta\) 在整体伸缩下唯一（伸缩吸收进 \(Z_\delta\)）。
定义缩放势
\[
\psi_\delta(x):=\log\frac{u_\delta(x)}{w(x)}\qquad(x\in X),
\]
其在加常数意义下唯一。
则映射 \(w\mapsto R_w(\delta)\) 在单纯形内点处 Fr\'echet 可微；对任意切向扰动 \(\dot w\)（满足 \(\sum_{x\in X}\dot w(x)=0\)）成立
\[
DR_w(\delta)[\dot w]=2\sum_{x\in X}\dot w(x)\,\psi_\delta(x).
\]
因此 \(R_w(\delta)\) 的梯度（切空间意义下）可取为
\[
\nabla R_w(\delta)\equiv 2\psi_\delta\pmod{\RR\mathbf 1}.
\]
\end{theorem}
\begin{proof}
当 \(P_X=P_Y=w\) 时，互信息满足恒等式
\(
I_P(X;Y)=H(w)+H(w)-H(P)
=2H(w)-H(P)
\)。
因此
\[
R_w(\delta)
=2H(w)-\sup\Bigl\{H(P):\ P_X=P_Y=w,\ \mathbb P_P(X=Y)\ge 1-\delta\Bigr\}.
\]
上确界问题对 \(P\) 为严格凹规划，且在 \(\delta<1-p_2(w)\) 时满足 Slater 条件；故极大元 \(P^\star_\delta\) 唯一，且 KKT 乘子在加常数规约下唯一。
\par
对极大问题写拉格朗日函数（以 \(\alpha,\beta\) 为两条边缘约束乘子、以 \(\lambda\) 为对角约束乘子）并对 \(P(x,y)\) 作驻点条件，可得
\[
P^\star_\delta(x,y)\propto \exp\!\bigl(-\alpha_x-\beta_y+\lambda\,\mathbf 1_{\{x=y\}}\bigr).
\]
令 \(u_\delta(x):=\exp(-\alpha_x)\) 且由交换对称性取 \(\alpha=\beta\)（相差常数吸收进 \(Z_\delta\)），即得到定理陈述中的指数族表示。
\par
记
\(
F(w):=\sup\{H(P):P_X=P_Y=w,\ \mathbb P(X=Y)\ge 1-\delta\}
\)。
由于极大元唯一且满足常规正则性，包络定理给出
\[
DF(w)[\dot w]=\sum_{x\in X}(\alpha_x+\beta_x)\,\dot w(x)
=-2\sum_{x\in X}\dot w(x)\,\log u_\delta(x),
\]
其中最后一步使用 \(\alpha_x=\beta_x=-\log u_\delta(x)\)（加常数不影响切向配对）。
又因 \(H(w)=-\sum_x w(x)\log w(x)\)，对切向扰动有 \(DH(w)[\dot w]=-\sum_x\dot w(x)\log w(x)\)。
代回 \(R_w(\delta)=2H(w)-F(w)\) 得
\[
DR_w(\delta)[\dot w]
=2DH(w)[\dot w]-DF(w)[\dot w]
=2\sum_{x\in X}\dot w(x)\Bigl(\log u_\delta(x)-\log w(x)\Bigr)
=2\sum_{x\in X}\dot w(x)\,\psi_\delta(x),
\]
即所示。证毕。
\end{proof}

\begin{corollary}[\texorpdfstring{$\sqrt N$}{sqrtN} 比特窗的显式方差常数]\label{cor:pom-microcanonical-query-distortion-critical-window-variance-psi}
在定理 \ref{thm:pom-microcanonical-query-distortion-critical-window-clt} 的设定下，若 \(\delta\in(0,1-p_2(w))\) 且 \(w\) 全支撑，则可取
\(
g_\delta\equiv 2\psi_\delta\pmod{\RR\mathbf 1}
\)
（\(\psi_\delta\) 见定理 \ref{thm:pom-diagonal-rate-frechet-derivative-scaling-potential}）。
从而临界窗的方差常数具有闭式
\[
\sigma_{\beta,\delta}^2
=\beta(1-\beta)\ \Var_{X\sim w}\!\bigl(2\psi_\delta(X)\bigr)
=4\beta(1-\beta)\Bigl(\EE_w[\psi_\delta(X)^2]-\EE_w[\psi_\delta(X)]^2\Bigr).
\]
\end{corollary}
\begin{proof}
由定理 \ref{thm:pom-diagonal-rate-frechet-derivative-scaling-potential}，\(w\mapsto R_w(\delta)\) 的梯度在切空间意义下可取 \(g_\delta=2\psi_\delta\)。
将其代入定理 \ref{thm:pom-microcanonical-query-distortion-critical-window-clt} 的
\(
\sigma_{\beta,\delta}^2=\beta(1-\beta)\ \widetilde g_\delta^{\mathsf T}(\mathrm{diag}(w)-ww^{\mathsf T})\widetilde g_\delta
\)
并注意 \(\widetilde g_\delta\) 即为对 \(w\) 去均值后的 \(2\psi_\delta\)，从而二次型等于 \(\Var_w(2\psi_\delta(X))\)。
证毕。
\end{proof}

% (User input window)

\paragraph{可解情形、双侧夹逼与协同不等式}
本段把定理 \ref{thm:pom-diagonal-rate-explicit-param-and-derivative} 的一般指数族结构在若干对称场景中完全显式化，并给出仅依赖有限统计量的通用下界与可达上界；最后刻画线性混合耦合的最优性条件与均匀字母表上的协同严格不等式。

\begin{theorem}[二点字母表的闭式与唯一最优耦合]\label{thm:pom-diagonal-rate-binary-closed-form}
令 \(X=\{0,1\}\)，取 \(w(0)=p\in(0,1)\)、\(w(1)=1-p\)。
记
\[
p_2(w)=p^2+(1-p)^2,\qquad \delta_0:=1-p_2(w)=2p(1-p).
\]
则对任意 \(\delta\in[0,1]\) 有
\[
R_w(\delta)=
\begin{cases}
0, & \delta\ge \delta_0,\\[4pt]
\displaystyle (p-\tfrac{\delta}{2})\log\frac{p-\delta/2}{p^2}
+\bigl((1-p)-\tfrac{\delta}{2}\bigr)\log\frac{(1-p)-\delta/2}{(1-p)^2}
+\delta\log\frac{\delta/2}{p(1-p)}, & 0\le \delta<\delta_0,
\end{cases}
\]
且当 \(0\le\delta<\delta_0\) 时最优耦合唯一并为
\[
P^\star_\delta=
\begin{pmatrix}
p-\delta/2 & \delta/2\\
\delta/2 & (1-p)-\delta/2
\end{pmatrix}.
\]
\leanverified{paper\_pom\_diagonal\_rate\_binary\_closed\_form}
\end{theorem}
\begin{proof}
在 \(2\times 2\) 情形，边缘约束 \(P_X=P_Y=w\) 的耦合集合为一参数族：
令 \(a:=P(0,0)\)，则
\[
P(a)=
\begin{pmatrix}
a & p-a\\
p-a & 1-2p+a
\end{pmatrix}.
\]
其对角质量为 \(\mathbb P(X=Y)=1-2p+2a\)，故约束 \(\mathbb P(X=Y)\ge 1-\delta\) 等价于 \(a\ge p-\delta/2\)。
由于边缘固定，互信息满足 \(I_{P(a)}(X;Y)=D_{\mathrm{KL}}(P(a)\,\|\,w\otimes w)\)，且作为 \(a\) 的函数严格凸并在独立耦合 \(a=p^2\) 处取唯一极小 \(0\)。
当 \(\delta\ge \delta_0\) 时有 \(p-\delta/2\le p^2\)，独立耦合可行，故 \(R_w(\delta)=0\)。
当 \(0\le \delta<\delta_0\) 时可行区间的左端点满足 \(p-\delta/2>p^2\)，严格凸性迫使唯一极小点落在边界 \(a=p-\delta/2\)，从而得到 \(P^\star_\delta\)。
将 \(P^\star_\delta\) 代入 \(D_{\mathrm{KL}}(P^\star_\delta\,\|\,w\otimes w)=\sum_{x,y}P^\star_\delta(x,y)\log\frac{P^\star_\delta(x,y)}{w(x)w(y)}\) 即得闭式表达。证毕。
\end{proof}

\begin{corollary}[二点字母表的线性混合结构]\label{cor:pom-diagonal-rate-binary-linear-mixture}
\leanverified{paper\_pom\_diagonal\_rate\_binary\_linear\_mixture}
在定理 \ref{thm:pom-diagonal-rate-binary-closed-form} 的设定下，若 \(0\le\delta<\delta_0\)，则最优耦合可唯一写成
\[
P^\star_\delta(x,y)=(1-\varepsilon)\,w(x)w(y)+\varepsilon\,w(x)\mathbf 1_{\{x=y\}},
\qquad
\varepsilon:=1-\frac{\delta}{\delta_0}\in(0,1].
\]
\end{corollary}
\begin{proof}
由定理 \ref{thm:pom-diagonal-rate-binary-closed-form} 的矩阵形式直接计算得
\(
P^\star_\delta(0,1)=P^\star_\delta(1,0)=\delta/2
\)。
另一方面，线性混合右端的非对角项为 \((1-\varepsilon)p(1-p)\)，令其等于 \(\delta/2\) 得 \(\varepsilon=1-\delta/(2p(1-p))=1-\delta/\delta_0\)。
用该 \(\varepsilon\) 代回对角项即与 \(P^\star_\delta\) 一致。
若存在另一组系数给出同一分解，则由非对角项立刻推出 \(\varepsilon\) 唯一。证毕。
\end{proof}

\begin{theorem}[均匀分布的闭式与唯一对称最优耦合]\label{thm:pom-diagonal-rate-uniform-closed-form}
令有限集 \(X\) 满足 \(|X|=n\ge 2\)，取均匀分布 \(U_n(x)=1/n\)。
记 \(\delta_0:=1-1/n\)。
则对任意 \(\delta\in[0,1]\) 有
\[
R_{U_n}(\delta)=
\begin{cases}
0, & \delta\ge 1-\frac{1}{n},\\[6pt]
(1-\delta)\log\bigl(n(1-\delta)\bigr)+\delta\log\!\Bigl(\dfrac{\delta n}{n-1}\Bigr), & 0\le\delta<1-\frac{1}{n},
\end{cases}
\]
且当 \(0\le\delta<1-\frac1n\) 时最优耦合唯一并满足
\[
P^\star_\delta(x,x)=\frac{1-\delta}{n},\qquad
P^\star_\delta(x,y)=\frac{\delta}{n(n-1)}\quad(x\ne y).
\]
\leanverified{paper\_pom\_diagonal\_rate\_uniform\_closed\_form}
\end{theorem}
\begin{proof}
对固定 \(\delta\)，可行集合在 \(X\) 的置换群作用下不变，而 \(P\mapsto I_P(X;Y)\) 在该凸集上严格凸。
对任意可行 \(P\) 取置换平均 \(\bar P\) 保持边缘与对角质量不变，并使互信息不增；因此最优解必为置换不变耦合。
在置换不变类中，双边缘均匀迫使对角项为常数、非对角项为常数，并由 \(\sum_x P(x,x)=1-\delta\) 唯一确定 \(P^\star_\delta\)。
代入
\(
I_{P^\star_\delta}(X;Y)=\sum_{x,y}P^\star_\delta(x,y)\log\frac{P^\star_\delta(x,y)}{1/n^2}
\)
得到所示闭式。证毕。
\end{proof}

\begin{corollary}[熵型重写]\label{cor:pom-diagonal-rate-uniform-entropy-rewrite}
在定理 \ref{thm:pom-diagonal-rate-uniform-closed-form} 的设定下，令二元熵
\[
H_2(\delta):=-(1-\delta)\log(1-\delta)-\delta\log\delta.
\]
则对 \(0\le\delta<1-\tfrac1n\) 有
\[
R_{U_n}(\delta)=\log n-H_2(\delta)-\delta\log(n-1).
\]
\end{corollary}
\begin{proof}
将定理 \ref{thm:pom-diagonal-rate-uniform-closed-form} 的闭式展开并整理即得。证毕。
\end{proof}

\begin{theorem}[均匀字母表上的协同严格不等式]\label{thm:pom-diagonal-rate-uniform-synergy-inequality}
令 \(n\ge 2\)、\(k\ge 2\)，并把 \(U_{n^k}\) 视为 \(n^k\) 元均匀分布。
对任意 \(\delta\in\bigl(0,1-1/n^k\bigr)\) 有严格不等式
\[
R_{U_{n^k}}(\delta)\ <\ k\,R_{U_n}(\delta).
\]
\leanverified{paper\_pom\_diagonal\_rate\_uniform\_synergy\_inequality}
\end{theorem}
\begin{proof}
由推论 \ref{cor:pom-diagonal-rate-uniform-entropy-rewrite}，
\[
R_{U_{n^k}}(\delta)=k\log n-H_2(\delta)-\delta\log(n^k-1),
\qquad
kR_{U_n}(\delta)=k\log n-kH_2(\delta)-k\delta\log(n-1).
\]
故差为
\[
kR_{U_n}(\delta)-R_{U_{n^k}}(\delta)
=(k-1)H_2(\delta)+\delta\bigl(\log(n^k-1)-k\log(n-1)\bigr).
\]
其中 \((k-1)H_2(\delta)>0\)，且 \(n^k-1>(n-1)^k\) 蕴含 \(\log(n^k-1)-k\log(n-1)>0\)，从而差严格为正。证毕。
\end{proof}

\begin{theorem}[对角粗粒化下界与线性可达上界]\label{thm:pom-diagonal-rate-bernoulli-lower-and-mixture-upper}
令 \(X\) 为有限字母表，取全支撑分布 \(w\in\Delta(X)\)，并记
\[
p_2(w):=\sum_{x\in X}w(x)^2,\qquad \delta_0:=1-p_2(w),\qquad H(w):=-\sum_{x\in X}w(x)\log w(x).
\]
则对任意 \(\delta\in[0,1]\) 成立下界
\[
R_w(\delta)\ \ge\ D_{\mathrm{KL}}\!\bigl(\mathrm{Bern}(1-\delta)\,\|\,\mathrm{Bern}(p_2(w))\bigr),
\]
其中 \(D_{\mathrm{KL}}(\mathrm{Bern}(a)\|\mathrm{Bern}(b)):=a\log\frac{a}{b}+(1-a)\log\frac{1-a}{1-b}\)。
\par
并且对任意 \(\delta\in[0,\delta_0]\)，令
\[
\varepsilon(\delta):=1-\frac{\delta}{\delta_0}\in[0,1],\qquad
P_\delta(x,y):=(1-\varepsilon(\delta))\,w(x)w(y)+\varepsilon(\delta)\,w(x)\mathbf 1_{\{x=y\}},
\]
则 \(P_\delta\) 可行且满足
\[
R_w(\delta)\ \le\ I_{P_\delta}(X;Y)\ \le\ \varepsilon(\delta)\,H(w)
=\Bigl(1-\frac{\delta}{1-p_2(w)}\Bigr)H(w).
\]
当 \(\delta\ge\delta_0\) 时有 \(R_w(\delta)=0\)。
\end{theorem}
\begin{proof}
对下界，令 \(f(x,y)=\mathbf 1_{\{x=y\}}\) 并记 \(Z=f(X,Y)\)。
对任意满足 \(P_X=P_Y=w\) 的耦合 \(P\)，有
\(
I_P(X;Y)=D_{\mathrm{KL}}(P\|w\otimes w)
\)；
由数据处理不等式
\[
D_{\mathrm{KL}}(P\|w\otimes w)\ \ge\ D_{\mathrm{KL}}\!\bigl(\mathrm{Law}_P(Z)\,\|\,\mathrm{Law}_{w\otimes w}(Z)\bigr).
\]
而 \(\mathrm{Law}_{w\otimes w}(Z)=\mathrm{Bern}(p_2(w))\)，并且可行性给出 \(\mathbb P_P(Z=1)=\mathbb P_P(X=Y)\ge 1-\delta\)。
当 \(1-\delta\le p_2(w)\) 时右端可取为 \(0\)；当 \(1-\delta>p_2(w)\) 时，\(\mathbb P_P(Z=1)\ge 1-\delta\) 的下确界在边界处取得，得到所示下界。
\par
对上界，直接计算得 \(P_\delta\) 的边缘仍为 \(w\)，且其对角质量为
\(
\sum_x P_\delta(x,x)=(1-\varepsilon)p_2(w)+\varepsilon=1-\delta
\)，故可行。
由于
\(
I_{P_\delta}(X;Y)=D_{\mathrm{KL}}(P_\delta\|w\otimes w)
\)
且 \(D_{\mathrm{KL}}(\cdot\|w\otimes w)\) 对第一变量凸，有
\[
D_{\mathrm{KL}}\!\bigl((1-\varepsilon)(w\otimes w)+\varepsilon\,\mathrm{Diag}(w)\,\big\|\,w\otimes w\bigr)
\le \varepsilon\,D_{\mathrm{KL}}\!\bigl(\mathrm{Diag}(w)\,\|\,w\otimes w\bigr),
\]
其中 \(\mathrm{Diag}(w)(x,y):=w(x)\mathbf 1_{\{x=y\}}\)。
而
\(
D_{\mathrm{KL}}(\mathrm{Diag}(w)\|w\otimes w)=\sum_x w(x)\log\frac{w(x)}{w(x)^2}=H(w)
\)，
从而 \(I_{P_\delta}(X;Y)\le \varepsilon H(w)\)。
最后一句由命题 \ref{prop:pom-diagonal-rate-zero-threshold-collision} 立刻推出。证毕。
\end{proof}

\begin{theorem}[线性混合耦合的最优性判别]\label{thm:pom-diagonal-rate-mixture-optimality-criterion}
令 \(X\) 为有限集（\(|X|\ge 2\)），取全支撑分布 \(w\in\Delta(X)\)，并令 \(\delta\in(0,\delta_0)\)（\(\delta_0:=1-p_2(w)\)）。
令 \(P_\delta\) 为定理 \ref{thm:pom-diagonal-rate-bernoulli-lower-and-mixture-upper} 的线性混合耦合。
则 \(P_\delta\) 为 \(R_w(\delta)\) 的最优耦合当且仅当满足以下两者之一：
\[
|X|=2,\qquad\text{或}\qquad w\ \text{为均匀分布}.
\]
在其余情形（即 \(|X|\ge 3\) 且 \(w\) 非均匀）恒有严格不等式 \(I_{P_\delta}(X;Y)>R_w(\delta)\)。
\leanverified{paper\_pom\_diagonal\_rate\_mixture\_optimality\_criterion}
\end{theorem}
\begin{proof}
由定理 \ref{thm:pom-diagonal-rate-explicit-param-and-derivative}，当 \(\delta\in(0,\delta_0)\) 时最优耦合唯一且必具有结构
\[
P^\star_\delta(x,y)=\frac{u_xu_y\bigl(1+\kappa\,\mathbf 1_{\{x=y\}}\bigr)}{Z}
\qquad(\kappa>0,\ u\in\RR_{>0}^X).
\]
若 \(P_\delta\) 为最优，则由唯一性必有 \(P_\delta=P^\star_\delta\)。
特别地，对任意 \(x\ne y\) 有 \(P_\delta(x,y)=(1-\varepsilon)w(x)w(y)\)，而结构式给出 \(P^\star_\delta(x,y)=u_xu_y/Z\)，从而
\[
u_xu_y\ \propto\ w(x)w(y)\qquad(x\ne y).
\]
当 \(|X|\ge 3\) 时，取互异 \(x,y,z\) 得
\(
u_x^2\propto w(x)^2
\)，故存在常数 \(c>0\) 使 \(u_x=c\,w(x)\) 对全部 \(x\) 成立。
代回对角项，结构式要求
\[
\frac{P_\delta(x,x)}{w(x)^2}=\frac{(1-\varepsilon)w(x)^2+\varepsilon w(x)}{w(x)^2}
\]
对 \(x\) 不变；这等价于 \(w(x)\) 为常数，故 \(w\) 必为均匀分布。
反之，当 \(w\) 为均匀分布时，定理 \ref{thm:pom-diagonal-rate-uniform-closed-form} 给出的唯一最优耦合恰为该线性混合形式，因而 \(P_\delta\) 最优。
\par
当 \(|X|=2\) 时由定理 \ref{thm:pom-diagonal-rate-binary-closed-form} 与推论 \ref{cor:pom-diagonal-rate-binary-linear-mixture} 可知线性混合耦合与最优族重合。
综上得到充要条件；其余情形由唯一性立刻推出严格次优。证毕。
\end{proof}

\input{sections/body/pom/parts/part__pom-diagonal-rate-two-orbit-coarsegrain-tensor}

\endinput


% (User input window)

\paragraph{端点展开、曲率审计与 Hellinger 普适极限}
沿用定理 \ref{thm:pom-diagonal-rate-explicit-param-and-derivative} 的最优耦合结构与对偶敏感性公式。
本段记录速率函数 \(R_w(\delta)\) 在两端点 \(\delta\downarrow 0\) 与 \(\delta\uparrow (1-p_2(w))\) 的显式渐近律，并给出由有限碰撞矩可审计的端点曲率常数与稀有错配下的 Hellinger 极限形状。

\begin{theorem}[小失真端点律：\texorpdfstring{$\delta\downarrow0$}{delta->0} 时对偶变量对数发散与速率端点展开]\label{thm:pom-diagonal-rate-small-distortion-hellinger-endpoint}
令有限字母表 \(X\)（\(|X|\ge2\)），取全支撑分布 \(w\in\Delta(X)\)。
记
\[
R_w(\delta):=\inf\{I_P(X;Y):P_X=P_Y=w,\ \mathbb P(X=Y)\ge 1-\delta\},\qquad \delta\in[0,1].
\]
定义 Hellinger 量
\[
A_{1/2}(w):=\sum_{x\in X}\sqrt{w(x)}\quad(>1).
\]
令 \(\lambda(\delta)\) 为最优对偶变量（满足 \(\frac{\mathrm d}{\mathrm d\delta}R_w(\delta)=-\lambda(\delta)\)，参见定理 \ref{thm:pom-diagonal-rate-explicit-param-and-derivative}）。
则当 \(\delta\downarrow 0\) 时成立
\[
\lambda(\delta)=\log\frac{A_{1/2}(w)^2-1}{\delta}+o(1),
\]
并且
\[
R_w(\delta)
=H(w)-\delta\log\frac{1}{\delta}-\delta\log\bigl(A_{1/2}(w)^2-1\bigr)+O(\delta).
\]
\leanverified{paper\_pom\_diagonal\_rate\_small\_distortion\_hellinger\_endpoint}
\end{theorem}
\begin{proof}
由定理 \ref{thm:pom-diagonal-rate-explicit-param-and-derivative}，最优耦合满足
\[
P^\star_\delta(x,y)=\frac{u_xu_y\bigl(1+\kappa\,\mathbf 1_{\{x=y\}}\bigr)}{Z},
\qquad \lambda=\log(1+\kappa).
\]
在 \(\delta\downarrow 0\) 时，对角约束逼近确定耦合，故必有 \(\kappa\to\infty\)。
利用尺度不变性，可取一族等价规范使得
\(
\kappa\sum_x u_x^2=1+o(1)
\)（从而 \(Z=(\sum_x u_x)^2+\kappa\sum_x u_x^2=1+o(1)\)）。
在该规范下，边缘约束方程给出
\[
u_x=\frac{1}{\sqrt{\kappa}}\sqrt{w(x)}+o(\kappa^{-1/2}),
\qquad
\sum_x u_x=\frac{A_{1/2}(w)}{\sqrt{\kappa}}+o(\kappa^{-1/2}),
\qquad
\sum_x u_x^2=\frac{1}{\kappa}+o(\kappa^{-1}).
\]
于是对角质量
\[
\mathbb P(X=Y)
=\frac{(1+\kappa)\sum_x u_x^2}{(\sum_x u_x)^2+\kappa\sum_x u_x^2}
=1-\frac{A_{1/2}(w)^2-1}{\kappa}+o(\kappa^{-1}),
\]
从而
\(
\kappa=\frac{A_{1/2}(w)^2-1}{\delta}\,(1+o(1))
\)
并给出
\(
\lambda(\delta)=\log(1+\kappa)=\log\frac{A_{1/2}(w)^2-1}{\delta}+o(1)
\)。
由 \(R'_w(\delta)=-\lambda(\delta)\) 在 \(\delta\downarrow 0\) 的积分并用 \(R_w(0)=H(w)\)（见正文中的端点说明）即得所示展开。证毕。
\end{proof}

\begin{proposition}[Rényi-\texorpdfstring{$\tfrac12$}{1/2} 熵与 Hellinger 量：严格恒等式与张量可加性]\label{prop:pom-renyi-half-hellinger-identity-tensor-additivity}
令有限字母表 \(X\)（\(|X|\ge 2\)），取全支撑分布 \(w\in\Delta(X)\)，并沿用
\[
A_{1/2}(w):=\sum_{x\in X}\sqrt{w(x)}.
\]
定义 Rényi-\(\tfrac12\) 熵（取自然对数）为
\[
H_{1/2}(w):=\frac{1}{1-\tfrac12}\log\sum_{x\in X}w(x)^{1/2}.
\]
则有严格恒等式
\[
H_{1/2}(w)=2\log A_{1/2}(w).
\]
并且对任意有限字母表 \(Y\) 与全支撑分布 \(v\in\Delta(Y)\) 成立张量可加性
\[
H_{1/2}(w\otimes v)=H_{1/2}(w)+H_{1/2}(v).
\]
\leanverified{paper\_pom\_renyi\_half\_hellinger\_identity\_tensor\_additivity}
\end{proposition}
\begin{proof}
由定义
\(
H_{1/2}(w)=2\log\sum_x w(x)^{1/2}=2\log A_{1/2}(w)
\)。
另一方面，
\[
A_{1/2}(w\otimes v)
=\sum_{(x,y)\in X\times Y}\sqrt{w(x)v(y)}
=\Bigl(\sum_{x\in X}\sqrt{w(x)}\Bigr)\Bigl(\sum_{y\in Y}\sqrt{v(y)}\Bigr)
=A_{1/2}(w)\,A_{1/2}(v),
\]
代回 \(H_{1/2}=2\log A_{1/2}\) 得张量可加性。证毕。
\end{proof}

\begin{corollary}[小失真端点常数的严格乘法律：由 \texorpdfstring{$H_{1/2}$}{H1/2} 生成]\label{cor:pom-small-distortion-constant-tensor-multiplicative}
在定理 \ref{thm:pom-diagonal-rate-small-distortion-hellinger-endpoint} 的设定下，定义端点常数
\[
C_{1/2}(w):=\lim_{\delta\downarrow 0}\ \delta\,e^{\lambda(\delta)}
=A_{1/2}(w)^2-1
=e^{H_{1/2}(w)}-1\ >0.
\]
则对任意两全支撑分布 \(w,v\) 有严格恒等式
\[
C_{1/2}(w\otimes v)=C_{1/2}(w)+C_{1/2}(v)+C_{1/2}(w)C_{1/2}(v),
\qquad
1+C_{1/2}(w\otimes v)=(1+C_{1/2}(w))(1+C_{1/2}(v)).
\]
并且令小失真节省量 \(\Delta_w(\delta):=H(w)-R_w(\delta)\)，则当 \(\delta\downarrow 0\) 时
\[
\Delta_w(\delta)=\delta\log\frac{C_{1/2}(w)}{\delta}+O(\delta).
\]
\end{corollary}
\leanverified{paper\_pom\_small\_distortion\_constant\_tensor\_multiplicative}
\begin{proof}
由定理 \ref{thm:pom-diagonal-rate-small-distortion-hellinger-endpoint} 的对偶展开 \(\lambda(\delta)=\log\frac{A_{1/2}(w)^2-1}{\delta}+o(1)\)，取指数并乘以 \(\delta\) 得
\(
\delta e^{\lambda(\delta)}\to A_{1/2}(w)^2-1
\)，从而 \(C_{1/2}(w)=A_{1/2}(w)^2-1\)。
再由命题 \ref{prop:pom-renyi-half-hellinger-identity-tensor-additivity} 的乘法恒等式 \(A_{1/2}(w\otimes v)=A_{1/2}(w)A_{1/2}(v)\)，得
\[
1+C_{1/2}(w\otimes v)=A_{1/2}(w\otimes v)^2=A_{1/2}(w)^2A_{1/2}(v)^2=(1+C_{1/2}(w))(1+C_{1/2}(v)),
\]
展开即为所示乘法律。
最后一式由定理 \ref{thm:pom-diagonal-rate-small-distortion-hellinger-endpoint} 取差 \(\Delta_w=H-R_w\) 直接得到。证毕。
\end{proof}

\begin{theorem}[小失真协同常数的唯一张量生成元]\label{thm:pom-small-distortion-synergy-unique-additive-generator}
\leanverified{paper\_pom\_small\_distortion\_synergy\_unique\_additive\_generator}
设 \(F\) 为定义在全体有限字母表的全支撑分布上的函数族，满足：
\begin{enumerate}
  \item \textup{（置换不变性）}对任意有限集 \(X\) 与任意置换 \(\pi:X\to X\)，有 \(F(w\circ\pi^{-1})=F(w)\)；
  \item \textup{（张量可加性）}对任意有限集 \(X,Y\) 与全支撑 \(w\in\Delta(X),v\in\Delta(Y)\)，有 \(F(w\otimes v)=F(w)+F(v)\)；
  \item \textup{（连续性）}在每个 \(\Delta(X)\) 的内点上连续。
\end{enumerate}
若存在连续函数 \(g:\mathrm{Im}(F)\to\RR_{>0}\) 使对一切全支撑 \(w\) 成立
\[
C_{1/2}(w)=g(F(w)),
\]
其中 \(C_{1/2}\) 为推论 \ref{cor:pom-small-distortion-constant-tensor-multiplicative} 的端点常数，
则存在常数 \(c>0\) 使得对所有 \(w\) 有
\[
F(w)=\frac{1}{c}\,H_{1/2}(w),
\qquad
g(z)=e^{cz}-1.
\]
\end{theorem}
\begin{proof}
由张量可加性与推论 \ref{cor:pom-small-distortion-constant-tensor-multiplicative} 的乘法律，对任意全支撑 \(w,v\) 有
\[
1+g\!\bigl(F(w)+F(v)\bigr)
=1+g\!\bigl(F(w\otimes v)\bigr)
=1+C_{1/2}(w\otimes v)
=(1+C_{1/2}(w))(1+C_{1/2}(v))
=\bigl(1+g(F(w))\bigr)\bigl(1+g(F(v))\bigr).
\]
令 \(h(z):=\log\bigl(1+g(z)\bigr)\)（其在 \(\mathrm{Im}(F)\) 上有意义且连续），则 \(h\) 在加法半群 \(\mathrm{Im}(F)\) 上满足 Cauchy 型方程
\[
h(z_1+z_2)=h(z_1)+h(z_2).
\]
连续性迫使 \(h\) 为线性：存在常数 \(c>0\) 使 \(h(z)=cz\)，从而 \(g(z)=e^{cz}-1\)。
另一方面，由 \(C_{1/2}(w)=g(F(w))\) 得
\[
H_{1/2}(w)=\log\bigl(1+C_{1/2}(w)\bigr)=\log\bigl(1+g(F(w))\bigr)=h(F(w))=cF(w),
\]
即 \(F(w)=H_{1/2}(w)/c\)。证毕。
\end{proof}

\begin{corollary}[临界比特密度在小失真端的局部展开]\label{cor:pom-bcrit-small-distortion-expansion}
沿用临界比特密度的定义
\(
b_{\mathrm{crit}}(\beta,\delta)=(1-\beta)R_w(\delta)/\log 2
\)。
在定理 \ref{thm:pom-diagonal-rate-small-distortion-hellinger-endpoint} 的设定下，当 \(\delta\downarrow 0\) 时
\[
b_{\mathrm{crit}}(\beta,\delta)
=\frac{1-\beta}{\log 2}\Bigl(H(w)-\delta\log\frac{C_{1/2}(w)}{\delta}+O(\delta)\Bigr),
\]
其中端点常数 \(C_{1/2}(w)=e^{H_{1/2}(w)}-1\)。
因此在 \(\delta=0\) 的局部尺度上，截距由 \(H(w)\) 给出，而首阶协同项由 \(H_{1/2}(w)\)（等价地 \(C_{1/2}(w)\)）唯一生成。
\end{corollary}
\begin{proof}
由定义 \(b_{\mathrm{crit}}(\beta,\delta)=(1-\beta)R_w(\delta)/\log 2\)，直接代入定理 \ref{thm:pom-diagonal-rate-small-distortion-hellinger-endpoint} 的端点展开并用 \(C_{1/2}(w)=A_{1/2}(w)^2-1\) 即得。证毕。
\end{proof}

\input{sections/body/pom/parts/part__pom-diagonal-rate-half-moment-critical-slope}

\begin{theorem}[缩放势的两端点渐近]\label{thm:pom-diagonal-rate-scaling-potential-endpoints}
令有限字母表 \(X\)（\(|X|\ge 2\)），取全支撑分布 \(w\in\Delta(X)\)。
对 \(\delta\in(0,\delta_0)\)（\(\delta_0:=1-p_2(w)\)），令 \(\psi_\delta\) 为定理 \ref{thm:pom-diagonal-rate-frechet-derivative-scaling-potential} 的缩放势。
则当 \(\delta\downarrow 0\) 时存在标量 \(c(\delta)\) 使
\[
\psi_\delta(x)= -\frac12\log w(x)+c(\delta)+o(1),\qquad x\in X,
\]
并且对任意 \(x,y\in X\) 有
\[
\psi_\delta(x)-\psi_\delta(y)\ \longrightarrow\ -\frac12\log\frac{w(x)}{w(y)}.
\]
另一方面，令 \(\varepsilon:=\delta_0-\delta\downarrow 0^+\)，并记
\[
C(w):=p_2(w)+p_2(w)^2-2p_3(w),\qquad p_3(w):=\sum_{x\in X}w(x)^3.
\]
则存在标量 \(c_\varepsilon\) 使
\[
\psi_\delta(x)=\frac{\varepsilon}{C(w)}\Bigl(\frac{p_2(w)}{2}-w(x)\Bigr)+c_\varepsilon+O(\varepsilon^2),
\qquad x\in X.
\]
\end{theorem}
\begin{proof}
当 \(\delta\downarrow 0\) 时，由定理 \ref{thm:pom-diagonal-rate-small-distortion-hellinger-endpoint} 的证明可取规范使 \(u_x=\kappa^{-1/2}\sqrt{w(x)}+o(\kappa^{-1/2})\)，从而
\[
\psi_\delta(x)=\log\frac{u_x}{w(x)}=-\frac12\log w(x)-\frac12\log\kappa+o(1).
\]
取 \(c(\delta):=-(1/2)\log\kappa\) 即得第一式；差分极限由常数项抵消直接推出。
\par
当 \(\delta\uparrow\delta_0\) 时，定理 \ref{thm:pom-diagonal-rate-zero-rate-endpoint-quadratic-curvature} 的证明在规范 \(\sum_x u_x=1\) 下给出展开
\(
u_x=w(x)+\kappa w(x)\bigl(p_2(w)-w(x)\bigr)+O(\kappa^2)
\)
以及
\(
\varepsilon=\delta_0-\delta=\kappa C(w)+O(\kappa^2)
\)。
因此
\[
\psi_\delta(x)=\log\Bigl(1+\kappa\bigl(p_2(w)-w(x)\bigr)+O(\kappa^2)\Bigr)
=\kappa\bigl(p_2(w)-w(x)\bigr)+O(\kappa^2).
\]
消去 \(\kappa=\varepsilon/C(w)+O(\varepsilon^2)\) 得
\(
\psi_\delta(x)=\frac{\varepsilon}{C(w)}\bigl(p_2(w)-w(x)\bigr)+O(\varepsilon^2)
\)。
最后将常数项 \(\frac{\varepsilon}{C(w)}\frac{p_2(w)}{2}\) 吸收进 \(c_\varepsilon\) 即得所示形式。证毕。
\end{proof}
\leanverified{paper\_pom\_diagonal\_rate\_scaling\_potential\_endpoints}

\leanverified{paper\_pom\_diagonal\_rate\_zero\_rate\_endpoint\_quadratic\_curvature}
\begin{theorem}[零速率端点的二次起跳：\texorpdfstring{$\delta\uparrow(1-p_2(w))$}{delta->delta0} 时曲率仅依赖 \texorpdfstring{$p_2,p_3$}{p2,p3}]\label{thm:pom-diagonal-rate-zero-rate-endpoint-quadratic-curvature}
记幂和 \(p_k(w):=\sum_{x\in X}w(x)^k\)，并令
\[
\delta_0:=1-p_2(w).
\]
则 \(\delta\ge\delta_0\) 时独立耦合 \(w\otimes w\) 可行且给出 \(R_w(\delta)=0\)（命题 \ref{prop:pom-diagonal-rate-zero-threshold-collision}）。
当 \(\delta\uparrow\delta_0\) 时，
\[
R_w(\delta)=\frac{(\delta_0-\delta)^2}{2\,C(w)}+O\bigl((\delta_0-\delta)^3\bigr),
\]
其中曲率常数
\[
C(w):=p_2(w)+p_2(w)^2-2p_3(w)
=\sum_{x\in X}w(x)^2(1-w(x))^2+2\sum_{x<y}w(x)^2w(y)^2>0.
\]
\end{theorem}
\begin{proof}
在 \(\delta\uparrow\delta_0\) 时，最优解逼近独立耦合，对偶参数 \(\kappa(\delta)=e^{\lambda(\delta)}-1\downarrow 0\)（定理 \ref{thm:pom-diagonal-rate-explicit-param-and-derivative}）。
利用尺度不变性取规范 \(\sum_x u_x=1\)。
写
\(
u_x=w(x)+\kappa a_x+O(\kappa^2)
\)。
由边缘约束
\(
w(x)=u_x(1+\kappa u_x)/(1+\kappa\sum_a u_a^2)
\)
的一阶比较可得
\(
a_x=w(x)\bigl(p_2(w)-w(x)\bigr)
\)，并因此
\(
\sum_x u_x^2=p_2(w)+2\kappa\bigl(p_2(w)^2-p_3(w)\bigr)+O(\kappa^2)
\)。
对角质量为
\[
\mathbb P(X=Y)=\frac{(1+\kappa)\sum_x u_x^2}{1+\kappa\sum_x u_x^2}
=p_2(w)+\kappa\,C(w)+O(\kappa^2),
\]
从而
\(
\delta_0-\delta=\mathbb P(X=Y)-p_2(w)=\kappa C(w)+O(\kappa^2)
\)。
另一方面，由敏感性公式
\(
R'_w(\delta)=-\log(1+\kappa)=-\kappa+O(\kappa^2)
\)。
消去 \(\kappa\) 得
\(
R'_w(\delta)=-(\delta_0-\delta)/C(w)+O((\delta_0-\delta)^2)
\)，对 \(\delta\) 积分即得二次起跳展开。证毕。
\end{proof}

\begin{proposition}[解析凸性与曲率表示：\texorpdfstring{$R_w$}{R} 在 \texorpdfstring{$(0,\delta_0)$}{(0,delta0)} 上严格凸]\label{prop:pom-diagonal-rate-analytic-strict-convexity-curvature}
令 \(\delta\in(0,\delta_0)\)，记最优对偶参数 \(\kappa(\delta)=e^{\lambda(\delta)}-1\in(0,\infty)\)。
则 \(R_w(\delta)\) 在 \((0,\delta_0)\) 上为实解析且严格凸，并满足端点斜率
\[
R_w'(0^+)=-\infty,\qquad R_w'(\delta_0^-)=0.
\]
并且对角指示变量
\[
Z:=\mathbf 1_{\{X=Y\}}\qquad\bigl((X,Y)\sim P^\star_\delta\bigr)
\]
满足 \(\EE[Z]=1-\delta\)，且二阶导数具有涨落—耗散形式
\[
R_w''(\delta)=-\frac{\kappa'(\delta)}{1+\kappa(\delta)}>0,
\qquad
R_w''(\delta)=\left(\frac{\mathrm d}{\mathrm d\lambda}\mathbb P_{P^\star_\delta}(X=Y)\right)^{-1}
=\frac{1}{\Var_{P^\star_\delta}(Z)}.
\]
\end{proposition}
\begin{proof}
由定理 \ref{thm:pom-diagonal-rate-explicit-param-and-derivative}，在 \((0,\delta_0)\) 上最优解唯一且随参数光滑变化；\(\kappa\mapsto \mathbb P(X=Y)\) 为严格增的实解析函数，故由隐函数定理得 \(\kappa(\delta)\) 在 \((0,\delta_0)\) 上实解析且严格递减。
由
\(
R_w'(\delta)=-\log(1+\kappa(\delta))
\)
得
\(
R_w''(\delta)=-\kappa'(\delta)/(1+\kappa(\delta))>0
\)，从而严格凸。
端点斜率由 \(\delta\downarrow 0\Rightarrow \kappa\to\infty\) 与 \(\delta\uparrow\delta_0\Rightarrow \kappa\to 0\) 直接给出。
由 \(\lambda=\log(1+\kappa)\) 与约束 \(\EE_{P^\star_\delta}[Z]=1-\delta\) 的链式求导得
\[
\frac{\mathrm d}{\mathrm d\delta}\lambda(\delta)
=-\left(\frac{\mathrm d}{\mathrm d\lambda}\EE_{P^\star_\delta}[Z]\right)^{-1}.
\]
而在指数族 \(P^\star_\delta\propto u_xu_y e^{\lambda Z}\) 下，配分函数的二阶导数恒等式给出
\(
\frac{\mathrm d}{\mathrm d\lambda}\EE[Z]=\Var(Z)
\)。
结合 \(R_w'(\delta)=-\lambda(\delta)\)（定理 \ref{thm:pom-diagonal-rate-explicit-param-and-derivative}）即得
\(
R_w''(\delta)=-\lambda'(\delta)=1/\Var(Z)
\)，从而严格凸。证毕。
\end{proof}

\begin{proposition}[端点的有限矩闭包：\texorpdfstring{$\delta_0$}{delta0} 邻域内 Taylor 系数仅依赖有限阶碰撞矩]\label{prop:pom-diagonal-rate-finite-moment-closure}
记 \(\varepsilon:=\delta_0-\delta\ge 0\)。
在 \(\varepsilon\) 足够小的邻域内，\(R_w(\delta)\) 可展开为
\[
R_w(\delta)=\sum_{j\ge 2} a_j(w)\,\varepsilon^j,
\qquad
a_2(w)=\frac{1}{2C(w)}.
\]
并且对任意整数 \(j\ge 2\)，系数 \(a_j(w)\) 可表示为 \(w\) 的对称多项式，且可由有限幂和 \(\{p_2(w),p_3(w),\dots,p_{j+1}(w)\}\) 的有理函数唯一确定。
\end{proposition}
\begin{proof}
在 \(\kappa\to 0\) 的解析区间内，最优缩放向量 \(u_x\) 可作 \(\kappa\) 的幂级数展开，其每一阶系数为 \(\{w(x)\}\) 的对称多项式；
因此 \(\varepsilon(\kappa)\) 与 \(R_w(\kappa)\) 均可写为 \(\kappa\) 的幂级数，其第 \(j\) 阶系数仅涉及至多 \(w(x)^{j+1}\) 的加权求和。
由 Newton--Girard 恒等式，任意此类对称多项式可用有限幂和 \(p_2,\dots,p_{j+1}\) 表达；
对 \(\varepsilon(\kappa)\) 在零点附近解析可逆后消去 \(\kappa\) 即得所示闭包与唯一性。证毕。
\end{proof}

\begin{corollary}[折叠可审计曲率：对任意固定 \texorpdfstring{$m$}{m}，二次起跳由 \texorpdfstring{$(\mathrm{MOM}_2,\mathrm{MOM}_3)$}{(MOM2,MOM3)} 完全决定]\label{cor:pom-diagonal-rate-folded-curvature-auditable}
在折叠语义下令推前分布 \(w_m(x)=d_m(x)/2^m\)，并记碰撞矩
\[
S_q(m):=\sum_{x\in X_m}d_m(x)^q.
\]
则
\[
p_2(w_m)=\sum_x w_m(x)^2=2^{-2m}S_2(m),\qquad
p_3(w_m)=\sum_x w_m(x)^3=2^{-3m}S_3(m),
\qquad
\delta_0(m):=1-p_2(w_m).
\]
记
\[
C_m:=p_2(w_m)+p_2(w_m)^2-2p_3(w_m).
\]
则当 \(\delta\uparrow \delta_0(m)\) 时成立
\[
R_{w_m}(\delta)=\frac{(\delta_0(m)-\delta)^2}{2C_m}+O\bigl((\delta_0(m)-\delta)^3\bigr).
\]
特别地，曲率常数 \(C_m\) 仅依赖于 \(\{S_2(m),S_3(m)\}\)，因而在本文的有限核编译口径下可由 \(\mathrm{MOM}_2,\mathrm{MOM}_3\) 的有限状态计数直接审计，而无需重建全体 \(w_m\)。
\end{corollary}
\begin{proof}
对固定 \(m\) 将定理 \ref{thm:pom-diagonal-rate-zero-rate-endpoint-quadratic-curvature} 应用于 \(w=w_m\)，并代入 \(p_k(w_m)=2^{-mk}S_k(m)\) 即得。证毕。
\end{proof}

\input{sections/body/pom/parts/part__pom-diagonal-rate-endpoint-three-kernel-half-moment-nonreplacement}

\begin{theorem}[稀有错配的 Hellinger 普适性：\texorpdfstring{$\delta\downarrow0$}{delta->0} 时最优耦合的非对角极限形状]\label{thm:pom-diagonal-rate-rare-mismatch-hellinger-universality}
令 \(P^\star_\delta\) 为实现 \(R_w(\delta)\) 的唯一最优耦合，并记 \(A_{1/2}(w)=\sum_x\sqrt{w(x)}\)。
当 \(\delta\downarrow 0\) 时，对任意 \(x\ne y\) 有
\[
\frac{1}{\delta}\,P^\star_\delta(x,y)\ \longrightarrow\ \frac{\sqrt{w(x)w(y)}}{A_{1/2}(w)^2-1},
\]
并且对角项满足
\[
P^\star_\delta(x,x)
=w(x)-\delta\,\frac{\sqrt{w(x)}\bigl(A_{1/2}(w)-\sqrt{w(x)}\bigr)}{A_{1/2}(w)^2-1}+o(\delta).
\]
令背景分布
\(
\pi_\delta(x):=u_x/\sum_y u_y
\)
（\(u\) 为定理 \ref{thm:pom-diagonal-rate-explicit-param-and-derivative} 的最优缩放向量），则
\[
\pi_\delta(x)\ \longrightarrow\ \pi_{1/2}(x):=\frac{\sqrt{w(x)}}{A_{1/2}(w)}.
\]
并且符号依赖保持率 \(p_x(\delta):=\mathbb P(Y=x\mid X=x)\) 满足
\[
p_x(\delta)
=1-\delta\,\frac{A_{1/2}(w)-\sqrt{w(x)}}{\sqrt{w(x)}\,(A_{1/2}(w)^2-1)}+o(\delta).
\]
\end{theorem}
\begin{proof}
由定理 \ref{thm:pom-diagonal-rate-explicit-param-and-derivative} 的指数族表示，\(P^\star_\delta(x,y)=u_xu_y(1+\kappa\mathbf 1_{\{x=y\}})/Z\)。
在 \(\delta\downarrow 0\) 时 \(\kappa\to\infty\)，并由定理 \ref{thm:pom-diagonal-rate-small-distortion-hellinger-endpoint} 的证明可取同一规范使
\(
u_x=\kappa^{-1/2}\sqrt{w(x)}+o(\kappa^{-1/2})
\)
且
\(
\kappa\sim (A_{1/2}(w)^2-1)/\delta
\)，同时 \(Z=1+o(1)\)。
于是对 \(x\ne y\),
\[
P^\star_\delta(x,y)=\frac{u_xu_y}{Z}
=\frac{1}{\kappa}\sqrt{w(x)w(y)}\,(1+o(1))
=\delta\,\frac{\sqrt{w(x)w(y)}}{A_{1/2}(w)^2-1}\,(1+o(1)),
\]
得第一式。
对角项由边缘约束
\(
w(x)=\sum_y P^\star_\delta(x,y)
\)
与非对角极限求和得到二式。
背景分布极限由 \(u_x\propto \sqrt{w(x)}\) 直接给出。
最后，由
\(
p_x(\delta)=P^\star_\delta(x,x)/w(x)
\)
代入对角展开即得保持率展开。证毕。
\end{proof}

\input{sections/body/pom/parts/part__pom-diagonal-rate-keep-resample-channel}

\input{sections/body/pom/parts/part__pom-diagonal-rate-maxent-markov-spectrum-synergy}

\input{sections/body/pom/parts/part__pom-diagonal-rate-absorbing-laguerre-hitting-time}

\begin{proposition}[最优缩放与保持率的同序性]\label{prop:pom-diagonal-rate-scaling-keep-comonotone}
\leanverified{paper\_pom\_diagonal\_rate\_scaling\_keep\_comonotone}
取全支撑分布 \(w\in\Delta(X)\) 与 \(\delta\in(0,\delta_0)\)（\(\delta_0:=1-p_2(w)\)）。
令 \(P^\star_\delta\) 为实现 \(R_w(\delta)\) 的唯一最优耦合，并写成定理 \ref{thm:pom-diagonal-rate-explicit-param-and-derivative} 的形式
\[
P^\star_\delta(x,y)=\frac{u_xu_y\bigl(1+\kappa\,\mathbf 1_{\{x=y\}}\bigr)}{Z},
\qquad
\kappa>0,\ u\in\RR_{>0}^X.
\]
定义条件核 \(K_\delta(y\mid x):=P^\star_\delta(x,y)/w(x)\)、保持率
\[
p_\delta(x):=K_\delta(x\mid x)=\mathbb P(Y=x\mid X=x),
\]
以及背景分布（缩放归一化）
\[
\pi_\delta(x):=\frac{u_x}{\sum_{z\in X}u_z}.
\]
若 \(w(x)\ge w(y)\)，则必有
\[
u_x\ge u_y,\qquad \pi_\delta(x)\ge \pi_\delta(y),\qquad p_\delta(x)\ge p_\delta(y).
\]
并且当 \(w\) 非均匀时，上述不等式对至少一对 \((x,y)\) 为严格。
\end{proposition}
\begin{proof}
令 \(A:=\sum_z u_z\)。
由边缘自洽关系 \(w(x)=u_x(A+\kappa u_x)/Z\) 得到对每个 \(x\) 的二次方程
\[
\kappa u_x^2+A u_x=Zw(x).
\]
右端仅随 \(w(x)\) 变化而系数 \((\kappa,A,Z)\) 对所有 \(x\) 相同，故其正根 \(u_x=f(w(x))\)（\(f\) 严格递增）直接给出 \(w(x)\ge w(y)\Rightarrow u_x\ge u_y\)；
归一化得到 \(\pi_\delta\) 同序。
另一方面，
\[
p_\delta(x)
=\frac{P^\star_\delta(x,x)}{w(x)}
=\frac{(1+\kappa)u_x^2/Z}{u_x(A+\kappa u_x)/Z}
=\frac{(1+\kappa)u_x}{A+\kappa u_x},
\]
其对 \(u_x\) 严格递增，因此 \(u_x\ge u_y\Rightarrow p_\delta(x)\ge p_\delta(y)\)。
若 \(w\) 非均匀，则存在 \(x,y\) 使 \(w(x)>w(y)\)，由 \(f\) 严格递增即得严格性。证毕。
\end{proof}

\begin{theorem}[缩放分布的主序插值：\texorpdfstring{$\pi_\delta$}{pi} 从 \texorpdfstring{$\pi_{1/2}$}{pi1/2} 到 \(w\) 的单调链]\label{thm:pom-diagonal-rate-pi-majorization-chain}
\leanverified{paper\_pom\_diagonal\_rate\_pi\_majorization\_chain}
令 \(w\in\Delta(X)\) 全支撑并记 \(\delta_0:=1-p_2(w)\)。
对 \(\delta\in(0,\delta_0)\) 令 \(\pi_\delta\) 为命题 \ref{prop:pom-diagonal-rate-scaling-keep-comonotone} 的背景分布，并定义
\[
\pi_{1/2}(x):=\frac{\sqrt{w(x)}}{\sum_{z\in X}\sqrt{w(z)}}.
\]
则当 \(\delta\downarrow 0\) 时有 \(\pi_\delta\to \pi_{1/2}\)，当 \(\delta\uparrow\delta_0\) 时有 \(\pi_\delta\to w\)。
并且对每个 \(\delta\in(0,\delta_0)\) 成立主序夹逼
\[
\pi_{1/2}\ \prec\ \pi_\delta\ \prec\ w,
\]
其中 \(\prec\) 为 Hardy--Littlewood--P\'olya 主序；若 \(w\) 非均匀则上述不等式为严格。
进一步地，若 \(0<\delta_1<\delta_2<\delta_0\)，则有单调主序链
\[
\pi_{\delta_1}\ \prec\ \pi_{\delta_2}.
\]
\end{theorem}
\begin{proof}
端点极限由定理 \ref{thm:pom-diagonal-rate-rare-mismatch-hellinger-universality}（\(\delta\downarrow 0\)）与
定理 \ref{thm:pom-diagonal-rate-zero-rate-endpoint-quadratic-curvature} 的 \(\kappa\downarrow 0\) 展开（在规约 \(\sum_x u_x=1\) 下 \(u\to w\)，从而 \(\pi_\delta\to w\)）给出。
\par
为证明主序关系，按 \(w\) 的非增序排列字母表，使得 \(w_{(1)}\ge\cdots\ge w_{(n)}\)，并相应排列 \(\pi_\delta\)。
由定理 \ref{thm:pom-diagonal-rate-explicit-param-and-derivative} 的二次关系可写成
\(
\kappa u_x^2+A u_x=Zw(x)
\)，
从而
\(
u_x/w(x)=\frac{2Z}{A+\sqrt{A^2+4\kappa Zw(x)}}
\)
随 \(w(x)\) 严格递减。
因此序列 \(\bigl(\pi_{\delta,(i)}/w_{(i)}\bigr)_{i=1}^n\) 严格递增（常数 \(\sum_z u_z\) 只作整体缩放）。
由此对任意 \(k\in\{1,\dots,n\}\)，初段加权平均满足
\[
\frac{\sum_{i=1}^k \pi_{\delta,(i)}}{\sum_{i=1}^k w_{(i)}}
=\frac{\sum_{i=1}^k w_{(i)}\cdot(\pi_{\delta,(i)}/w_{(i)})}{\sum_{i=1}^k w_{(i)}}
\le \frac{\sum_{i=1}^n w_{(i)}\cdot(\pi_{\delta,(i)}/w_{(i)})}{\sum_{i=1}^n w_{(i)}}
=1,
\]
其中不等式由“向递增序列加入不小于当前平均的新项会增加平均”的基本事实给出。
于是 \(\sum_{i=1}^k \pi_{\delta,(i)}\le \sum_{i=1}^k w_{(i)}\)（\(k=1,\dots,n-1\)），即 \(\pi_\delta\prec w\)。
若 \(w\) 非均匀，则 \(\pi_{\delta,(i)}/w_{(i)}\) 非常数，故至少对某个 \(k\) 严格不等，得到严格主序。
\par
对单调链，注意 \(\delta\mapsto \kappa(\delta)\) 严格递减（命题 \ref{prop:pom-diagonal-rate-analytic-strict-convexity-curvature}），而对固定 \(w\) 上述比值谱
\(
w\mapsto u_x/w(x)
\)
在 \(\kappa\) 上呈严格单调的扁平化：\(\kappa\) 越大，对大权重坐标的相对压缩越强，从而 \(\pi_\delta\) 在主序意义下越接近均匀端点 \(\pi_{1/2}\)。
据此对 \(0<\delta_1<\delta_2<\delta_0\) 得 \(\pi_{\delta_1}\prec \pi_{\delta_2}\)，并由端点极限推出 \(\pi_{1/2}\prec \pi_\delta\)。
证毕。
\end{proof}

\input{sections/body/pom/parts/part__pom-diagonal-rate-schur-concavity-tp2}
\input{sections/body/pom/parts/part__pom-diagonal-rate-schur-concavity}

\endinput


\endinput


\begin{theorem}[固定 \(t\) 的热力学极限：完全齐次对称多项式]\label{thm:pom-microcanonical-fold-bayes-success-thermo-limit}
令 \(w(x):=d(x)/N\)。
若 \(t\) 固定而 \(N\to\infty\)，则
\[
\mathrm{Succ}^{\mathrm{Bayes}}_{N-t}(d)\ \longrightarrow\ \mathsf{h}_t(w)
:=\sum_{\substack{r_x\in\ZZ_{\ge 0}\\ \sum_x r_x=t}}\ \prod_{x\in X_m} w(x)^{r_x}.
\]
并且 \(\mathsf{h}_t(w)\) 的生成函数为
\[
\boxed{\;
\sum_{t\ge 0}\mathsf{h}_t(w)\,z^t=\prod_{x\in X_m}\frac{1}{1-w(x)z}\; } .
\]
\end{theorem}
\begin{proof}
对固定 \(t\)，有
\(
(d(x))_{r_x}=d(x)^{r_x}+O(d(x)^{r_x-1})
\)
且 \((N)_t=N^t+O(N^{t-1})\)。
将定理 \ref{thm:pom-microcanonical-fold-bayes-success-Nminus-t} 的闭式除以 \(N^t\) 并用 \(d(x)=Nw(x)\) 展开，
得极限为 \(\mathsf{h}_t(w)\)。
\par
生成函数由完全齐次对称多项式的标准定义给出：
对每个 \(x\) 有 \((1-w(x)z)^{-1}=\sum_{r\ge 0}w(x)^r z^r\)，对所有 \(x\) 相乘并取 \(z^t\) 系数即得。证毕。
\end{proof}

\leanverified{paper\_pom\_microcanonical\_fold\_ht\_uniform\_min}
\begin{theorem}[均匀极小原理：任意 \(t\) 下 \texorpdfstring{$\mathsf{h}_t$}{h\_t} 在均匀分布处取最小]\label{thm:pom-microcanonical-fold-ht-uniform-min}
设 \(n=\card{X_m}\ge 1\)，令 \(U\in\Delta(X_m)\) 为均匀分布 \(U(x)=1/n\)。
则对任意 \(w\in\Delta(X_m)\) 与任意整数 \(t\ge 0\)，成立
\[
\boxed{\;
\mathsf{h}_t(w)\ \ge\ \mathsf{h}_t(U)=\frac{\binom{n+t-1}{t}}{n^t}\; } ,
\]
且等号当且仅当 \(w\equiv U\)。
\end{theorem}
\begin{proof}
取任意 \(z\in(0,1)\)，令 \(\phi_z(x):=-\log(1-zx)\)。
则 \(\phi_z\) 在 \(x\in[0,1]\) 上严格凸，且由 Jensen 不等式
\[
\sum_{x\in X_m}\phi_z(w(x))
\ \ge\
n\,\phi_z\!\left(\frac{1}{n}\sum_x w(x)\right)
= -n\log\!\left(1-\frac{z}{n}\right).
\]
指数化得
\(
\prod_x(1-w(x)z)^{-1}\ge (1-z/n)^{-n}
\)。
两边展开为系数非负的幂级数，逐项比较 \(z^t\) 系数即得 \(\mathsf{h}_t(w)\ge \mathsf{h}_t(U)\)；
而
\(
(1-z/n)^{-n}=\sum_{t\ge 0}\binom{n+t-1}{t}(z/n)^t
\)
给出闭式值。
严格凸性推出等号充要条件为 \(w\equiv U\)。证毕。
\end{proof}

\begin{theorem}[碰撞矩完备性：\texorpdfstring{$\mathsf{h}_t(w)$}{h\_t(w)} 由有限幂和显式生成]\label{thm:pom-microcanonical-fold-ht-from-power-sums}
\leanverified{paper\_pom\_microcanonical\_fold\_ht\_from\_power\_sums}
对分布 \(w\in\Delta(X_m)\) 定义幂和
\(
p_k(w):=\sum_{x\in X_m}w(x)^k
\)
（其中 \(p_1(w)=1\)）。
令
\[
\mathcal{H}_w(z):=\sum_{t\ge 0}\mathsf{h}_t(w)\,z^t=\prod_{x\in X_m}(1-w(x)z)^{-1}.
\]
则作为形式幂级数有严格恒等式
\[
\boxed{\;
\log \mathcal{H}_w(z)=\sum_{k\ge 1}\frac{p_k(w)}{k}\,z^k\; } .
\]
从而对任意固定 \(t\)，\(\mathsf{h}_t(w)\) 仅依赖于 \(\{p_2(w),\dots,p_t(w)\}\)，并可由截断指数级数的 \(z^t\) 系数唯一确定。
\end{theorem}
\begin{proof}
由恒等式 \(-\log(1-u)=\sum_{k\ge 1}u^k/k\)（作为形式幂级数）与 \(u=w(x)z\)，有
\[
\log \mathcal{H}_w(z)
=-\sum_{x}\log(1-w(x)z)
=\sum_{x}\sum_{k\ge 1}\frac{w(x)^k}{k}z^k
=\sum_{k\ge 1}\frac{p_k(w)}{k}z^k.
\]
由于 \(p_1(w)=1\) 固定，故 \(z^t\) 系数仅依赖于 \(p_2,\dots,p_t\)。证毕。
\end{proof}

\begin{corollary}[低阶显式公式]\label{cor:pom-microcanonical-fold-ht-low-order}
记 \(p_k:=p_k(w)\)，则
\[
\mathsf{h}_2=\frac{1+p_2}{2},\qquad
\mathsf{h}_3=\frac{1+3p_2+2p_3}{6},
\]
\[
\mathsf{h}_4=\frac{1+6p_2+3p_2^2+8p_3+6p_4}{24},
\]
\[
\mathsf{h}_5=\frac{1+10p_2+15p_2^2+20p_3+20p_2p_3+30p_4+24p_5}{120}.
\]
\end{corollary}
\begin{proof}
由定理 \ref{thm:pom-microcanonical-fold-ht-from-power-sums} 将
\(
\mathcal{H}_w(z)=\exp(\sum_{k\ge 1}p_k z^k/k)
\)
截断到相应阶并直接展开，取系数即可。证毕。
\end{proof}

\begin{theorem}[折叠矩谱到微观可逆性的编译：\texorpdfstring{$(N-t)$}{N-t} 查询 Bayes 识别率为有限多项式]\label{thm:pom-microcanonical-fold-bayes-success-polynomial-in-moments}
在折叠语义下取 \(d=d_m\)、\(N=2^m\)、\(w=w_m\)，并记碰撞矩
\(
S_k(m):=\sum_{x\in X_m}d_m(x)^k
\)。
则对每个固定整数 \(t\ge 0\)，当 \(m\to\infty\) 时
\[
\boxed{\;
\mathrm{Succ}^{\mathrm{Bayes}}_{2^m-t}(d_m)=\mathsf{h}_t(w_m)+o(1)\;}
\]
并且由 \(\,p_k(w_m)=\sum_x w_m(x)^k=2^{-mk}S_k(m)\) 可知 \(\mathsf{h}_t(w_m)\) 可由
\(\{S_2(m),\dots,S_t(m)\}\) 通过推论 \ref{cor:pom-microcanonical-fold-ht-low-order} 的多项式公式显式计算。
\end{theorem}
\begin{proof}
由定理 \ref{thm:pom-microcanonical-fold-bayes-success-thermo-limit} 得
\(
\mathrm{Succ}^{\mathrm{Bayes}}_{N-t}(d)=\mathsf{h}_t(w)+o(1)
\)
（固定 \(t\)，令 \(N=2^m\to\infty\)）。
在 Fold 口径下 \(w=w_m\) 且 \(p_k(w_m)=2^{-mk}S_k(m)\)（见 \S\ref{subsubsec:pom-real-q-pressure-thermo-frontier}），再由定理 \ref{thm:pom-microcanonical-fold-ht-from-power-sums} 与推论 \ref{cor:pom-microcanonical-fold-ht-low-order} 得到“有限多项式可计算性”。证毕。
\end{proof}

\begin{corollary}[二、三阶例式]\label{cor:pom-microcanonical-fold-bayes-success-examples-t2-t3}
在定理 \ref{thm:pom-microcanonical-fold-bayes-success-polynomial-in-moments} 的设定下，有
\[
\mathrm{Succ}^{\mathrm{Bayes}}_{2^m-2}(d_m)
=\frac{1+p_2(w_m)}{2}+o(1)
=\frac{1+2^{-2m}S_2(m)}{2}+o(1),
\]
\[
\mathrm{Succ}^{\mathrm{Bayes}}_{2^m-3}(d_m)
=\frac{1+3p_2(w_m)+2p_3(w_m)}{6}+o(1)
=\frac{1+3\cdot 2^{-2m}S_2(m)+2\cdot 2^{-3m}S_3(m)}{6}+o(1).
\]
\end{corollary}

\begin{theorem}[残余置换熵：\texorpdfstring{$(N-t)$}{N-t} 查询的不可消去不确定性由 \texorpdfstring{$\mathsf{h}_t(w)$}{h\_t(w)} 定量]\label{thm:pom-microcanonical-fold-residual-min-entropy}
\leanverified{paper\_pom\_microcanonical\_fold\_residual\_min\_entropy}
在均匀微正则先验下，查询 \(N-t\) 点后尚未被观测的自由度仅为未查询集 \(U\) 上的标签排列；
其后验最大概率即 Bayes 最优精确识别率。
令
\[
\mathcal{R}_t(w):=-\log \mathsf{h}_t(w).
\]
则在 \(t\) 固定且 \(N\to\infty\) 时，\(N-t\) 查询后“残余可消去不确定性”的最优 min-entropy 极限等于 \(\mathcal{R}_t(w)\)。
并且 \(\mathcal{R}_t(w)=0\) 当且仅当 \(w\) 退化为单点分布；
而 \(\mathcal{R}_t(w)\) 在均匀分布处取最大值
\[
\boxed{\;
\mathcal{R}_t(U)= -\log\!\left(\frac{\binom{n+t-1}{t}}{n^t}\right)\; } .
\]
\end{theorem}
\begin{proof}
由定理 \ref{thm:pom-microcanonical-fold-bayes-success-thermo-limit}，
\(
\mathrm{Succ}^{\mathrm{Bayes}}_{N-t}(d)\to \mathsf{h}_t(w)
\)，取负对数得到 \(\mathcal{R}_t(w)\)。
若 \(w\) 为单点分布，则 \(\mathsf{h}_t(w)=1\) 从而 \(\mathcal{R}_t(w)=0\)；
反之若 \(\mathsf{h}_t(w)=1\) 则 \(p_2(w)=1\)（由推论 \ref{cor:pom-microcanonical-fold-ht-low-order} 的 \(\mathsf{h}_2\)）从而 \(w\) 必为单点分布。
均匀分布处的最大性由定理 \ref{thm:pom-microcanonical-fold-ht-uniform-min}（均匀最小 \(\mathsf{h}_t\)）立即推出。证毕。
\end{proof}

\begin{theorem}[主纤维极点定理：\texorpdfstring{$t\to\infty$}{t->infty} 时 \texorpdfstring{$\mathsf{h}_t(w)$}{h\_t(w)} 的指数率与约化常数]\label{thm:pom-microcanonical-fold-dominant-pole-asymptotics}
令
\[
w_\star:=\max_{x\in X_m} w(x),
\qquad
r:=\card{\{x\in X_m:\ w(x)=w_\star\}}\ge 1,
\qquad
C(w):=\prod_{x:\ w(x)<w_\star}\left(1-\frac{w(x)}{w_\star}\right)^{-1}.
\]
则生成函数 \(\mathcal{H}_w(z)=\prod_x(1-w(x)z)^{-1}\) 在 \(z=1/w_\star\) 处具有 \(r\) 阶极点，并且当 \(t\to\infty\) 时
\[
\boxed{\;
\mathsf{h}_t(w)=\frac{C(w)}{(r-1)!}\,t^{\,r-1}\,w_\star^{\,t}\,(1+o(1))\; } .
\]
因此
\[
\boxed{\;
\lim_{t\to\infty}\frac{1}{t}\log \mathsf{h}_t(w)=\log w_\star\; } ,
\]
而多项式前因子完全由约化乘积常数 \(C(w)\) 与极点阶 \(r\) 决定。
\end{theorem}
\begin{proof}
将达到最大质量的 \(r\) 个因子抽出，写
\[
\mathcal{H}_w(z)=(1-w_\star z)^{-r}\cdot G_w(z),
\qquad
G_w(z):=\prod_{x:\ w(x)<w_\star}(1-w(x)z)^{-1}.
\]
则 \(G_w\) 在 \(z=1/w_\star\) 处解析且 \(G_w(1/w_\star)=C(w)\)。
并且
\[
(1-w_\star z)^{-r}=\sum_{t\ge 0}\binom{t+r-1}{r-1}w_\star^{\,t}z^t.
\]
由系数提取与解析主导项估计（最近奇点主导）得
\[
\mathsf{h}_t(w)=\binom{t+r-1}{r-1}w_\star^{\,t}C(w)+o\!\left(t^{r-1}w_\star^{\,t}\right),
\]
再用 \(\binom{t+r-1}{r-1}\sim t^{r-1}/(r-1)!\) 即得所示渐近。
对极限 \(\frac{1}{t}\log \mathsf{h}_t(w)\) 取对数并除以 \(t\) 即得 \(\log w_\star\)。证毕。
\end{proof}

\paragraph{近完全查询的统一渐近：\texorpdfstring{$t\to\infty$}{t->infty} 时的 Bayes 成功率与标度相图}
定理 \ref{thm:pom-microcanonical-fold-bayes-success-thermo-limit} 给出固定 \(t\) 的热力学极限：
\((N-t)\) 查询的 Bayes 成功率收敛到完全齐次对称多项式 \(\mathsf{h}_t(w)\)。
下述结果表明：当 \(t\) 随 \(N\) 增长但仍远小于 \(N\) 的特定统一尺度内，
该逼近依然有效，并与定理 \ref{thm:pom-microcanonical-fold-dominant-pole-asymptotics} 的主极点律拼接为一条三段标度相图。

\begin{theorem}[近完全查询的统一热力学逼近]\label{thm:pom-microcanonical-fold-bayes-success-near-complete-uniform}
\leanverified{paper\_pom\_microcanonical\_fold\_bayes\_success\_near\_complete\_uniform}
设字母表大小 \(n=\card{X_m}\) 固定，且存在分布 \(w\in\Delta(X_m)\)（各坐标严格正）使 \(d(x)/N\to w(x)\)。
取 \(t=t(N)\to\infty\) 且满足
\[
t=o(N^{2/3}),
\qquad
k=N-t.
\]
则在均匀微正则先验 \(F\sim\mathrm{Unif}(\mathfrak{Fold}_m(d))\) 下，\((N-t)\) 查询的 Bayes 最优精确恢复成功率满足
\[
\boxed{\;
\mathrm{Succ}^{\mathrm{Bayes}}_{N-t}(d)
=\mathsf{h}_t(w)\,(1+o(1))\;}
\qquad(N\to\infty),
\]
其中 \(\mathsf{h}_t(w)\) 定义见定理 \ref{thm:pom-microcanonical-fold-bayes-success-thermo-limit}。
\end{theorem}
\begin{proof}
由定理 \ref{thm:pom-microcanonical-fold-bayes-success-Nminus-t}，
在 \(k=N-t\) 查询后未查询集上的剩余计数向量 \(R=(R_x)_{x\in X_m}\) 服从多元超几何分布，
且
\[
\mathrm{Succ}^{\mathrm{Bayes}}_{N-t}(d)=\EE\!\left[\frac{\prod_x R_x!}{t!}\right],
\qquad
0\le \frac{\prod_x R_x!}{t!}\le 1.
\]
令 \(\widetilde R\sim\mathrm{Mult}(t;w)\)。
由定理 \ref{thm:pom-microcanonical-fold-bayes-success-thermo-limit} 的定义可知
\(
\EE[\prod_x \widetilde R_x!/t!]=\mathsf{h}_t(w)
\)。
因此
\[
\left|\mathrm{Succ}^{\mathrm{Bayes}}_{N-t}(d)-\mathsf{h}_t(w)\right|
\le D_{\mathrm{TV}}\!\left(\mathrm{Law}(R),\ \mathrm{Law}(\widetilde R)\right).
\]
当 \(n\) 固定且 \(w\) 严格正时，
对超几何 pmf 与多项 pmf 的对数似然比作 Stirling 展开（对每个坐标的阶乘比展开到二阶，并控制余项），
在 \(t=o(N^{2/3})\) 下余项为 \(o(1)\) 且在典型区域上一致；
再结合多项分布的高斯集中（\(\widetilde R=tw+O_{\mathbb{P}}(\sqrt t)\)）可得上述对数似然比在概率意义下趋于 \(0\)，并由此推出总变差距离为 \(o(1)\)。
代回即得结论。证毕。
\end{proof}

\begin{corollary}[近完全查询的主指数与约化常数]\label{cor:pom-microcanonical-fold-bayes-success-near-complete-pole-law}
在定理 \ref{thm:pom-microcanonical-fold-bayes-success-near-complete-uniform} 的条件下，若进一步 \(t\to\infty\)，
则由定理 \ref{thm:pom-microcanonical-fold-dominant-pole-asymptotics}（记号 \(w_\star,r,C(w)\) 同该定理）得到
\[
\boxed{\;
\mathrm{Succ}^{\mathrm{Bayes}}_{N-t}(d)
=\frac{C(w)}{(r-1)!}\,t^{\,r-1}\,w_\star^{\,t}\,(1+o(1))\;}
\]
从而
\[
\boxed{\;
\log \mathrm{Succ}^{\mathrm{Bayes}}_{N-t}(d)
=t\log w_\star+(r-1)\log t+\log\!\left(\frac{C(w)}{(r-1)!}\right)+o(1)\;}
\qquad(N\to\infty).
\]
\end{corollary}

\begin{corollary}[全恢复成功率的三段标度相图]\label{cor:pom-microcanonical-fold-bayes-success-phase-diagram}
在推论 \ref{cor:pom-microcanonical-fold-bayes-success-near-complete-pole-law} 的条件下（因此 \(n\) 与 \(w\) 固定且 \(w_\star<1\)），
在 \(t=o(N^{2/3})\) 的近完全查询区间内有如下三段标度：
\begin{enumerate}
  \item \textbf{常数量级（\(t=O(1)\)）}：\(\mathrm{Succ}^{\mathrm{Bayes}}_{N-t}(d)\to \mathsf{h}_t(w)\in(0,1]\)。
  \item \textbf{幂律量级（\(t=\lfloor \lambda\log N\rfloor\)）}：对任意固定 \(\lambda>0\)，
  \[
  \mathrm{Succ}^{\mathrm{Bayes}}_{N-t}(d)
  =N^{\lambda\log w_\star+o(1)}(\log N)^{r-1+o(1)}.
  \]
  特别地，斜率 \(\lambda_\ast:=1/|\log w_\star|\) 对应主阶 \(\mathrm{Succ}^{\mathrm{Bayes}}_{N-t}(d)=N^{-1+o(1)}(\log N)^{r-1+o(1)}\)。
  \item \textbf{指数衰减（\(t/\log N\to\infty\)）}：若 \(t/\log N\to\infty\) 且仍有 \(t=o(N^{2/3})\)，则
  \[
  \frac{1}{t}\log \mathrm{Succ}^{\mathrm{Bayes}}_{N-t}(d)\ \to\ \log w_\star<0.
  \]
\end{enumerate}
\end{corollary}

\endinput


综上，固定纤维谱的微正则先验把“宏观纤维分布 \(w\)”与“微观可逆性”之间的关系压缩为一条完全可计算的对称函数链：
在 \(N-t\) 查询的热力学极限下，Bayes 最优精确识别率由 \(\mathsf{h}_t(w)\) 给出；
其低阶仅依赖有限碰撞矩 \(p_2(w),\dots,p_t(w)\)（在 Fold 口径下等价于 \(S_2(m),\dots,S_t(m)\)），
而其大阶渐近由主纤维质量 \(w_\star\) 与约化常数 \(C(w)\) 统摄。
在近完全查询的统一尺度 \(t=o(N^{2/3})\) 内，该描述仍保持有效，并由此诱导出由 \(w_\star\) 控制的常数/幂律/指数三段标度。

\endinput
